\documentclass[11pt,a4paper]{article}
\usepackage[utf8]{inputenc}
\usepackage[T1]{fontenc}
\usepackage[english]{babel}
\usepackage[margin=2.2cm]{geometry}
\usepackage{amsmath,amssymb,amsthm}
\usepackage{parskip}
\usepackage{enumitem}
\setlist{nosep,leftmargin=1.4em}
\usepackage{booktabs}
\usepackage{array}
\usepackage{xcolor}
\definecolor{linkblue}{RGB}{20,60,150}
\definecolor{defbg}{RGB}{240,244,252}
\definecolor{defframe}{RGB}{100,130,190}
\definecolor{markcol}{RGB}{176,84,0}
\definecolor{markbg}{RGB}{253,246,236}
\usepackage[most]{tcolorbox}
\newtcolorbox{defbox}[1][]{breakable,colback=defbg,colframe=defframe,boxrule=0.6pt,arc=1.5mm,left=2mm,right=2mm,top=1mm,bottom=1mm,title={#1},fonttitle=\bfseries}
\newtcolorbox{poznbox}[1][]{breakable,colback=markbg,colframe=markcol,boxrule=0.6pt,arc=1.5mm,left=2mm,right=2mm,top=1mm,bottom=1mm,title={#1},fonttitle=\bfseries}
\newtheorem{theorem}{Theorem}[section]
\theoremstyle{definition}
\newtheorem{definition}{Definition}[section]
\usepackage[colorlinks=true,linkcolor=linkblue,urlcolor=linkblue]{hyperref}
\usepackage{algorithm}
\usepackage{algpseudocode}

\floatname{algorithm}{Algorithm}
\renewcommand{\algorithmicrequire}{\textbf{Input:}}
\renewcommand{\algorithmicensure}{\textbf{Output:}}

% Quotation macro carried over from the Czech original (babel-czech \uv).
\providecommand{\uv}[1]{``#1''}

% Marker for material that is not in the NDBI023 / NAIL116 slides.
\newcommand{\mimo}{\textcolor{markcol}{\textsf{\footnotesize[$\blacklozenge$~beyond the slides]}}}
\newcommand{\mimoq}[1]{\textcolor{markcol}{\textsf{\footnotesize[$\blacklozenge$~beyond the slides: #1]}}}

\title{\textbf{Data Mining}\\[2mm]
\large Study Text for the Final State Examination}
\author{Final State Examination in Artificial Intelligence, MFF UK}
\date{August 2026}

\begin{document}
\maketitle

\begin{center}
\begin{minipage}{0.92\textwidth}
\small
This text covers the examination topic \emph{Data Mining} of the master's final state
examination in Artificial Intelligence at MFF UK. Its chapters correspond \textbf{one to one}
to the sentences of the official wording of the topic (see the mapping table below). It is
based on the slides of the lecture courses \textbf{NDBI023 Data Mining} and
\textbf{NAIL116 Social Networks and their Analysis} (I.~Mrázová, MFF UK), supplemented with the
standard material of the textbooks Han, Kamber, Pei: \emph{Data Mining --- Concepts and
Techniques}, Tan, Steinbach, Kumar: \emph{Introduction to Data Mining}, Aggarwal: \emph{Data
Mining --- The Textbook}, and for networks Easley, Kleinberg: \emph{Networks, Crowds and
Markets}, Newman: \emph{Networks} and Zafarani, Abbasi, Liu: \emph{Social Media Mining}. It
contains definitions, pseudocode, derivations, numerical examples and comparison tables; every
chapter ends with a summary for the oral examination. This is an English translation of the
Czech study text \emph{Dobývání znalostí}.
\end{minipage}
\end{center}

\begin{center}
\begin{minipage}{0.92\textwidth}
\begin{poznbox}[Scope and marking]
\small
The length of the chapters reflects the weight of the material: most space goes to the
data mining methods (Chapter~\ref{sec:metody}) and to social network analysis
(Chapter~\ref{sec:sna}) --- that is, to what both lecture courses treat in most detail and what
the topic names most concretely.

The marker \mimo{} flags material that is \textbf{in the slides of neither course} (the slides
are available in full) but has been kept in the text --- either because \emph{the official
wording of the topic names it}, or because the surrounding exposition would not make sense
without it. It is an offer for deletion: whatever is not marked is backed by the slides. Where
the marker is followed by a colon, it states precisely what the reservation applies to ---
typically that a summarising table or framework is mine, while the individual items inside it do
come from the slides.

The slides cover: data analysis and clustering, association rules, decision trees and
ensembles, SVM with classifier evaluation, Bayesian models and ELM (NDBI023); basic concepts,
influence modelling, homophily and affiliation, link analysis, and complex networks with
community detection and link prediction (NAIL116). The most conspicuous gap is the
\textbf{KDD process and methodologies} (Chapter~\ref{sec:paradigmata}). The chapter on
\textbf{characteristic and discriminant rules} (Chapter~\ref{sec:pravidla}) is built
primarily on the slides of the course \emph{Internet and classification methods}
(Holeňa, lectures \texttt{ikm1}--\texttt{ikm6}, above all \texttt{ikm5}); only the
delimitation of the two kinds of rules with the t-weight/d-weight measures follows
Han--Kamber there. The slides \texttt{nmr1}--\texttt{nmr5} from the same folder belong
to a~different course and do not cover the subject matter of the chapter (see the note
at its start).
\end{poznbox}
\end{minipage}
\end{center}

\vspace{2mm}

\section*{Mapping of the examination topic to chapters}
\addcontentsline{toc}{section}{Mapping of the examination topic to chapters}

\begin{center}
\small
\begin{tabular}{@{}>{\raggedright\arraybackslash}p{8.6cm}>{\raggedright\arraybackslash}p{7.4cm}@{}}
\toprule
\textbf{Sentence of the official topic} & \textbf{Chapter in this text} \\
\midrule
Basic paradigms of data mining.
  & \ref{sec:paradigmata}~Basic paradigms of data mining \\
Data preparation; attribute selection and methods for the analysis of their relevance.
  & \ref{sec:priprava}~Data preparation, attribute selection and relevance analysis \\
Methods for data mining; association rules, approaches based on the principle of
supervised learning, and cluster analysis.
  & \ref{sec:metody}~Methods for data mining (\ref{sec:metody}.1~association rules,
    \ref{sec:metody}.2~supervised learning, \ref{sec:metody}.3~cluster analysis) \\
Methods for the extraction of characteristic discriminant rules and the measurement of their
interestingness.
  & \ref{sec:pravidla}~Characteristic and discriminant rules and their interestingness \\
Representation, evaluation and visualisation of the acquired knowledge.
  & \ref{sec:reprez}~Representation, evaluation and visualisation of knowledge \\
Models for social network analysis; centrality measures, community detection.
  & \ref{sec:sna}~Models for social network analysis, centrality and community detection \\
Practical use of techniques for data mining and social network analysis.
  & \ref{sec:praxe}~Practical use of the techniques \\
\bottomrule
\end{tabular}
\end{center}

\vspace{1mm}
\noindent\small\textit{A note on terminology.} The Czech original gives the English equivalent
in italics in parentheses at the first occurrence; in this English version the italicised terms
are kept where they mark a technical term being introduced.\normalsize

\vspace{2mm}
\tableofcontents
\newpage

%=====================================================================
\section{Basic paradigms of data mining}
\label{sec:paradigmata}
%=====================================================================

This chapter delimits the field, its process, and the basic classification of tasks and data.
It is the frame into which all the other chapters fit: without it one cannot explain why data
preparation is a chapter of its own, and why the methods are divided precisely into association
rules, supervised learning and clustering.

\subsection{Motivation and delimitation of the field}

Organisations accumulate data in volumes that far exceed a human's capacity to inspect them.
Classical querying (SQL, OLAP) can answer a question we already know how to ask;
\emph{data mining} answers the questions we did not know how to ask --- it looks in the
data for structures whose existence we did not know about in advance.

An illustrative example: a system went through 20 million procedures billed to a health
insurance company, flagged 214 thousand of them as odd, and passed the 14 thousand most
suspicious ones on for manual review. The reviewing physicians, who are able to inspect on the
order of hundreds of cases, immediately marked 8 thousand of the submitted cases as fraudulent.
The machine did not replace the expert here --- it \emph{focused} their attention.

\begin{defbox}[Knowledge discovery in databases]
\emph{Knowledge discovery in databases} (KDD) is the non-trivial process of extracting implicit,
previously unknown and potentially useful information from data.

All four words matter:
\begin{itemize}
\item \textbf{non-trivial} --- it is not a mere query or aggregation;
\item \textbf{implicit} --- the knowledge is not written down explicitly in the data;
\item \textbf{previously unknown} --- a result everybody already knows is worthless;
\item \textbf{potentially useful} --- it must lead to an action or a decision.
\end{itemize}
\end{defbox}

The term \emph{data mining} (DM) refers, in the narrow sense, only to one step of the KDD
process (the actual application of the analytical algorithms); in common practice, however, the
two terms are used interchangeably. The field took shape in the 1990s at the intersection of
three disciplines:

\begin{itemize}
\item \textbf{artificial intelligence} --- in particular machine learning methods (rule
      induction, trees, neural networks);
\item \textbf{database systems} --- storage of and efficient access to large data sets, data
      warehouses, information retrieval;
\item \textbf{statistics} --- modelling and analysis of the dependencies found in the data,
      hypothesis testing.
\end{itemize}

The fourth --- and in practice decisive --- component is the need to use the processed data to
support (strategic) decision-making in an enterprise. Without it, data mining reduces
to an academic exercise.

\subsection{The KDD process}

\mimoq{the five-phase framework as such; the activities inside the phases are in the slides}
KDD is an \textbf{interactive and iterative} process: the results of
one step routinely lead back to the preceding steps. The classical scheme (Fayyad et al.) has
five phases:

\begin{center}
\small
\begin{tabular}{@{}p{3.3cm}p{4.2cm}p{7.6cm}@{}}
\toprule
\textbf{Phase} & \textbf{Input $\to$ output} & \textbf{Content} \\
\midrule
Selection & original data $\to$ selected data & selection of the relevant records and
  attributes from the data warehouse \\
Preprocessing & selected data $\to$ preprocessed data & cleaning, missing values, noise,
  outliers, integration of sources \\
Transformation & preprocessed data $\to$ transformed data & normalisation, discretisation,
  derived attributes, dimensionality reduction \\
Data mining proper & transformed data $\to$ patterns & application of the algorithms (trees,
  associations, clusters, \dots) \\
Interpretation & patterns $\to$ knowledge & evaluation from the end user's point of view,
  visualisation, deployment \\
\bottomrule
\end{tabular}
\end{center}

The goal of the first three phases (collectively \emph{data preparation},
Chapter~\ref{sec:priprava}) is to build, out of data stored in a complex structure, \textbf{one
table} containing the relevant information about the objects under study (a bank's clients,
customers, \dots). This step typically consumes 60--80~\% of the total project time.

\textbf{Seven steps through a manager's eyes.} \mimo{} An organisational view of the same
process (the impulse is a problem, the output is the knowledge to solve it):
\begin{enumerate}
\item \textbf{team} --- experts on the domain, on the data, and on KDD methods;
\item \textbf{problem specification} --- \uv{increase profit} $\to$ \uv{predict customer churn};
\item \textbf{data collection} --- including external data on the environment (season,
      campaigns, weather);
\item \textbf{choice of methods};
\item \textbf{preprocessing};
\item \textbf{mining} --- repeatedly; the results tune the parameters of later runs;
\item \textbf{interpretation} --- an analytical report, possibly an action.
\end{enumerate}

\subsection{Types of data mining tasks}

Tasks are classified from two independent points of view: \emph{what we want from the result}
and \emph{whether a target variable is available}.

\mimo{} \textbf{By what we want from the result:}
\begin{itemize}
\item \textbf{Classification / prediction} --- classify new cases, resp.\ estimate future
      development (weather, share prices, electricity consumption). What matters is
      \emph{accuracy}, not simplicity; the result tends to be a lot of knowledge, hard to
      interpret.
\item \textbf{Description} (characterisation, profile) --- capture the dominant structure or
      relations in the data (the profile of a creditworthy client). What matters is
      \emph{comprehensibility} with full coverage of the concept; the knowledge is smaller
      and less accurate.
\item \textbf{Search for \uv{nuggets}} --- find new, \emph{surprising} knowledge (an unusual
      combination of goods in a basket, a suspicious pattern of transactions); it \emph{need
      not} cover the whole concept.
\end{itemize}

\begin{defbox}[Descriptive vs.\ predictive tasks]
\textbf{Descriptive} tasks characterise the general properties of the data: characterisation
and discrimination, association and correlation analysis, clustering, outlier detection,
evolution analysis.

\textbf{Predictive} tasks perform inference over the data in order to predict: classification
(discrete target attribute), regression (continuous target attribute), time series prediction.
\end{defbox}

\textbf{By the availability of a target variable:} \emph{supervised learning} ---
classification, regression; \emph{unsupervised learning} --- clustering, association rules,
outlier detection; \emph{semi-supervised} --- a small labelled and a large unlabelled part of
the data.

This division is also the key to Chapter~\ref{sec:metody}: the topic names in it exactly one
representative of unsupervised learning on transactional data (association rules), supervised
learning, and unsupervised learning on metric data (cluster analysis).

\subsection{Types of data}

\textbf{By attribute type} --- this determines which operations are meaningful and which methods
can be used:

\begin{center}
\small
\begin{tabular}{@{}p{3.3cm}p{2.2cm}p{5.2cm}p{4.0cm}@{}}
\toprule
\textbf{Scale} & \textbf{Operations} & \textbf{Example} & \textbf{Statistics} \\
\midrule
Nominal & $=$, $\neq$ & colour, kind of fruit & mode, frequencies \\
Ordinal & $=$, $<$ & grade, \{low, middle, high\} & median, quantiles \\
Interval-scaled & $+$, $-$ & temperature in $^{\circ}$C, date & mean, standard deviation \\
Ratio-scaled & $\times$, $/$ & income, weight, frequency & geometric mean \\
\bottomrule
\end{tabular}
\end{center}

For the interval scale the key point is that \emph{intervals keep the same importance
throughout the scale}: the difference in age between 10 and 20 is the same as that between 40
and 50. The ratio scale additionally has an absolute zero, so ratios are meaningful;
ratio-scaled attributes with exponential behaviour (e.g.\ the total amount of microorganisms
$Ae^{Bt}$) are first log-transformed and then treated like interval-scaled ones. A
\emph{nominal} attribute with $v$ values is converted, for numeric methods, into $v$ binary
attributes (value 1 for the corresponding category, 0 otherwise); an \emph{ordinal} attribute is
commonly treated as interval-scaled.

\textbf{By the structure of the data set:}
\begin{itemize}
\item \textbf{Tabular / relational data} --- objects $\times$ attributes; the standard input of
      most methods.
\item \textbf{Transactional data} --- records of variable length (the shopping basket); the
      input of association rules.
\item \textbf{Temporal data} --- e.g.\ time series of stock prices. A typical task is to predict
      future values; the series is transformed into a table by a sliding window: inputs
      $y(t_0),y(t_1),y(t_2),y(t_3)$ $\to$ output $y(t_4)$; $y(t_1),\dots,y(t_4)\to y(t_5)$
      and so on.
\item \textbf{Spatial data} --- geographic information systems; an implicit neighbourhood
      relation (the attribute values of adjacent objects do not differ too much), useful
      especially for interpretation.
\item \textbf{Structural data} --- e.g.\ chemical compounds written as SMILES strings or as
      files of facts in Prolog; a task might be the classification of substances by their
      structure.
\item \textbf{Text and web data, multimedia, graphs and networks} --- these require a specific
      representation (the vector model, graph models, see Chapter~\ref{sec:sna}).
\end{itemize}

\subsection{Data mining methodologies}
\label{ssec:metodologie}

\mimo{} The purpose of a methodology is to provide a uniform framework and to guide
applications regardless of the line of business --- to allow the sharing and transfer of
experience from successful projects. Methodologies arise either at the large BI tool vendors
(SEMMA by SAS Institute, 5A by SPSS), or as neutral, software-independent standards (CRISP-DM).

\textbf{SEMMA} (SAS Enterprise Miner), five phases: \emph{Sample} --- selection of data for
modelling (sampling, imputation, partitioning into training, test and validation sets);
\emph{Explore} --- visual exploration and reduction of the data (outliers, multidimensional
histograms, clustering, Kohonen self-organising maps); \emph{Modify} --- preparation of objects,
values and variables (selection, creation, transformation); \emph{Model} --- application of the
methods (trees, regression, neural networks); \emph{Assess} --- evaluation of reliability and
usefulness; \emph{the user must understand them}.

\textbf{CRISP-DM} (CRoss-Industry Standard Process for Data Mining) --- an open process model
neutral with respect to industry, tool and application; the development was funded by the ESPRIT
initiative, and the consortium consisted of ISL (later SPSS, today IBM), Teradata, NCR,
Daimler AG and the insurance company OHRA. Six phases whose \textbf{order is not strict} ---
the results of one phase determine the choice of the next steps and some steps must be repeated
(the process is drawn as a cycle with an outer cycle of deployment and monitoring):

\begin{enumerate}
\item \textbf{Business Understanding} --- business objectives and success criteria; an inventory
      of resources (people, information, software, computing capacity), assumptions and
      constraints, risks, costs and benefits; the data mining goals; the project plan.
\item \textbf{Data Understanding} --- collection of the initial data; its description (quantity,
      properties, format); exploration and visualisation (distributions of the key attributes,
      relations of pairs of attributes, important subpopulations, simple statistics);
      verification of data quality.
\item \textbf{Data Preparation} --- selection, cleaning, construction of derived attributes,
      integration, aggregation, formatting; performed iteratively and in varying order.
\item \textbf{Modeling} --- choice of technique (with regard to its assumptions, which may force
      further adjustments of the data), test design, building the model with various parameters,
      assessing the model.
\item \textbf{Evaluation} --- the results against the \emph{business} success criteria; approval
      of the models; review of the process; determination of the next steps.
\item \textbf{Deployment} --- presentation in a form easily interpretable and deployable by the
      user; deployment plan, monitoring and maintenance; final report and project review
      (documentation of the difficulties and blind alleys).
\end{enumerate}

\textbf{ASUM-DM} (IBM, 2015) responds to the main criticism of CRISP-DM --- its weak support for
project management. Its phases are \emph{Analyze}, \emph{Design}, \emph{Configure \& Build},
\emph{Deploy}, \emph{Operate \& Optimize}, with \emph{Project Management} running across all of
them.

\begin{center}
\small
\begin{tabular}{@{}llll@{}}
\toprule
& \textbf{SEMMA} & \textbf{CRISP-DM} & \textbf{ASUM-DM} \\
\midrule
Origin & SAS Institute & ESPRIT / consortium & IBM (2015) \\
Neutrality & tied to SAS EM & tool-neutral & tied to IBM \\
Business phase & missing & yes (phases 1 and 5) & yes \\
Project management & no & weak & explicit, cross-cutting \\
Number of phases & 5 & 6 & 5 + PM \\
\bottomrule
\end{tabular}
\end{center}

\subsection{Summary for the examination}

\begin{itemize}
\item \textbf{KDD} $=$ the non-trivial process of extracting implicit, previously unknown and
      potentially useful information from data; interactive and iterative.
\item \textbf{Five KDD phases}: selection $\to$ preprocessing $\to$ transformation $\to$ mining
      $\to$ interpretation. Data preparation consumes 60--80~\% of the time.
\item \textbf{Roots of the field}: artificial intelligence (machine learning), databases,
      statistics $+$ the need for decision support. \textbf{OLAP} answers questions posed in
      advance, KDD looks for patterns we do not know about.
\item \textbf{Three types of task}: classification/prediction (accuracy $>$ simplicity),
      description (comprehensibility, full coverage), search for nuggets (surprise, need not
      cover).
\item \textbf{Descriptive vs.\ predictive} tasks; \textbf{supervised vs.\ unsupervised} vs.\
      \textbf{semi-supervised}.
\item \textbf{Attribute scales}: nominal, ordinal, interval, ratio --- they determine the
      admissible operations and the methods. Nominal $\to$ $v$ binary attributes; a ratio-scaled
      attribute with exponential behaviour $\to$ take logarithms.
\item \textbf{Methodologies}: SEMMA (Sample, Explore, Modify, Model, Assess), CRISP-DM (6
      phases, order not strict, cyclic), ASUM-DM (CRISP-DM $+$ project management).
\item The most frequent follow-up question: \emph{why is the order of the CRISP-DM phases not
      strict?} --- because the results of a phase determine the next steps; e.g.\ the data
      quality found leads back to a reformulation of the goal.
\end{itemize}

%=====================================================================
\newpage
\section{Data preparation, attribute selection and relevance analysis}
\label{sec:priprava}
%=====================================================================

The sentence of the topic joins two things into one chapter: \emph{data preparation} (what has
to be done so that the data is processable at all) and \emph{attribute selection with relevance
analysis} (how to pick, out of many attributes, those that carry information). The second part
is really a special case of the first --- dimensionality reduction --- but the topic names it
separately, because knowledge of concrete criteria and the ability to compute them are expected.

\subsection{Why data preparation is crucial}

Preprocessing is \emph{of key importance for the success of the application}, it is difficult and
time-consuming, and collaboration with a domain expert is an advantage. The goals are:
\begin{itemize}
\item to select (or create), from the available data, the data relevant for the task at hand,
\item to represent it in a form suitable for the chosen algorithm,
\item the final state being \emph{(one)} data table containing the attribute values of the
      objects.
\end{itemize}

\mimo{} Data quality is judged along several dimensions: \emph{accuracy}, \emph{completeness},
\emph{consistency}, \emph{timeliness}, \emph{believability}, \emph{interpretability}. The
principle \uv{garbage in, garbage out} holds without exception: the best algorithm on bad data
produces only confident nonsense.

\subsection{Data cleaning: missing values}

The basic strategies (in order from the simplest to the most accurate):

\begin{enumerate}
\item \textbf{Ignore objects} with any missing value. Simple, but with many attributes we lose
      most of the data; moreover it introduces bias if the missing values are not random.
\item \textbf{Replace with a new value} \uv{I don't know}. The missing value becomes a
      full-fledged category --- appropriate when the absence itself carries information (an
      unfilled income field in a loan application).
\item \textbf{Replace with an existing attribute value}: the most frequent value (the mode; for
      numeric attributes the mean or median), proportionally to the ratio of all values
      (randomly according to the empirical distribution), or any value.
\item \textbf{Impute the value using a model} --- train a predictive model (regression, decision
      tree, $k$-NN) that predicts the missing attribute from the others. The most accurate, the
      most computationally demanding, and it risks introducing an artificial dependence between
      attributes.
\end{enumerate}

Some methods can handle missing values themselves: C4.5 ignores them when constructing the tree
(the splitting criterion is computed only from the known values) and estimates them at
classification time from other patterns; CART does not count missing values towards the decision
criterion; naive Bayes can classify even incompletely described patterns
(Chapter~\ref{sec:metody}).

\subsection{Data cleaning: noise and outliers}

\emph{Noise} is a random error or variance in a measured quantity. Smoothing techniques:
\textbf{binning} (values are divided into bins and replaced by the bin mean, median or boundary),
\textbf{regression} (values are replaced by the prediction of a regression function),
\textbf{clustering} (values outside clusters are marked as outliers) and a \textbf{combination
of machine and human} (suspicious values are submitted for review).

\begin{defbox}[Range, quartiles, interquartile range, outlier]
The \emph{range} of a set of data is the difference between the highest and the lowest value.

The \emph{quartiles} $Q_1,Q_2,Q_3$ split the data into four parts: 25~\% of the observations
have a value less than the lower quartile $Q_1$, 50~\% below the median $Q_2$, and 75~\% below
the upper quartile $Q_3$.

The \emph{interquartile range}: $\mathrm{IQR}=Q_3-Q_1$.

An \emph{outlier} is an extreme value that lies outside the overall pattern of the data. The
commonly used rule: an outlier is any value with
\[ x > Q_3 + 1.5\cdot \mathrm{IQR} \qquad\text{or}\qquad x < Q_1 - 1.5\cdot \mathrm{IQR}. \]
\end{defbox}

\textbf{Computing the quartiles} for $n$ values sorted in ascending order (discrete data): for
$Q_1$ divide $n$ by 4 --- when $n/4$ is a whole number, take the midpoint of the corresponding
term and the term above; when it is not, round up and pick the corresponding term. For $Q_3$
proceed analogously with $3n/4$. For continuous data, interpolate instead of rounding.

\begin{defbox}[Example --- box plot and outliers]
The blood glucose of 30 females (mmol/l), stem-and-leaf diagram (key: $2\mid 1$ means 2.1):
\begin{center}\ttfamily
2 | 2 2 3 3 5 7 \hfill (6)\\
3 | 1 2 6 7 7 7 8 8 8 8 9 9 9 \hfill (13)\\
4 | 0 0 0 0 4 5 6 7 8 \hfill (9)\\
5 | 1 5 \hfill (2)
\end{center}
$n=30$, $30/4=7.5$ $\to$ round up to 8, the 8th number is 3.2, hence $Q_1=3.2$. The median is
$Q_2=3.8$. For $Q_3$: $3\cdot 30/4=22.5\to 23$, the 23rd number is 4.0, hence $Q_3=4.0$.

$\mathrm{IQR}=4.0-3.2=0.8$; lower limit $=3.2-1.5\cdot 0.8=2.0$; upper limit
$=4.0+1.5\cdot 0.8=5.2$.

\textbf{The outlier is the value 5.5.} The highest value which is not an outlier is 5.1 --- in
the box plot a whisker leads to it, and the outlier is marked with a cross.
\end{defbox}

The multiplier 1.5 is a convention; with a multiplier of 1 we get a stricter criterion. For male
turtles with $Q_1=400$~g, $Q_3=580$~g we have $\mathrm{IQR}=180$ and the limits
$400-180=220$ and $580+180=760$ --- none of the values 400, 260, 550, 640~g is an outlier, and
the largest weight without being an outlier is 760~g. For females with $Q_1=260$, $Q_3=340$,
$\mathrm{IQR}=80$ the limits are 180 and 420, so the outliers are the values 170~g and 440~g.

\subsection{Data integration}
\label{ssec:integrace}

\mimo{} In \emph{data integration} we merge several sources into one table. The main problems:
the \textbf{entity identification problem} (the attribute \texttt{cust\_id} in one database
corresponds to \texttt{customer\_number} in another; metadata and record linkage are needed);
\textbf{redundancy} (an attribute derivable from others --- annual income $=$ 12$\times$monthly;
detected by correlation analysis for numeric and by the $\chi^2$ test for categorical
attributes, see Section~\ref{ssec:filtr}); \textbf{value conflicts} (different units, scales,
encodings); \textbf{duplicates} (the same object several times with slightly different values).

\subsection{Transformation and standardisation of attributes}

In Euclidean space, standardisation of attributes is recommended so that all attributes can have
an \emph{equal impact} on the computation of distances. For the points $x_i=(0.1;\,20)$ and
$x_j=(0.9;\,720)$ we get
\[ \mathrm{dist}(x_i,x_j)=\sqrt{(0.9-0.1)^2+(720-20)^2}=700.000457, \]
so the distance is entirely dominated by the difference $720-20=700$ in the second attribute;
the first attribute has practically no influence. Standardisation forces the attributes into a
common value range.

\begin{defbox}[Basic standardisation transformations]
Let $f$ be an attribute and $x_{if}$ its value for the $i$-th object.

\textbf{Decimal scaling} to $[-1,1]$: divide the values by the smallest power of 10 that keeps
all the transformed values within $[-1,1]$.

\textbf{Min-max normalisation} (range standardisation) to $[0,1]$:
\[ \mathrm{range}(x_{if})=\frac{x_{if}-\min(f)}{\max(f)-\min(f)}, \]
and to $[-1,1]$ then $\mathrm{range}_{[-1,1]}(x_{if})=2\cdot \mathrm{range}(x_{if})-1$.

\textbf{Standardisation according to the mean absolute difference} --- zero mean and unit mean
absolute difference:
\[ m_f=\tfrac1n\big(x_{1f}+\cdots+x_{nf}\big),\qquad
s_f=\tfrac1n\big(|x_{1f}-m_f|+\cdots+|x_{nf}-m_f|\big),\qquad
z(x_{if})=\frac{x_{if}-m_f}{s_f}. \]
\end{defbox}

The mean absolute difference is \emph{robust} to outliers (the deviations are not squared), so
outliers compress the remaining data less.

\subsection{Discretisation of numeric attributes}

A number of methods (association rules, ID3, Bayesian classification with discrete
attributes) require categorical input. \emph{Discretisation} divides the range of a numeric
attribute into intervals. The methods differ in four respects: the \emph{strategy} (top-down
division of intervals, or bottom-up joining of intervals), the \emph{number} of intervals, the
\emph{type} of intervals, and the \emph{criterion} of interval quality (minimum classification
error, entropy, the $\chi^2$ test).

\begin{itemize}
\item \textbf{Equi-width} --- the range $[\min(f),\max(f)]$ is divided into a preset number
      $k$ of intervals of the same length $(\max(f)-\min(f))/k$. The simplest (it needs
      neither the classes nor sorting), but sensitive to outliers and to a skewed
      distribution: a single outlier stretches the range and most objects then end up in a
      single interval.
\item \textbf{Equi-depth} \mimo{} --- the values are sorted in ascending order and a boundary
      is placed after every $n/k$ objects, so each interval contains the same number of
      objects. Robust to outliers, but the boundaries need not respect natural clusters of
      the values.
\item \textbf{Using the class information} --- the boundaries are chosen so that the intervals
      are as pure as possible with respect to the target class. \emph{Top-down} (splitting):
      a candidate boundary is every place between adjacent sorted values where the class
      changes; the boundary with the smallest weighted sum of entropies of the resulting
      parts is selected and the splitting recurses as long as it brings an improvement (the
      same way C4.5 chooses the threshold of a numeric attribute,
      Chapter~\ref{sec:metody}). \emph{Bottom-up} (merging): initially every value has its
      own interval and the adjacent pair with the most similar class distributions (the
      smallest $\chi^2$) is merged repeatedly, until the $\chi^2$ of every adjacent pair
      exceeds a chosen significance threshold.
\end{itemize}

\subsubsection{Fuzzy discretisation}

Sharp boundaries have an unpleasant consequence: two nearly identical values just on the two
sides of a boundary fall into different categories. \emph{Fuzzy discretisation} blurs the
boundaries --- the characteristic function $\mu$ of an interval does not drop from 1 to 0 in a
step at the boundary, but linearly over a transition zone. Suitable intervals are found by the
procedure \textbf{Interval}: the attribute values are sorted in ascending order and maximal
sequences of values with the same class are concatenated into intervals with $\mu=1$; stretches
with mixed classes (and the gaps between the intervals) receive the class \uv{?}. Every
interval $Int_i$ of class \uv{?} is then resolved according to its neighbours: if
$Int_{i-1}$ and $Int_{i+1}$ belong to the same class, all three are joined into one crisp
interval with $\mu=1$; otherwise \emph{two overlapping fuzzy intervals}
$Int_{i-1}\cup Int_i$ and $Int_i\cup Int_{i+1}$ arise --- each has $\mu=1$ on its original
crisp part, and over the \uv{?} zone (between the upper limit $x_{i-1}$ of $Int_{i-1}$ and the
lower limit $x_{i+1}$ of $Int_{i+1}$) the characteristic functions decrease and increase
linearly, respectively:
\[ \mu_{i-1}(x)=\frac{x_{i+1}-x}{x_{i+1}-x_{i-1}},\qquad
   \mu_{i+1}(x)=\frac{x-x_{i-1}}{x_{i+1}-x_{i-1}},\qquad x\in[x_{i-1},x_{i+1}]. \]

One numeric value may thus correspond to \emph{two} adjacent categories whose characteristic
functions sum to 1. An object with a value in the transition zone enters the data as a
\uv{fuzzy object} with a weight given by the product of the characteristic functions of all
attributes,
\[ w = \prod_j \mu_j < 1. \]
Discretisation in general may lose information and introduce \emph{contradictions} --- objects
with the same discretised input values that belong to different classes.

\subsubsection{Grouping the values of categorical attributes (KEX)}

A categorical attribute with too many values is handled by grouping its values; the goal is to
form groups $\mathrm{Grp}$ for which $P(\mathrm{Class}\mid A(\mathrm{Grp}))$ differs
significantly from the a-priori $P(\mathrm{Class})$. The algorithm: \emph{(1)}~build an ordered
list of the attribute values; \emph{(2)}~for each value, compute the class frequencies among
the objects with this value and let the procedure \textbf{Assign} give it a code --- if all
such objects belong to the same class, the code of this class; if the class distribution for
this value differs \emph{significantly} (by the $\chi^2$ test) from the a-priori one, the code
of the most frequent class; otherwise the code \uv{?}; \emph{(3)}~the procedure \textbf{Group}
joins the values with the same code into groups --- at most as many groups arise as there are
classes, plus one (\uv{?}).

\subsection{Data reduction}

\subsubsection{Too many objects}

Batch processing becomes problematic. The remedies: \emph{(a)}~work only with a \textbf{sample}
of objects; \emph{(b)}~store the data in a way that allows access without having to keep it all
in main memory; \emph{(c)}~build \textbf{several models over subsets} and then combine them
(ensembles, Chapter~\ref{sec:metody}).

The key issue is the \textbf{representativeness} of the selected data --- the sample should
capture the nature of the data as well as possible, in particular as regards the correspondence
of the attribute value distributions on the sample and on the entire data.

\textbf{Unbalanced classes} in the original data lead to a tendency to prefer the majority
class. The remedies: different \emph{weights} for different types of misclassification
(cost-sensitive learning) or selecting objects from different classes with \emph{different
probabilities} (oversampling the minority, undersampling the majority).

\subsubsection{Too many attributes}

Two fundamentally different routes:

\begin{itemize}
\item \textbf{Reduction through transformation} --- from the existing attributes we create
      \emph{fewer new} ones, typically as linear combinations of the original ones (principal
      component analysis, PCA). Drawbacks: the values of \emph{all} the original attributes have
      to be known, and the new attributes may lack a clear interpretation.
\item \textbf{Reduction through selection} (\emph{feature selection}) --- from the existing
      attributes we choose only the most important ones. The interpretation is preserved and the
      unnecessary attributes need not be measured at all.
\end{itemize}

\mimo{} PCA seeks orthogonal directions of maximal variance (the eigenvectors of the covariance
matrix with the largest eigenvalues); it is an \emph{unsupervised} transformation, so it may
discard a direction that is crucial for prediction but has small variance.

\subsection{Attribute relevance analysis: filters and wrappers}
\label{ssec:filtr}

We are looking for those attributes that best contribute to a proper classification of objects.

\begin{defbox}[Filters vs.\ wrappers]
\begin{itemize}
\item \textbf{Filter} --- for each attribute compute a characteristic of its contribution to
      classification, order, select the best. \emph{Independent} of the model, fast.
\item \textbf{Wrapper} --- build a model over a \emph{subset} of the attributes and evaluate
      it; use the best subset. Expensive, but takes the interaction of the attributes
      \emph{with the concrete} model into account.
\item Searching the subsets: \emph{bottom-up} (forward selection --- adding attributes)
      vs.\ \emph{top-down} (backward elimination --- removing them).
\end{itemize}
\end{defbox}

\textbf{Limitation of filters:} each attribute is evaluated \emph{separately} $\Rightarrow$
attribute interactions are hard to capture (a criterion for whole sets runs into combinatorial
explosion). \mimo{} A filter thus lets through a pair of strongly correlated attributes and
discards an attribute useful only \emph{in combination} with another (an XOR relation).

\subsubsection{Criteria based on contingency tables}
\label{ssec:kontingencni}

The relation of an input attribute $A$ (values $v_1,\dots,v_R$) and a target attribute $C$
(classes $v_1,\dots,v_S$) is described by a contingency table with the frequencies $a_{kl}$ of
the combination $A(v_k)\wedge C(v_l)$, with the marginals
\[ r_k=\sum_{l=1}^{S}a_{kl},\qquad s_l=\sum_{k=1}^{R}a_{kl},\qquad
n=\sum_{k=1}^{R}\sum_{l=1}^{S}a_{kl}. \]
Probability estimates: $P(A(v_k)\wedge C(v_l))=a_{kl}/n$, $P(A(v_k))=r_k/n$,
$P(C(v_l))=s_l/n$. The expected frequency under independence is $e_{kl}=r_k s_l/n$.

\begin{defbox}[$\chi^2$ as a relevance criterion]
\[ \chi^2=\sum_{k=1}^{R}\sum_{l=1}^{S}\frac{(a_{kl}-e_{kl})^2}{e_{kl}}
   = n\sum_{k=1}^{R}\sum_{l=1}^{S}\frac{\big(a_{kl}-\tfrac{r_ks_l}{n}\big)^2}{r_ks_l}
   \qquad\text{(the higher the value, the better).}\]
An alternative form convenient for hand computation:
$\chi^2=\sum_k\sum_l \frac{a_{kl}^2}{e_{kl}}-n$.

\textbf{As a statistical test:} with the validity of the null hypothesis of independence
\[ H_0:\ P(A=v_k\wedge C=v_l)=P(A=v_k)\,P(C=v_l)\quad \forall k,l \]
the $\chi^2$ statistic has a distribution with $(R-1)(S-1)$ degrees of freedom. If
$\chi^2\ge\chi^2_{(R-1)(S-1)}(\alpha)$, the null hypothesis is rejected at the chosen level of
significance $\alpha$ and an alternative hypothesis of \emph{dependence} is accepted.
\end{defbox}

\begin{defbox}[Example --- $\chi^2$ test on a four-field table]
\begin{center}
\begin{tabular}{@{}llrrr@{}}
\toprule
& & \multicolumn{2}{c}{Provision of credit} & \\
& & YES & NO & $r_k$ \\
\midrule
Income height & high & 50 & 0 & 50 \\
              & low & 30 & 40 & 70 \\
\midrule
$s_l$ & & 80 & 40 & 120 \\
\bottomrule
\end{tabular}
\end{center}

Expected frequencies: $e_{11}=50\cdot 80/120=33.3$; $e_{12}=50\cdot 40/120=16.7$;
$e_{21}=70\cdot 80/120=46.7$; $e_{22}=70\cdot 40/120=23.3$. The differences are therefore
$+16.7$, $-16.7$, $-16.7$, $+16.7$ and the value of the statistic is $\chi^2=42.857$.

At the significance level $\alpha=0.05$ the critical value is
$\chi^2_{(1)}(0.05)=3.84$. Since $42.857>3.84$, \textbf{there is a dependence between the income
height and the provision of credit}.
\end{defbox}

\begin{defbox}[Fisher's exact test]
The $\chi^2$ test can be used only for sufficiently large frequencies --- it must hold that
$(r_k s_l)/n\ge 5$ for all $k,l$. For four-field tables with low frequencies \emph{Fisher's
test} is used. The probability that a four-field table has, for the given marginals
$r_1,r_2,s_1,s_2$, exactly the frequencies $a_{kl}$:
\[ p=\frac{r_1!\,r_2!\,s_1!\,s_2!}{n!\,a_{11}!\,a_{12}!\,a_{21}!\,a_{22}!}
     =\frac{\binom{s_1}{a_{11}}\binom{s_2}{a_{12}}}{\binom{n}{r_1}}. \]
The probabilities are summed over all considered values of the actual frequencies at the given
marginals (let $m=\min(r_1,s_1)$):
\[ P=\sum_{i=0}^{m}\frac{r_1!\,r_2!\,s_1!\,s_2!}{n!\,(a_{11}-i)!\,(a_{12}+i)!\,(a_{21}+i)!\,(a_{22}-i)!}. \]
If $P\le\alpha$, the null hypothesis is rejected at the level of significance $\alpha$.

\textbf{Example.} For a table with the marginals $r=(12,18)$, $s=(6,24)$, $n=30$ we get
$P_0=0.031$ and $P_1=0.173$, hence $P=P_0+P_1=0.204>0.05$ --- \textbf{the null hypothesis of
independence is accepted}. (Check: the sum $P_0,\dots,P_6$
$=0.031+0.173+0.340+0.302+0.128+0.024+0.002=1.000$.)
\end{defbox}

\subsubsection{Criteria based on entropy}

\begin{defbox}[Entropy, mutual information, information measure of dependence]
The \textbf{entropy} of an attribute $A$ as the weighted average of the entropies of its values
(the \emph{smaller}, the better):
\[ H(A)=\sum_{k=1}^{R}\frac{r_k}{n}\,H\big(A(v_k)\big),\qquad
   H\big(A(v_k)\big)=-\sum_{l=1}^{S}\frac{a_{kl}}{r_k}\log_2\frac{a_{kl}}{r_k}. \]

\textbf{Mutual information:}
\[ \mathrm{MI}(A,C)=\sum_{k=1}^{R}\sum_{l=1}^{S}\frac{a_{kl}}{n}
   \log_2\frac{n\,a_{kl}}{r_k s_l}. \]

The \textbf{information measure of dependence} --- normalised mutual information (the
\emph{bigger}, the better):
\[ \mathrm{ID}(A,C)=\frac{\mathrm{MI}(A,C)}{H(C)}
   =\frac{\mathrm{MI}(A,C)}{-\sum_{l=1}^{S}\frac{s_l}{n}\log_2\frac{s_l}{n}}. \]
\end{defbox}

The normalisation in $\mathrm{ID}$ is essential: $\mathrm{MI}$ alone systematically favours
attributes with many values (in the limit, the object identifier is the \uv{most informative}
attribute). The same bias is addressed by the information gain ratio in C4.5
(Chapter~\ref{sec:metody}).

\subsubsection{Criteria for numeric attributes: correlation}
\label{ssec:korelace}

The criteria above ($\chi^2$, entropy) work with \emph{categorical} attributes over a contingency
table. For a pair of \textbf{numeric} attributes, \emph{correlation analysis} is used instead ---
and it serves two purposes: measuring the \textbf{relevance} of an input with respect to
a~numeric output, and detecting \textbf{redundancy} between a~pair of inputs
(see Sect.~\ref{ssec:integrace}).

\begin{defbox}[Pearson correlation coefficient]
A~measure of the \textbf{linear} relationship between two numeric quantities; it is the sample
covariance normalised by the standard deviations:
\[ r=\frac{\mathrm{cov}(X,Y)}{s_X\,s_Y}
   =\frac{\sum_{i=1}^{n}(x_i-\bar x)(y_i-\bar y)}
          {\sqrt{\sum_{i=1}^{n}(x_i-\bar x)^2\ \sum_{i=1}^{n}(y_i-\bar y)^2}}\ \in[-1,1]. \]
$r=\pm 1$ means the points lie \emph{exactly} on a~line (with a~positive, resp.\ negative
slope); $r=0$ means there is no \emph{linear} relationship between them. The square $r^2$ (the
\emph{coefficient of determination}) gives the proportion of the variance of $Y$ explained by
the linear dependence on~$X$.

\textbf{Relation to regression} (Sect.~\ref{ssec:regrese}): the slope is
$\beta=\mathrm{cov}(X,Y)/s_X^2$, hence $r=\beta\,s_X/s_Y$ --- the same information, merely
scaled to be dimensionless.
\end{defbox}

\begin{poznbox}[Three pitfalls of the Pearson coefficient]
\textbf{(1) It measures only a~linear relationship.} For data distributed symmetrically around
zero and $y=x^2$ we get $r=0$, even though $Y$ depends on $X$ \emph{functionally}. Zero
correlation therefore \textbf{does not imply independence} (the converse does hold:
independence $\Rightarrow r=0$).

\textbf{(2) It is sensitive to outliers.} A~single extreme point can drive $r$ towards one even
where the rest of the data shows no trend at all.

\textbf{(3) Correlation is not causation.} The relationship may be produced by a~hidden third
variable; distinguishing them requires an experiment or a~temporal ordering (cf.\ influence
versus correlation in Chapter~\ref{sec:sna}).
\end{poznbox}

\begin{defbox}[Spearman rank correlation coefficient $\rho$]
It is the \textbf{Pearson coefficient computed from the ranks} instead of the values. If there
are no tied ranks in the data, it simplifies to
\[ \rho = 1-\frac{6\sum_{i=1}^{n} d_i^2}{n\,(n^2-1)}\ \in[-1,1], \]
where $d_i$ is the difference of the ranks of the $i$-th object in $X$ and in $Y$.

\textbf{What this buys you:}
\begin{itemize}
\item It captures \emph{any \textbf{monotone}} relationship, not merely a~linear one.
\item It is \textbf{robust to outliers} --- ranking compresses the extremes.
\item It is \textbf{invariant under monotone transformations} ($\log$, $\sqrt{\ }$, \dots);
      Pearson is not.
\item It also works for \textbf{ordinal} attributes, whose values can be ordered but not
      subtracted (grades, satisfaction levels).
\end{itemize}
\end{defbox}

\textbf{A~numerical example.} Five objects, $x=(1,2,3,4,5)$ and $y=(1,2,4,8,16)$ --- the
relationship is \emph{perfectly monotone}, yet exponential, i.e.\ non-linear. The ranks of both
variables coincide, so all $d_i=0$ and $\rho=1$ (a~perfect match). Pearson, by contrast, gives
$r\doteq 0{,}93$: it \emph{penalises} the non-linearity. Had we taken the logarithm of $y$
($\log_2 y = 0,1,2,3,4$), Pearson would give $r=1$ as well --- which is exactly the invariance
at work: Spearman returns $1$ straight away, Pearson only after the right transformation.

\textbf{Which measure to choose.} Pearson where the relationship is approximately linear and the
data are free of outliers (and where we want to read $r^2$ as the explained variance); Spearman
where only the \emph{strength and direction of the trend} matter, where the data contain
extremes, or where the attributes are merely ordinal. \mimo{} A~third common variant is
\emph{Kendall's $\tau$}, which counts the proportion of concordant and discordant pairs instead
of rank differences.

\textbf{Use in data reduction.} A~\emph{correlation matrix} of all pairs of input attributes is
built; pairs with $|r|$ close to 1 carry practically the same information and one of them can be
removed. Mind, however, the caveat from Sect.~\ref{ssec:filtr}: pairwise correlation alone cannot
recognise an attribute that is useful only \emph{in combination} with another.

\subsection{Relations between attributes: regression}
\label{ssec:regrese}

Whereas contingency tables describe the relation of \emph{categorical} attributes, for numeric
ones regression analysis is used: \emph{linear regression} determines the \emph{parameters} of
a linear relationship between two numerical quantities.

\textbf{Reasons for regression analysis:} \emph{(1)}~it is very costly to measure the output
$\Rightarrow$ predict it from inputs that are easily available; \emph{(2)}~the input values are
known earlier than the output $\Rightarrow$ estimate the output; \emph{(3)}~controlled input
values might help find the correct behaviour of the output; \emph{(4)}~there can be a causal
relationship between input and output which we want to find.

\begin{defbox}[Linear regression for one input variable by least squares]
Patterns $[x_1,y_1],\dots,[x_p,y_p]$, regression equation $y=\beta x+\alpha$, quadratic error
\[ E=\sum_{i=1}^{p}\varepsilon_i^2=\sum_{i=1}^{p}\big(y_i-\hat y_i\big)^2
   =\sum_{i=1}^{p}\big(y_i-\beta x_i-\alpha\big)^2. \]
Setting the partial derivatives to zero:
\[ \frac{\partial E}{\partial\alpha}=-2\sum_i\big(y_i-\beta x_i-\alpha\big)=0,\qquad
   \frac{\partial E}{\partial\beta}=-2\sum_i\big(y_i-\beta x_i-\alpha\big)x_i=0, \]
whence, after rearranging, $\alpha=\bar y-\beta\bar x$ and
\[ \beta=\frac{\sum_{i=1}^{p}(x_i-\bar x)(y_i-\bar y)}{\sum_{i=1}^{p}(x_i-\bar x)^2}
   \qquad(\text{i.e.\ the sample covariance divided by the sample variance in }x). \]
Prediction: $\hat y=\beta x+\alpha$.
\end{defbox}

\textbf{Multi-dimensional linear regression.} We assume a linear dependence of the explained
quantity $Y$ on several explaining quantities $X_1,\dots,X_m$:
$\hat y=\beta_0+\beta_1x_1+\cdots+\beta_mx_m$, in matrix notation $\hat Y=X\beta$, where $X$ is
the extended matrix of input patterns (the first column being ones) and
$\beta=(\beta_0,\dots,\beta_m)$. The quadratic deviation is
$E=(Y-X\beta)^{\!\top}(Y-X\beta)$; from $\partial E/\partial\beta=0$ we obtain the normal
equations $X^{\!\top}\!X\beta=X^{\!\top}Y$, hence
\[ \beta=\big(X^{\!\top}\!X\big)^{-1}X^{\!\top}Y. \]
For complex practical tasks this computation is expensive $\Rightarrow$ approximative
(iterative) solutions are used.

\textbf{Non-linear regression} assumes a more complex functional dependence (quadratic,
exponential, \dots). Its most important case is \emph{logistic regression}, which models a
categorical (two-valued) output --- since it is at the same time a classifier, it is treated in
Chapter~\ref{sec:metody}, Section~\ref{ssec:logreg}.

\subsection{Summary for the examination}

\begin{itemize}
\item \textbf{The goal of data preparation}: select the relevant data, represent it for the
      chosen algorithm, the result being \emph{one} table of objects $\times$ attributes.
\item \textbf{Missing values}: ignore the object / a new value \uv{I don't know} / an existing
      value (mode, proportionally, arbitrarily) / impute with a model. Be able to say when each
      choice is appropriate (absence as information $\to$ \uv{I don't know}).
\item \textbf{Outliers}: $\mathrm{IQR}=Q_3-Q_1$, an outlier lies outside
      $[Q_1-1.5\,\mathrm{IQR},\,Q_3+1.5\,\mathrm{IQR}]$. Be able to compute the quartiles from
      ascending data (the $n/4$ and $3n/4$ rule, rounding up) and to draw a box plot.
\item \textbf{Standardisation}: decimal scaling, min-max to $[0,1]$ and $[-1,1]$, z-score by the
      mean absolute difference (robust to outliers). The
      motivation: without it the distance is dominated by the attribute with the largest range.
\item \textbf{Discretisation}: equi-width (same length), equi-depth (same number of objects),
      using the class (entropy top-down by splitting, $\chi^2$ bottom-up by merging);
      \textbf{fuzzy discretisation} --- the procedure Interval, blurred boundaries, two
      categories per value with characteristic functions summing to 1, the weight of an
      object is $\prod_j\mu_j$. \textbf{KEX} groups the values of a categorical attribute
      (Assign: a class code by the $\chi^2$ test against the a-priori class distribution,
      otherwise \uv{?}; Group: values with the same code).
\item \textbf{Data reduction}: too many objects $\to$ sampling (representativeness!), efficient
      storage, several models over subsets; unbalanced classes $\to$ error weights or different
      selection probabilities. Too many attributes $\to$ \emph{transformation} (PCA: fewer new
      attributes as linear combinations, but all the original ones must be known and the
      interpretation is lost) vs.\ \emph{selection} (a subset of the original ones, the
      interpretation is preserved).
\item \textbf{Filters vs.\ wrappers}: a filter computes a characteristic for each attribute
      separately and is independent of the model (fast, but blind to interactions); a wrapper
      builds a model over a subset and evaluates it (expensive, takes interactions into account).
      Search \emph{bottom-up} (adding) or \emph{top-down} (removing).
\item \textbf{Filter criteria}: $\chi^2$ (bigger $=$ better), entropy $H(A)$ (smaller $=$
      better), $\mathrm{MI}(A,C)$, $\mathrm{ID}(A,C)=\mathrm{MI}/H(C)$ (bigger $=$ better). Be
      able to justify the normalisation in $\mathrm{ID}$: without it, attributes with many
      values are favoured.
\item \textbf{The $\chi^2$ test}: $(R-1)(S-1)$ degrees of freedom, the condition
      $r_ks_l/n\ge 5$; for low frequencies \textbf{Fisher's exact test}. Be able to go through
      the computation on a four-field table (the credit example: $\chi^2=42.857>3.84$
      $\Rightarrow$ dependence).
\item \textbf{Regression}: the least squares method, $\beta$ as covariance
      divided by variance, $\alpha=\bar y-\beta\bar x$; multi-dimensionally
      $\beta=(X^{\!\top}\!X)^{-1}X^{\!\top}Y$.
\end{itemize}

%=====================================================================
\newpage
\section{Methods for data mining}
\label{sec:metody}
%=====================================================================

The most extensive sentence of the topic. It names three groups of methods that cover three
different situations: \emph{association rules} for transactional data without a target attribute,
\emph{supervised learning} for data with a target attribute, and \emph{cluster analysis} for
metric data without a target attribute. The chapter is divided into three sections accordingly.

\subsection{Association rules and frequent patterns}
\label{ssec:ar}

\subsubsection{Market basket analysis}

\emph{Market basket analysis} (MBA) is the archetypal task: \uv{which items occur together in the
basket?} The results are expressed in the form of rules and can be applied immediately --- in
planning and the layout of shops, in offering coupons and time-limited discounts, or in
\uv{packaging} products.

An \textbf{item} is a product or service on offer, a \textbf{transaction} contains one or more
items. The \textbf{frequency table} indicates the number of co-occurrences of any two items
within the transactions; the values on its diagonal correspond to the number of transactions
that contain the respective item.

\begin{defbox}[Example --- frequency table]
Transactions in a grocery shop:
\begin{center}
\small
\begin{tabular}{@{}rl@{\hspace{1.2cm}}lrrrr@{}}
\toprule
Cust. & Items & & bread & butter & milk & coffee \\
\midrule
1 & bread, butter          & bread & 4 & 3 & 1 & 2 \\
2 & milk, bread, butter    & butter & 3 & 4 & 1 & 2 \\
3 & bread, coffee          & milk & 1 & 1 & 1 & 0 \\
4 & bread, butter, coffee  & coffee & 2 & 2 & 0 & 3 \\
5 & coffee, butter         &       &   &   &   &   \\
\bottomrule
\end{tabular}
\end{center}
The character of the sales is apparent from the frequency table at once: \emph{bread and butter
tend to be bought together}, \emph{milk is never purchased together with coffee}.
\end{defbox}

The rules have the form \texttt{IF} \emph{condition} \texttt{THEN} \emph{conclusion}. Association
rules should be:
\begin{itemize}
\item \textbf{easy to understand} --- once a relationship is found, it can be easily verified;
\item \textbf{applicable} --- they comprise useful information that can lead to further
      interventions;
\end{itemize}
and conversely they should \emph{not} be \textbf{trivial} (everyone already knows the result) nor
\textbf{inexplicable} (there is no explanation for them and they do not lead to an action).

\subsubsection{Support, confidence and lift}

\begin{defbox}[Support, confidence, lift]
\textbf{Support} --- how often can we use the rule:
\[ \mathrm{sup}(R)=\frac{\text{number of transactions containing } i \text{ and } j}
   {\text{number of all transactions}}\cdot 100\,\%. \]

\textbf{Confidence} --- how reliable are the results of the rule:
\[ \mathrm{conf}(R)=\frac{\text{number of transactions containing } i \text{ and } j}
   {\text{number of transactions containing } i}\cdot 100\,\%. \]

\textbf{Lift} --- how many times is it better to use the rule for prediction instead of just
assuming its result:
\[ \mathrm{lift}(R)=\frac{p(i\wedge j)}{p(i)\,p(j)}. \]
If $\mathrm{lift}<1$, the rule is \emph{worse than a random choice} for prediction and the
\textbf{negation of the conclusion} may yield a better rule (\texttt{IF} \emph{condition}
\texttt{THEN NOT} \emph{conclusion}).
\end{defbox}

\begin{defbox}[Example --- support, confidence, lift]
Over the data of the previous example (5 transactions):

\emph{Rule 1:} \uv{if the customer buys bread then he buys also butter.}
$\mathrm{sup}=3/5=60\,\%$, $\mathrm{conf}=3/4=75\,\%$.

\emph{Rule 2:} \uv{if the customer buys coffee then he buys also butter.}
$\mathrm{sup}=2/5=40\,\%$, $\mathrm{conf}=2/3=66\,\%$.

\emph{Rule 3:} \uv{if the customer buys milk then he buys also butter.}
$\mathrm{sup}=1/5=20\,\%$, $\mathrm{conf}=1/1=100\,\%$,
$\mathrm{lift}=\frac{1/5}{(1/5)\cdot(4/5)}=\frac54=1.25$.

Rule 3 has \emph{one hundred percent} confidence but a support of only 20~\% --- it rests on a
single transaction. The lift of 1.25 says that the gain over chance is only small (80~\% of the
customers buy butter anyway). This is the basic trap: \emph{high confidence alone is not enough}.
\end{defbox}

\subsubsection{Choice of items, taxonomies and virtual items}

Transaction data is often of worse quality and requires a more extensive preprocessing; moreover
the relevance of the items changes with time. The adequate level of processing is a compromise
between a growing number of combinations and the immediate applicability of the results (specific
items), while keeping a sufficiently high support.

A \textbf{taxonomy} is a hierarchy of categories. In the beginning more general items are used;
later, rules are generated for specific items, and only from the transactions that contain them.
For the results to be applicable, the items considered should occur in roughly the same number of
transactions: \emph{rare items are moved to higher levels} of the taxonomy (where they appear
relatively more frequently) and \emph{the more common items are left at lower levels} (in order
not to allow the most frequent items to dominate the rules).

\textbf{Virtual items} exceed the framework of a taxonomy: brands across product types (Calvin
Klein) or information about the transaction itself (day, time, customer). The risk is
\emph{redundancy} in the rules (a virtual item corresponds to exactly one taxonomy item, or
meets the general item in one rule). A typical use is the \emph{comparison of shops}: a virtual
item marks the type of the shop, MBA is applied to comparably sized data from the new and the
existing shops, and one watches the rules in which the virtual items appear.

\textbf{Pruning according to minimum support} eliminates the less frequent items --- the actions
undertaken should concern enough transactions. Two options: eliminate the rare items (and
subsequently the corresponding rules), or use a taxonomy to create general items (which must meet
the set threshold). A variable minimum support leads to a \emph{cascade effect}.

\subsubsection{The Apriori algorithm}

\begin{defbox}[Frequent itemset and the Apriori principle]
A \emph{frequent itemset} is a combination (conjunction) of categories that reaches a preset
frequency on the data --- the support \emph{minsup}.

The \textbf{downward closure property} --- the key idea for Apriori: \emph{any subset of a
frequent itemset is also a frequent itemset}. Equivalently: if an itemset is infrequent, so are
all its supersets.

When looking for combinations of length $k$ we therefore use the known combinations of length
$k-1$ --- a \emph{breadth-first} generation. To generate a combination of length $k$ we demand
that \emph{all} its subcombinations of length $k-1$ fulfil the frequency requirement.
\end{defbox}

\begin{defbox}[Algorithm APRIORI (R.~Agrawal)]
\begin{algorithmic}[1]
\State $L_1\gets$ all categories that reach at least the required frequency
\State $k\gets 2$
\While{$L_{k-1}\neq\emptyset$}
  \State $C_k\gets \textsc{Apriori-Gen}(L_{k-1})$ \Comment{candidate generation}
  \State $L_k\gets$ those combinations from $C_k$ that reached at least the required frequency
  \State $k\gets k+1$
\EndWhile
\end{algorithmic}

\textbf{The function \textsc{Apriori-Gen}$(L_{k-1})$:}
\begin{algorithmic}[1]
\For{all pairs of combinations $\mathit{Comb}_p,\mathit{Comb}_q\in L_{k-1}$}
  \If{$\mathit{Comb}_p$ and $\mathit{Comb}_q$ match in $k-2$ categories}
     \State add $\mathit{Comb}_p\cup \mathit{Comb}_q$ to $C_k$ \Comment{the join step}
  \EndIf
\EndFor
\For{each combination $\mathit{Comb}\in C_k$} \Comment{the pruning step}
  \If{any of its subcombinations of length $k-1$ is not contained in $L_{k-1}$}
     \State remove $\mathit{Comb}$ from $C_k$
  \EndIf
\EndFor
\end{algorithmic}
\end{defbox}

\textbf{Rule generation.} After finding the combinations with a satisfactory frequency, the
association rules are created (note that \emph{supercombination frequency} $\le$
\emph{combination frequency}). Each combination $\mathit{Comb}$ is divided into all possible
pairs of subcombinations $\mathit{Ant}$ and $\mathit{Suc}$ such that
\[ \mathit{Suc}=\mathit{Comb}\setminus \mathit{Ant},\qquad
   \mathit{Ant}\cap \mathit{Suc}=\emptyset,\quad \mathit{Ant}\cup \mathit{Suc}=\mathit{Comb}. \]
Then, for the rule $\mathit{Ant}\Rightarrow \mathit{Suc}$, the \emph{support} corresponds to the
frequency of the combination $\mathit{Comb}$ and the \emph{confidence} to the ratio of the
frequencies of the combinations $\mathit{Comb}$ and $\mathit{Ant}$.

\begin{center}
\small
\begin{tabular}{@{}p{7.6cm}p{8.0cm}@{}}
\toprule
\textbf{Advantages of MBA} & \textbf{Drawbacks of MBA} \\
\midrule
clear and understandable results --- IF-THEN rules that are immediately applicable &
  computational costs: the space of association rules is exponential, $O(2^m)$ for $m$ items \\
mining without target output values (without a-priori knowledge) & taxonomies and virtual items
  are necessary \\
can process data of a variable length & difficulty in determining an adequate number of items;
  the items should have about the same frequency \\
simple and easy-to-understand computation & \emph{it tends to produce excessive numbers of
  rules} --- thousands, tens of thousands, millions \\
\bottomrule
\end{tabular}
\end{center}

\subsubsection{Multiple minimum supports: MS-Apriori}

\textbf{The rare item problem:} a single \emph{minsup} does not fit frequent and rare items at
the same time (bread vs.\ a food processor) --- a high one misses rules with rare items, a low
one explodes combinatorially on the frequent ones.

\begin{defbox}[The main difference from Apriori]
Each item has its \emph{own minimum support} $\mathrm{MIS}(i)$ (\emph{minimum item support});
\textbf{the minsup of a rule is the minimum of the MIS values of its items}. To prevent very
frequent and very rare items from mixing in one itemset, the support difference is also
bounded: $\max_{i\in s}\mathrm{sup}(i)-\min_{i\in s}\mathrm{sup}(i)\le\varphi$.
\end{defbox}

\textbf{Consequence: downward closure no longer holds.} Counterexample:
$\mathrm{MIS}(1)=10\,\%$, $\mathrm{MIS}(2)=20\,\%$, $\mathrm{MIS}(3)=5\,\%$; the itemset
$\{1,2\}$ with support 9~\% is infrequent, but $\{1,2,3\}$ could be frequent --- $\{1,2\}$
cannot be discarded. MS-Apriori copes with this by three modifications: \emph{(a)}~it sorts
the items in \emph{ascending order of MIS} and tests the frequency of an itemset against the
MIS of its \emph{first} item ($c.\mathrm{count}/n\ge\mathrm{MIS}(c[1])$); \emph{(b)}~it
generates the level-2 candidates from the set of seeds $L$ (items satisfying the MIS of some
earlier item), not from $F_1$; \emph{(c)}~for rule generation it additionally counts
\texttt{tailCount} --- the support of the itemset without its first item --- because the
supports of the subsets are no longer guaranteed to be available. The single-minsup model is a
special case ($\mathrm{MIS}=100\,\%$ prohibits rules with the item).

\subsubsection{Class association rules (CAR)}

Usual mining finds all rules; sometimes only rules with a \emph{target class} in the
consequent are of interest (which words are associated with which topic of a document).

\begin{defbox}[Class association rule]
The transactions are additionally labelled with classes; a \emph{class association rule}
(CAR) is an implication $X\to y$, where $X\subseteq I$ (items) and $y\in Y$ (classes),
$I\cap Y=\emptyset$. Support and confidence are defined the same way as for usual association
rules. E.g.\ on a collection of documents:
Student, School $\to$ Education $[\mathrm{sup}=2/7,\ \mathrm{conf}=2/2]$.
\end{defbox}

CARs are mined \textbf{directly in one step} (by a modified Apriori) via \emph{rule-items}
$(\mathit{condset},y)$ with a support above \emph{minsup}. \emph{Different minsups for
different classes} can be specified (the multiple-supports idea); a class minsup of 101~\%
prohibits generating its rules.

\subsubsection{Sequential pattern mining}

Association rules ignore the \emph{order} of the transactions; the order matters e.g.\ in MBA
(a bed first, bed sheets later) or when looking for navigational patterns in web logs
(\emph{web usage mining}).

\begin{defbox}[Sequence, size, length, subsequence]
A \emph{sequence} $s=\langle a_1a_2\dots a_r\rangle$ is an ordered list of non-empty itemsets
(\emph{elements}). The \emph{size} $=$ the number of elements, the \emph{length} $=$ the
number of items ($k$-sequence). $s_1=\langle a_1\dots a_r\rangle$ is a \emph{subsequence} of
$s_2=\langle b_1\dots b_v\rangle$ if there exist indices $1\le j_1<\cdots<j_r\le v$ such that
$a_1\subseteq b_{j_1},\dots,a_r\subseteq b_{j_r}$. Beware: $\langle\{3\}\{8\}\rangle$ and
$\langle\{3,8\}\rangle$ do \emph{not} contain each other.
\end{defbox}

\textbf{The task:} find all sequences with a support (the fraction of the data sequences
containing them) of at least \emph{minsup} --- the \emph{sequential patterns}. The
\textbf{GSP} algorithm (\emph{Generalized Sequential Pattern}) works similarly to Apriori:
breadth-first, candidates of length $k$ by joining frequent sequences of length $k-1$.
Add-ons: a cause--effect analysis by introducing items for the transactions \emph{before} and
\emph{after} an event; \emph{dissociation rules} (IF $A$ AND NOT $B$) via complementary items
--- at the cost of twice as many items.

\subsubsection{FP-Growth and the FP-tree}

Apriori uses a \emph{generate-and-test} approach: candidate generation is expensive (both in
space and time) and support counting is expensive too --- subset checking is computationally
expensive and the database is scanned \emph{repeatedly} (I/O). \textbf{FP-Growth} allows frequent
itemset discovery \emph{without candidate generation} in two steps: \emph{(1)}~build a compact
data structure, the \emph{FP-tree} (using two passes over the data); \emph{(2)}~extract the
frequent itemsets directly from the FP-tree by traversing it.

\begin{defbox}[The FP-tree]
The nodes of the FP-tree correspond to items and have a counter. FP-Growth reads one transaction
at a time and maps it to a \emph{path} in the tree. A \emph{fixed order} of items is used, so
paths can overlap when the transactions share a \emph{prefix}; in that case the counters are
incremented. Pointers are maintained between the nodes containing the same item, creating singly
linked lists. The more paths overlap, the higher the compression, and the more likely the FP-tree
fits in memory.
\end{defbox}

\textbf{Construction (2 passes).} \emph{Pass 1:} scan the data and find the support of each item,
discard the infrequent items, sort the frequent ones in \emph{decreasing} order of support (e.g.\
$a,b,c,d,e$) --- this order is used when building the tree so that common prefixes can be shared.
\emph{Pass 2:} construct the tree. E.g.\ transaction $\{a,b\}$ creates the path
\texttt{null}$\to a\to b$ with counters 1; transaction $\{b,c,d\}$ creates a \emph{disjoint} path
\texttt{null}$\to b\to c\to d$ (it shares no common prefix with the first one), and a link is
added between the two $b$ nodes; transaction $\{a,c,d,e\}$ shares the common prefix $a$ with the
first one, so the paths overlap and the counter of node $a$ is incremented.

\textbf{The size of the FP-tree.} Usually smaller than the uncompressed data --- transactions
typically share items, and hence prefixes. \emph{Best case:} all transactions contain the same set
of items $\Rightarrow$ one path. \emph{Worst case:} every transaction has a unique set of items
$\Rightarrow$ the tree is at least as large as the original data and the storage requirements are
\emph{higher} (the pointers and counters have to be stored). The size depends on how the items
are ordered --- decreasing support is only a \emph{heuristic} and does not always lead to the
smallest tree.

\begin{defbox}[Extracting frequent itemsets from the FP-tree]
The algorithm proceeds \emph{bottom-up}, from the leaves towards the root, by \emph{divide and
conquer}: it first looks for frequent itemsets ending in $e$, then in $de$, \dots, then $d$,
$cd$, and so on.

\begin{enumerate}
\item The \emph{prefix path sub-tree} ending in a given item is obtained using the linked lists.
\item Check whether the item is frequent --- add the counts along the linked list. (In the example
      with $\mathit{minSup}=2$ we get $\mathrm{count}(e)=3$, so $\{e\}$ is extracted as a frequent
      itemset.)
\item If the item is frequent, find recursively the frequent itemsets ending in $de$, $ce$, $be$,
      $ae$. For that, the \emph{conditional FP-tree} must be obtained.
\end{enumerate}

The \textbf{conditional FP-tree} for the item $e$ is the tree that would be built if we only
considered the \emph{transactions containing} $e$ and then removed $e$ from all transactions. It
is obtained from the prefix path sub-tree in three steps:
\begin{enumerate}
\item update the support counts along the prefix paths to reflect the number of transactions
      containing $e$;
\item remove the nodes containing $e$ (the information about them is no longer needed);
\item remove the infrequent items from the prefix paths --- e.g.\ if $b$ has a support of 1
      (which really means that $be$ has a support of 1), then $be$ is infrequent and $b$ can be
      removed.
\end{enumerate}
The conditional tree is then used to find the itemsets ending in $de$, $ce$, $ae$ ($be$ is not
considered, because $b$ is not in the conditional tree); for each of them the prefix paths are
obtained again, frequent itemsets are extracted and a conditional tree is generated ---
recursively (e.g.\ $e\to de\to ade$).
\end{defbox}

\begin{center}
\small
\begin{tabular}{@{}p{7.6cm}p{8.0cm}@{}}
\toprule
\textbf{Advantages of FP-Growth} & \textbf{Disadvantages of FP-Growth} \\
\midrule
only 2 passes over the data set & the FP-tree may not fit in memory \\
it \uv{compresses} the data set & the FP-tree is expensive to build (but afterwards the frequent
  itemsets are read off easily) \\
no candidate generation, much faster than Apriori & pruning can only be done on single items;
  the support can be computed only after \emph{all} the data has been added to the tree \\
\bottomrule
\end{tabular}
\end{center}

\mimo{} As a supplement: a \emph{closed} frequent itemset is one no proper superset of which has
the same support; a \emph{maximal} one is one no proper superset of which is frequent. Maximal
itemsets are the most compact representation (but they lose the information about the supports of
the subsets), whereas closed itemsets are a compromise --- a lossless representation of all
supports.

\subsection{Approaches based on the principle of supervised learning}
\label{ssec:supervised}

\subsubsection{Formulation of the task}

\begin{defbox}[Classification and prediction]
A training set $\{(x_1,y_1),\dots,(x_p,y_p)\}$ is given, where $x_i$ is a vector of input
attribute values and $y_i$ the value of the \emph{target} attribute. The task is to find a model
$\hat y=f(x)$ that predicts the target attribute well also for previously unseen inputs.

If the target attribute is \emph{discrete}, this is \textbf{classification}; if it is
\emph{continuous}, \textbf{regression} (prediction). The solution always has two phases:
\emph{(1)}~\textbf{induction} (learning, training) of the model from the training data;
\emph{(2)}~\textbf{application} of the model to new data.
\end{defbox}

It is essential that the model be evaluated on \emph{new} data --- the ability to memorise the
training data is worthless (\emph{overfitting}). The methods for error estimation and the quality
measures are in Chapter~\ref{sec:reprez}, Section~\ref{ssec:eval}.

\subsubsection{Decision trees}

\begin{defbox}[Decision tree]
Let $D$ be a database $\{d_1,\dots,d_p\}$, $\{A_1,\dots,A_n\}$ a set of attributes and
$C=\{C_1,\dots,C_m\}$ a set of class labels. A \emph{decision tree} for $D$ is a tree where:
\begin{itemize}
\item each inner node is labelled by an attribute $A_a$;
\item each edge is labelled by a predicate (applicable to the attribute of the parent node);
\item each leaf has a class label $C_j$.
\end{itemize}
The tree divides the feature space into \emph{rectangular regions}; patterns are classified
according to the region they belong to.
\end{defbox}

\begin{center}
\small
\begin{tabular}{@{}p{7.4cm}p{8.2cm}@{}}
\toprule
\textbf{Advantages} & \textbf{Disadvantages} \\
\midrule
easy to implement and interpret & difficult to process continuous data (attribute categorisation
  needed $\to$ rectangular regions) \\
efficient & cannot always be used --- the class boundary may be diagonal and an axis-parallel
  tree can only approximate it in steps (in the bank-loan example the boundary runs across the
  \emph{account balance} $\times$ \emph{income} plane) \\
extraction of simple rules & difficult to process missing data \\
applicable also to big data sets & over-fitting $\Rightarrow$ \emph{pruning} required \\
  & mutual correlations among the attributes are not considered \\
\bottomrule
\end{tabular}
\end{center}

\textbf{Terminology:} the \emph{decision attributes} for the division of the data set label the
decision nodes of the tree; the \emph{decision predicates} label the edges; the \emph{stopping
criterion} is e.g.\ \uv{all patterns from the reduced set belong to the same class}.

\begin{defbox}[The TDIDT algorithm (Top Down Induction of Decision Trees)]
Induction of the tree top-down by \emph{divide and conquer}:
\begin{enumerate}
\item Choose one attribute as the \emph{root} of the partial (sub)tree.
\item Divide the data arriving at this node into subsets according to the values of the chosen
      attribute and add a node for each subset.
\item If there is a node such that the data patterns arriving at it do not all belong to the same
      class, repeat the process for it starting from step 1; else stop.
\end{enumerate}
\end{defbox}

\begin{defbox}[Entropy and the choice of the splitting attribute]
The objective is to choose the attribute that would best divide the patterns from different
classes. \emph{Entropy} as a measure of disorder (\uv{surprise}) in the system:
\[ H=-\sum_{j=1}^{m}p_j\log_2 p_j, \]
where $p_j$ is the probability of occurrence of class $j$ (its relative frequency on the set of
patterns considered) and $m$ the number of classes. Convention: $0\log_2 0=0$; conversion
$\log_2 x=\log_{10}x/0.301$.

For an attribute $A$: for each of its possible values $v$ compute the entropy $H(A(v))$ on the
set of examples covered by the category $A(v)$,
\[ H\big(A(v)\big)=-\sum_{j=1}^{m}p_j^{A(v)}\log_2 p_j^{A(v)}, \]
and the \emph{weighted average entropy} of the attribute as the weighted sum
\[ H(A)=\sum_{v}\frac{|A(v)|}{|D|}\,H\big(A(v)\big). \]
The attribute with the \emph{minimum} $H(A)$ is chosen.
\end{defbox}

\begin{defbox}[Example --- an application for a loan]
\begin{center}
\small
\begin{tabular}{@{}lllllc@{}}
\toprule
Client & Income & Account & Gender & Unemployed & Loan approved \\
\midrule
C1 & high & high & female & no & yes \\
C2 & high & high & male & no & yes \\
C3 & low & low & male & no & no \\
C4 & low & high & female & yes & yes \\
C5 & low & high & male & yes & yes \\
C6 & low & low & female & yes & no \\
C7 & high & low & male & no & yes \\
C8 & high & low & female & yes & yes \\
C9 & low & middle & male & yes & no \\
C10 & high & middle & female & no & yes \\
C11 & low & middle & female & yes & no \\
C12 & low & middle & male & no & yes \\
\bottomrule
\end{tabular}
\end{center}

\textbf{Step 1 --- choosing the attribute by minimum entropy.} The four-field table for
\emph{income} and \emph{loan approved}: high income 5 yes / 0 no, low income 3 yes / 4 no.
\begin{align*}
H(\text{income}(\text{high}))&=-\tfrac55\log_2\tfrac55-\tfrac05\log_2\tfrac05=0,\\
H(\text{income}(\text{low}))&=-\tfrac37\log_2\tfrac37-\tfrac47\log_2\tfrac47=0.9852,\\
H(\text{income})&=\tfrac{5}{12}\cdot 0+\tfrac{7}{12}\cdot 0.9852=0.5747.
\end{align*}
Analogously $H(\text{account})=\tfrac13\cdot 0+\tfrac13\cdot 1+\tfrac13\cdot 1=0.6667$,
$H(\text{gender})=\tfrac12\cdot 0.9183+\tfrac12\cdot 0.9183=0.9183$ and
$H(\text{unemployed})=\tfrac12\cdot 1+\tfrac12\cdot 0.65=0.825$.
The minimum belongs to \textbf{income} $\Rightarrow$ the first branching.

\textbf{Step 2.} Income(high): loan approved(yes) for all patterns $\to$ a leaf. Income(low):
7 clients, further branching; $H(\text{account})=\tfrac27 H(\dots)+\tfrac37 H(\dots)+\tfrac27
H(\dots)=0.3935$, $H(\text{gender})=0.965$, $H(\text{unemployed})=0.9792$ $\Rightarrow$ branch on
\textbf{account}.

\textbf{Step 3.} Account(high): yes for all; account(low): no for all; account(middle): 3
clients, $H(\text{gender})=\tfrac23\cdot 1+\tfrac13\cdot 0=0.6667$,
$H(\text{unemployed})=0$ $\Rightarrow$ the third branching on \textbf{unemployed}.

\textbf{The resulting tree, transformed into rules:}
\begin{align*}
&\texttt{IF } \text{income(high)} \texttt{ THEN } \text{loan approved(yes)}\\
&\texttt{IF } \text{income(low)}\wedge\text{account(high)} \texttt{ THEN }
   \text{loan approved(yes)}\\
&\texttt{IF } \text{income(low)}\wedge\text{account(low)} \texttt{ THEN }
   \text{loan approved(no)}\\
&\texttt{IF } \text{income(low)}\wedge\text{account(middle)}\wedge\text{unempl.(yes)}
   \texttt{ THEN } \text{loan approved(no)}\\
&\texttt{IF } \text{income(low)}\wedge\text{account(middle)}\wedge\text{unempl.(no)}
   \texttt{ THEN } \text{loan approved(yes)}
\end{align*}
\end{defbox}

\textbf{Transformation of a tree into rules} in general: each path between the root and a leaf
corresponds to one rule; the non-leaf nodes (together with their value for the respective edge)
appear on the left-hand side (in the body) of the rule, and the leaf node (the target class)
represents the right-hand side (the head).

\textbf{Factors important for the induction of a decision tree:}
\begin{itemize}
\item \emph{the choice of the decision attributes} --- it has an impact on the efficiency of the
      built tree; an \uv{informed} expert entry for the area considered is useful;
\item \emph{the order of the decision attributes} --- it eliminates redundant comparisons;
\item \emph{branching} --- the number of necessary divisions; more demanding for continuous values
      or many values;
\item \emph{the form of the built tree} --- balanced trees with as few levels as possible are more
      suitable (at the cost of more complex comparisons); some algorithms create binary trees
      only (often quite deep);
\item \emph{stopping criteria} --- correct classification of the training data; \emph{early
      stopping} prevents building large trees and overtraining (accuracy vs.\ efficiency);
      \emph{Occam's razor} (William of Occam, 1320) --- preference for the simplest hypothesis
      for $D$;
\item \emph{the training data and its quantity} --- too little $\Rightarrow$ the tree might not be
      reliable enough for more general data; too much $\Rightarrow$ the danger of overtraining;
\item \emph{pruning} --- a higher accuracy and efficiency for the tested data; removal of
      redundant comparisons, of entire subtrees and of irrelevant attributes; it can be done
      already during tree induction (which avoids forming excessively large trees), or on the
      already formed tree;
\item \emph{time and space complexity} --- depends on the number of training data $p$, the number
      of attributes $n$ and the form of the tree. Tree induction: $O(n\,p\log p)$; classification
      of $q$ patterns with a tree of depth $O(\log p)$: $O(q\log p)$.
\end{itemize}

\textbf{The ID3 algorithm} (Iterative Dichotomiser, Ross Quinlan, 1986): training of functions
with Boolean values, a \emph{greedy} method, builds the tree top-down; in each node it chooses the
attribute that best classifies the local training data --- the best attribute has the
\emph{biggest information gain}. The process continues until all the training data is classified
correctly, or until all the attributes have been used.

\begin{defbox}[The ID3 algorithm and information gain]
\textsc{ID3}(\emph{EXAMPLES}, \emph{TARGET\_ATTRIBUTE}, \emph{ATTRIBUTES}):
\begin{algorithmic}[1]
\State build a root $\mathit{ROOT}$ of the tree
\If{all \emph{EXAMPLES} are positive} \Return $\mathit{ROOT}$ with a single node labelled
  \uv{$+$} \EndIf
\If{all \emph{EXAMPLES} are negative} \Return $\mathit{ROOT}$ with a single node labelled
  \uv{$-$} \EndIf
\If{$\mathit{ATTRIBUTES}=\emptyset$} \Return $\mathit{ROOT}$ with a single node labelled by the
  most frequent value of \emph{TARGET\_ATTRIBUTE} in \emph{EXAMPLES} \EndIf
\State $A\gets$ the attribute from \emph{ATTRIBUTES} that best classifies \emph{EXAMPLES}
\State decision attribute for $\mathit{ROOT}\gets A$
\For{all possible values $v_i$ of the attribute $A$}
   \State attach a new edge to $\mathit{ROOT}$ corresponding to the predicate $A=v_i$
   \State $\mathit{EXAMPLES}_{v_i}\gets$ the subset of \emph{EXAMPLES} with the value $v_i$ for
          $A$
   \If{$\mathit{EXAMPLES}_{v_i}=\emptyset$}
      \State attach a leaf to the edge, labelled by the most frequent value of
             \emph{TARGET\_ATTRIBUTE}
   \Else
      \State attach the subtree \textsc{ID3}$(\mathit{EXAMPLES}_{v_i},
             \mathit{TARGET\_ATTRIBUTE}, \mathit{ATTRIBUTES}\setminus\{A\})$
   \EndIf
\EndFor
\State \Return $\mathit{ROOT}$
\end{algorithmic}

\textbf{Information gain} --- the expected entropy reduction after dividing the data according to
the attribute considered:
\[ \mathrm{GAIN}(D,A)=H(D)-\sum_{v\in \mathrm{VALUES}(A)}\frac{|D_v|}{|D|}H(D_v). \]
Since $H(D)$ is the same for all attributes, maximising the information gain is equivalent to
minimising the weighted entropy $H(A)$.
\end{defbox}

\textbf{The C4.5 algorithm} (Quinlan, 1993) --- a modification of ID3:
\begin{itemize}
\item \textbf{Missing data} --- ignored during tree construction (when evaluating the criterion
      for data division, only the known values of the given attribute are considered); during
      classification the missing values are estimated based on the values known for other
      patterns.
\item \textbf{Continuous data} --- categorised according to the values of the respective attribute
      on the patterns from the training set.
\item \textbf{Pruning by validation} --- the training data is divided into a training
      ($\tfrac23$) and a validation ($\tfrac13$) set. \emph{Reduced-error pruning}: subtrees are
      replaced by leaves labelled by the most frequent class of the respective training patterns;
      the new tree \emph{must not} perform worse on the validation set than the original one,
      otherwise the subtree is not replaced. \emph{Subtree raising}: a subtree is replaced by its
      most frequently used subtree, which is then raised to a higher level; the classification
      accuracy is tested.
\item \textbf{Rule post-pruning} --- a decision tree is inferred from the training set
      (over-training is allowed), converted into an equivalent set of rules (one rule per path
      from the root to a leaf), the rules are pruned (generalised) by eliminating all possible
      assumptions if this does not worsen the estimated classification accuracy, and are ordered
      by their expected accuracy; they are then used in this order to classify the patterns.
      Pruning is easier and more radical than for complete decision trees, and the resulting
      rules are transparent and easy to understand.
\item \textbf{Estimating the accuracy of rules} --- a standard approach uses a validation set;
      C4.5 evaluates the accuracy \emph{on the training data} as the ratio of patterns correctly
      classified by the rule to all patterns covered by the rule, $\mathit{acc}=k/n$. The standard
      deviation of the accuracy is assessed assuming a binomial distribution; the
      \emph{lower accuracy bound} for the chosen confidence interval is taken as the rule
      characteristic, at the 95~\% confidence level
      \[ \mathit{acc}-1.96\sqrt{\frac{\mathit{acc}(1-\mathit{acc})}{n}}. \]
\item \textbf{Information gain ratio} as the criterion for data division:
      \[ \mathrm{GAIN\_RATIO}(D,A)=\frac{\mathrm{GAIN}(D,A)}
         {-\sum_{v}\frac{|D_v|}{|D|}\log_2\frac{|D_v|}{|D|}}. \]
      The attribute with the \emph{biggest} information gain ratio is used for the division; the
      denominator penalises \emph{excessive branching} (attributes with many values).
\end{itemize}

\textbf{The C5.0 algorithm} --- a commercial successor of C4.5 for large databases: a more
efficient generation of rules and a higher accuracy thanks to \emph{boosting} (see the
ensemble methods below).

\begin{defbox}[CART --- Classification And Regression Trees]
Generates \emph{binary} trees --- during each data division just two edges are formed. The choice
of the best decision attribute uses entropy or the \textbf{Gini index}: for an attribute with $r$
values $v_1,\dots,v_r$, $n$ data patterns and $m$ classes
\[ G=\sum_{i=1}^{r}\frac{n_i}{n}G(v_i),\qquad G(v_i)=1-\sum_{j=1}^{m}p_j^2,
   \qquad \sum_{i=1}^{r}n_i=n, \]
where $n_i$ is the number of patterns taking on the value $v_i$ and $p_j$ the probability of class
$j$ in these $n_i$ patterns. \emph{Lower} values of $G$ imply better discrimination.

The division is performed according to the \uv{best} data point found for the respective node
$t$; all possible values $v$ of the decision attribute $A$ are checked (in the case of equality
the right-hand-side subtree is used). Let $P_L,P_R$ be the probabilities of processing by the
left-hand-side, resp.\ right-hand-side subtree, and $P(C_j\mid t_L),P(C_j\mid t_R)$ the
probabilities that a pattern belongs to class $C_j$ and is in the left-hand-side, resp.\
right-hand-side subtree of node $t$. The criterion for data division:
\[ \Phi=2P_LP_R\sum_{j=1}^{m}\big|P(C_j\mid t_L)-P(C_j\mid t_R)\big|, \]
which attains its maximum when all patterns from the class considered are grouped in the same
subtree.

\textbf{Characteristic properties of CART:} the order of the attributes corresponds to their
impact during classification; missing data is ignored and is not counted towards the evaluation of
the decision criteria; the stopping criterion --- no further division would improve the
classification accuracy; a high accuracy achieved on the training set does not have to reflect the
accuracy on the test data.
\end{defbox}

\subsubsection{Bayesian classification}

\begin{defbox}[Bayes' formula and the MAP hypothesis]
The probability that the hypothesis $H$ is true if we observe $E$:
\[ P(H\mid E)=\frac{P(E\mid H)\,P(H)}{P(E)}, \]
where $P(H)$ is the \emph{a-priori} probability of the hypothesis, $P(H\mid E)$ the
\emph{a-posteriori} (conditional) probability and $P(E)$ the probability of the event $E$.

The \emph{most probable hypothesis} $H_{MAP}$ (maximum a-posteriori) has the biggest a-posteriori
probability. Since the denominator $P(E)$ is the same for all hypotheses, it can be neglected:
\[ H_{MAP}=\arg\max_{H\in\mathcal H} P(H\mid E)=\arg\max_{H\in\mathcal H} P(E\mid H)\,P(H). \]
\end{defbox}

\begin{defbox}[Naive Bayesian classifier]
\textbf{Assumption:} the events $E_1,\dots,E_k$ are \emph{conditionally independent} subject to
the validity of the hypothesis $H$. In real-world tasks this assumption is rarely fulfilled ---
hence the name \uv{naive}. Then
\[ P(H\mid E_1,\dots,E_k)=\frac{P(E_1,\dots,E_k\mid H)P(H)}{P(E_1,\dots,E_k)}
   =\frac{P(H)}{P(E_1,\dots,E_k)}\prod_{i=1}^{k}P(E_i\mid H), \]
and we classify by the hypothesis with the biggest a-posteriori probability
\[ H_{MAP}=\arg\max_{H_t} P(H_t)\prod_{i=1}^{k}P(E_i\mid H_t). \]
The estimates from the data: $P(H_t)=P(\text{class }t)=\frac{\text{number of patterns from
}t}{\text{number of all patterns}}$ and $P(E_k\mid H_t)=\frac{\text{number of patterns from }t
\text{ matching }E_k}{\text{number of patterns from }t}$.

\textbf{Advantage:} the possibility to classify even \emph{incompletely described} patterns.
\textbf{Disadvantage:} a zero a-posteriori probability $P(E_k\mid H_t)$ if there is no
corresponding pattern in the training set, resp.\ underestimated terms in the case of a low
mutual incidence of $E_k$ and $H_t$.
\end{defbox}

\begin{defbox}[Example --- Bayesian classification of a loan]
\emph{The simple case:} $P(\text{borrow})=0.667$, $P(\text{reject})=0.333$,
$P(\text{high income}\mid\text{borrow})=0.91$,
$P(\text{high income}\mid\text{reject})=0.12$. Then
$0.91\cdot 0.667=0.607$ vs.\ $0.12\cdot 0.333=0.040$, hence
$H_{MAP}=\text{borrow}$.

\emph{The naive classifier} over the loan data (12 clients):
$P(\text{credit(yes)})=8/12=0.667$, $P(\text{credit(no)})=4/12=0.333$;
$P(\text{account(middle)}\mid\text{yes})=2/8=0.25$,
$P(\text{account(middle)}\mid\text{no})=2/4=0.5$;
$P(\text{unempl.(no)}\mid\text{yes})=5/8=0.625$,
$P(\text{unempl.(no)}\mid\text{no})=1/4=0.25$.

For a client with a middle balance on their account who is not unemployed:
\[ 0.667\cdot 0.25\cdot 0.625=0.1042 \quad\text{vs.}\quad
   0.333\cdot 0.5\cdot 0.25=0.0416, \]
so the client is assigned to the class \textbf{credit(yes)}.
\end{defbox}

\subsubsection{Support Vector Machines}

SVMs were invented by V.~Vapnik and his co-workers in the 1970s in Russia and became known to the
West in 1992. They are \emph{linear} classifiers that find a hyperplane separating two classes
(positive and negative); for non-linear separation, \emph{kernel} functions are used. SVM not only
has a rigorous theoretical foundation but also performs classification more accurately than most
other methods in applications, especially for high-dimensional data --- it belongs to the best
classifiers for text classification.

\begin{defbox}[Linear SVM, the separable case]
The training set $D=\{(x_1,y_1),\dots,(x_r,y_r)\}$, $x_i\in X\subseteq\mathbb R^n$,
$y_i\in\{1,-1\}$. SVM finds a linear function $f(x)=\langle w\cdot x\rangle+b$ such that
\[ y_i=\begin{cases}1 & \text{if }\langle w\cdot x_i\rangle+b\ge 0,\\
-1 & \text{if }\langle w\cdot x_i\rangle+b<0.\end{cases} \]
The hyperplane $\langle w\cdot x\rangle+b=0$ is called the \emph{decision boundary} (surface).
There are many such hyperplanes --- SVM looks for the one with the \emph{largest margin}; machine
learning theory says that this hyperplane minimises the error bound.

The (perpendicular) distance from a point $x_i$ to the hyperplane is
$\dfrac{|\langle w\cdot x_i\rangle+b|}{\|w\|}$, where
$\|w\|=\sqrt{\langle w\cdot w\rangle}$. For the margin hyperplanes $H_+$ and $H_-$,
\[ d_+=d_-=\frac{1}{\|w\|},\qquad \text{margin}=d_++d_-=\frac{2}{\|w\|}. \]

\textbf{The optimisation problem:}
\[ \text{minimise }\ \frac{\langle w\cdot w\rangle}{2}
   \quad\text{subject to}\quad y_i\big(\langle w\cdot x_i\rangle+b\big)\ge 1,\
   i=1,\dots,r, \]
which summarises the constraints $\langle w\cdot x_i\rangle+b\ge 1$ for $y_i=1$ and
$\langle w\cdot x_i\rangle+b\le -1$ for $y_i=-1$.
\end{defbox}

\begin{defbox}[The dual formulation and support vectors]
The problem is solved by the standard \emph{Lagrangian} method --- the constraints are pulled
into the objective function using \emph{Lagrange multipliers} $\alpha_i$, one per training
constraint:
\[ L_P=\frac{\langle w\cdot w\rangle}{2}-\sum_{i=1}^{r}\alpha_i
   \big[y_i(\langle w\cdot x_i\rangle+b)-1\big],\qquad \alpha_i\ge 0. \]

\textbf{Kuhn--Tucker conditions.} An optimal solution must satisfy them. In general they are only
\emph{necessary}, but for a \emph{convex} objective function with \emph{linear} constraints ---
which is exactly our case --- they are also \emph{sufficient}:
\begin{align*}
\frac{\partial L_P}{\partial w_j} &= w_j - \sum_{i=1}^{r} y_i\alpha_i x_{ij} = 0,
  \quad j=1,\dots,n && \text{stationarity in }w\\
\frac{\partial L_P}{\partial b} &= -\sum_{i=1}^{r} y_i\alpha_i = 0
  && \text{stationarity in }b\\
y_i\big(\langle w\cdot x_i\rangle+b\big)-1 &\ge 0, \quad i=1,\dots,r
  && \text{the original constraints}\\
\alpha_i &\ge 0, \quad i=1,\dots,r && \text{non-negativity of multipliers}\\
\alpha_i\big(y_i(\langle w\cdot x_i\rangle+b)-1\big) &= 0, \quad i=1,\dots,r
  && \text{\textbf{complementarity}}
\end{align*}
where $x_{ij}$ is the $j$-th component of the vector $x_i$ (in vector form
$w=\sum_i y_i\alpha_i x_i$).

\textbf{Why complementarity yields the support vectors:} the product is zero, so for every $i$ at
least one of the two factors must vanish. If $\alpha_i>0$, the second bracket must be zero, i.e.\
$y_i(\langle w\cdot x_i\rangle+b)=1$ --- the point lies \emph{exactly on} the margin hyperplane
$H_+$ or $H_-$. These are the \textbf{support vectors}. For all the other points $\alpha_i=0$, so
they do not enter the solution at all: an SVM is determined by a~small subset of the training data
only.

\textbf{From the primal to the dual program} --- two steps:
\begin{enumerate}
\item \emph{Set to zero the partial derivatives of $L_P$ with respect to the primal variables}
      $w$ and $b$. This gives the first two KKT conditions, rewritten as
      \[ w=\sum_{i=1}^{r} y_i\alpha_i x_i, \qquad \sum_{i=1}^{r} y_i\alpha_i = 0. \]
\item \emph{Substitute these relations back into $L_P$.} The primal variables $w$ and $b$ are
      thereby \textbf{eliminated} and only the multipliers $\alpha_i$ remain.
\end{enumerate}
The result is the \textbf{Wolfe dual} problem. \emph{Note:} it is a~\textbf{constrained} problem
--- the relations obtained in step~1 are part of it:
\[ \text{maximize }\ L_D=\sum_{i=1}^{r}\alpha_i-\frac12\sum_{i,j}y_iy_j\alpha_i\alpha_j
   \langle x_i\cdot x_j\rangle
   \quad\text{subject to}\quad \sum_{i=1}^{r} y_i\alpha_i=0,\ \ \alpha_i\ge 0. \]
The maximum of $L_D$ occurs at the same values of $w,b,\alpha_i$ as the minimum of $L_P$. The dual
is easier to solve than the primal (numerical techniques, in practice e.g.\ SMO) and, crucially,
it involves the input vectors \emph{only} through the dot products $\langle x_i\cdot x_j\rangle$
--- precisely what makes the kernel trick possible (see below).

\textbf{Recovering $w$ and $b$.} Solving the dual yields the $\alpha_i$; from them
\[ w=\sum_{i\in \mathit{sv}} y_i\alpha_i x_i, \qquad
   b = y_k - \langle w\cdot x_k\rangle \quad\text{for any support vector } x_k, \]
where the expression for $b$ follows from complementarity: for a~support vector
$y_k(\langle w\cdot x_k\rangle+b)=1$, and since $y_k\in\{1,-1\}$, hence $y_k^2=1$, it suffices to
multiply by $y_k$. Numerically, $b$ is usually \emph{averaged} over all support vectors.

\textbf{The decision boundary} and \textbf{testing} an instance $z$:
\[ \langle w\cdot x\rangle+b=\sum_{i\in \mathit{sv}}y_i\alpha_i\langle x_i\cdot x\rangle+b=0,
\qquad
\mathrm{sign}\Big(\sum_{i\in \mathit{sv}}\alpha_iy_i\langle x_i\cdot z\rangle+b\Big). \]
If the expression returns 1, $z$ is classified as positive; otherwise as negative.
\end{defbox}

\begin{defbox}[The kernel trick]
Real-life data may require non-linear separation. The data is therefore \emph{transformed} into
another (usually much higher-dimensional) space in which a \emph{linear} boundary separates the
classes: $\phi:X\to F$, $x\mapsto\phi(x)$; the transformed space $F$ is called the \emph{feature
space}, the original one the \emph{input space}.

\textbf{The problem with explicit transformation:} the number of dimensions of the feature space
can be huge for useful transformations even with a reasonable number of input attributes --- the
\emph{curse of dimensionality} makes explicit computation infeasible.

\textbf{The crucial observation:} in the dual formulation, both the construction of the optimal
hyperplane and the evaluation of the decision function require \emph{only the dot products}
$\langle\phi(x)\cdot\phi(z)\rangle$, never the mapped vector $\phi(x)$ itself. So if we have a way
to compute the dot product directly from the input vectors, we do not need to know $\phi$ at all
--- a \emph{kernel function}
\[ K(x,z)=\langle\phi(x)\cdot\phi(z)\rangle. \]
In all expressions the dot product is replaced by the kernel --- the \textbf{kernel trick}.

\emph{Example (polynomial kernel with $d=2$, 2-D input):}
\begin{align*}
\langle x\cdot z\rangle^2&=(x_1z_1+x_2z_2)^2=x_1^2z_1^2+2x_1z_1x_2z_2+x_2^2z_2^2\\
&=\big\langle (x_1^2,x_2^2,\sqrt2 x_1x_2)\cdot(z_1^2,z_2^2,\sqrt2 z_1z_2)\big\rangle
 =\langle\phi(x)\cdot\phi(z)\rangle.
\end{align*}
The kernel $\langle x\cdot z\rangle^d$ is therefore a dot product in a transformed feature space
without ever evaluating $\phi$. Whether a given function is a kernel (i.e.\ a dot product in some
feature space) is answered by \emph{Mercer's theorem}. Commonly used kernels: linear (the
identity), polynomial, Gaussian (RBF), sigmoid.
\end{defbox}

\textbf{Further properties of SVM.} It works only in a real-valued space --- categorical
attributes have to be converted to numeric values. It does only \emph{two-class} classification;
for multi-class problems, strategies such as \emph{one-against-rest} or \emph{error-correcting
output coding} are applied. The hyperplane produced by SVM is hard to understand by human users
and the kernels make the matter worse --- SVM is therefore commonly used in applications that do
not require human understanding. \mimo{} For non-linearly separable data, \emph{slack variables}
$\xi_i\ge 0$ are introduced and $\frac{\langle w\cdot w\rangle}{2}+C\sum_i\xi_i$ is minimised;
the parameter $C$ controls the trade-off between the width of the margin and the number of errors.

\subsubsection{Logistic regression}
\label{ssec:logreg}

Logistic regression is the most important case of non-linear regression
(Chapter~\ref{sec:priprava}) and at the same time a classifier: the dependent variable $y$ is
categorical, typically two-valued, and what is modelled is the \emph{probability} that $y$ takes
a certain value, depending on the combination of the values of the independent quantities.

\begin{defbox}[Logistic regression]
For binary $y\in\{0,1\}$ one models the logarithm of the \uv{conditional odds}:
\[ \ln\frac{P(y=1\mid x_1,\dots,x_m)}{1-P(y=1\mid x_1,\dots,x_m)}
   =\beta_0+\beta_1x_1+\cdots+\beta_mx_m=z, \]
whence
\[ P(y=1\mid x)=\frac{e^{z}}{1+e^{z}}=\frac{1}{1+e^{-z}}=\sigma(z)
   \qquad\text{(the sigmoid function).} \]

\textbf{Learning} by maximum likelihood:
\[ L=\prod_{i=1}^{p}P\big(y_i\mid x_{i,1},\dots,x_{i,m}\big)
   =\prod_{i=1}^{p}\Big(\frac{1}{1+e^{-z_i}}\Big)^{y_i}
     \Big(\frac{1}{1+e^{z_i}}\Big)^{1-y_i}. \]
Instead of maximising $L$, one minimises the \emph{cross-entropy} loss $E=-\log L$:
\[ E=-\sum_i\Big[y_i\log\sigma(z_i)+(1-y_i)\log\big(1-\sigma(z_i)\big)\Big]. \]
Since the derivative of the sigmoid is $\sigma'=\sigma(1-\sigma)$, after simplification we obtain
very simple adjustment rules
\[ \Delta\beta_j=-\eta\frac{\partial E}{\partial\beta_j}
   =\eta\sum_i\big(y_i-\sigma(z_i)\big)x_{i,j},\qquad
   \Delta\beta_0=\eta\sum_i\big(y_i-\sigma(z_i)\big), \]
where $\eta$ denotes the learning rate. The error thus enters the weight update directly as the
difference between the \emph{observed} and the \emph{predicted} value.
\end{defbox}

Logistic regression is a common choice for the base classifier wherever a probabilistic output
and interpretable coefficients are required --- e.g.\ in link prediction in social networks
(Chapter~\ref{sec:sna}).

\subsubsection{Discriminant analysis}
\label{ssec:da}

\begin{defbox}[Discriminant analysis]
\emph{Classification of patterns into known classes}: we search for the dependence of the nominal
quantity (that impacts the class membership) on the other numerical quantities. We assume that to
every class $C_j$, $j=1,\dots,m$ there exists a \emph{discriminative function} $g_j$; a pattern
$x=(x_1,\dots,x_m)$ belongs to class $C_j$ if and only if
\[ g_j(x)=\max_{k} g_k(x). \]

\textbf{Linear discriminant analysis:}
$g_j(x)=\beta_{j0}+\beta_{j1}x_1+\cdots+\beta_{jm}x_m$. For two classes, instead of $g_1,g_2$ we
can look for the single function $g=g_1-g_2$ and classify according to its \emph{sign}: \uv{$+$}
indicates class 1, \uv{$-$} class 2.

\textbf{An optimum classification in the sense of a minimum error} occurs when the discriminant
functions are $\propto$ the a-posteriori (conditional) probabilities of classifying the
observation $x$ into class $C_j$: $g_j(x)=P(C_j\mid x)=\dfrac{p(x\mid C_j)P(C_j)}{p(x)}$; for two
classes $g(x)=p(x\mid C_1)P(C_1)-p(x\mid C_2)P(C_2)$.
\end{defbox}

For a \emph{normal distribution} with different means $\mu_1,\mu_2$ and general covariance
matrices $S_1,S_2$ we get a \emph{quadratic} discriminant function. If the covariance matrices are
\emph{identical} ($S_1=S_2=S$), it reduces to a \emph{linear} one,
\[ g(x)=(\mu_1-\mu_2)^{\!\top}S^{-1}x-\tfrac12(\mu_1-\mu_2)^{\!\top}S^{-1}(\mu_1+\mu_2)
   -\ln\frac{P(C_2)}{P(C_1)}, \]
and if moreover both classes are equally probable and $S=I$, it is a simple comparison of the
distances from the means. For normal distributions the search for discriminant functions therefore
reduces to \textbf{estimating the means} (from the sample means) and the \textbf{covariance
matrices} (composed of corrected sample variances).

\subsubsection{Neural networks and extreme learning machines (ELM)}

\begin{defbox}[The ELM network (Extreme Learning Machine, G.-B.~Huang)]
A feedforward network with \emph{one} hidden layer ($n$ inputs, $L$ hidden neurons, 1 output):
$f_L(x)=\sum_{i=1}^{L}\beta_i\,G(a_i,b_i,x)$, where $G$ is e.g.\ a sigmoid
$g(\langle a_i\cdot x\rangle+b_i)$ or an RBF $g(b_i\|x-a_i\|)$.

\textbf{The key idea:} the parameters of the hidden neurons $(a_i,b_i)$ are \emph{not
adapted} --- they are generated \emph{randomly}, even without the knowledge of the training
data (with probability one, $\lim_{L\to\infty}\|f-f_L\|=0$). Only the output weights $\beta$
are learned, and \emph{once} --- by solving the system $H\beta=T$, where $H$ is the
\emph{output matrix of the hidden layer} (its $j$-th row $=$ the outputs of the hidden
neurons for the pattern $x_j$):
\begin{algorithmic}[1]
\State choose \emph{randomly} the parameters of the hidden neurons $(a_i,b_i)$, $i=1,\dots,L$
\State compute the output matrix of the hidden layer $H$
\State $\beta\gets H^{+}T$ \Comment{the Moore--Penrose pseudoinverse}
\end{algorithmic}
Regularisation (adding $I/C$ to the diagonal of $H^{\!\top}H$) gives a more stable solution
with better generalisation.
\end{defbox}

\textbf{Advantages of ELM:} three non-iterative steps --- very fast, no local minima or
initialisation issues; the hidden-neuron parameters do not depend on the data; any bounded
piecewise continuous transfer function suffices (gradient algorithms require a smooth one).

\subsubsection{Ensemble methods}

\begin{defbox}[Bagging (bootstrapped aggregating)]
Attempts to reduce the \emph{variance} of the classifier: if the variance of a classifier is
$\sigma^2$, then the variance of the average of $k$ independent and identically distributed
(i.i.d.) classifiers reduces to $\sigma^2/k$.

\textbf{Bootstrapping} approximates i.i.d.\ classifiers: the data points are sampled uniformly
from the original data \emph{with replacement}; the sample is drawn of approximately the same size
as the original data, may comprise duplicates, and contains a fraction of roughly
\[ 1-\Big(1-\frac1n\Big)^{\!n}\approx 1-\frac1e=63.2\,\% \]
distinct data points from the original data (the probability of a data point not being included is
$(1-1/n)^n$). A total of $k$ different bootstrapped samples of size $n$ are drawn independently, a
classifier is trained on each of them, and for a test instance the predicted class label is
reported as the \emph{dominant vote} of the different classifiers.

\textbf{Suitability:} decision trees are an ideal choice for bagging --- they have a \emph{low
bias} and \emph{high variance} (if grown sufficiently deep).
\textbf{Limitation:} bagging does \emph{not reduce the model bias}; for cases in which the bias is
the primary source of error it may even slightly worsen accuracy due to correlations between the
ensemble components (and hence the failed i.i.d.\ assumption).
\end{defbox}

\begin{defbox}[Random forests]
If $k$ classifiers, each with variance $\sigma^2$, have a positive pairwise \emph{correlation}
$\rho$ between them, then the variance of the averaged prediction is
\[ \rho\sigma^2+\frac{(1-\rho)\sigma^2}{k}. \]
The term $\rho\sigma^2$ is \emph{invariant} to the number of components $k$ --- it is precisely
this term that limits the performance gains from bagging. Bootstrapped decision trees, moreover,
tend to be positively correlated.

A \emph{random forest} is an ensemble of decision trees into whose model-building process
randomness has explicitly been inserted, so that the components are less correlated. Two options:
\emph{(1)}~\textbf{bagging} --- make the trees different by training them on slightly different
data (bootstrap samples); \emph{(2)}~a \textbf{limited number of features} --- at each node, a
\emph{random subset} of $m$ features is given to the tree and the information gain is computed
only on that subset. The second option also speeds up training (fewer features to search over) and
the trees need not be pruned; the reduced variance does not negatively affect the bias. The output
of the forest is the majority vote of the trees.

\textbf{The algorithm:} for each of the $N$ trees: create a new bootstrap sample of the training
set; use it to train a decision tree; at each node randomly select $m$ features, compute the
information gain only on that set and select the optimal one; repeat until the tree is complete.

\textbf{Parameters:} the number of features $m$ --- a subset size that is the square root of the
number of features seems to be common (random forests are not too sensitive to this parameter);
the number of trees --- building trees can be continued until the error stops decreasing.
\end{defbox}

\begin{defbox}[Boosting and AdaBoost (Freund \& Schapire, 1996)]
\textbf{The principle:} a \emph{weight} is associated with each training pattern and different
classifiers are trained using these weights; the weights are iteratively modified based on the
performance of the classifier. The models to be constructed therefore \emph{depend on the previous
ones} --- each classifier is constructed using the same algorithm on a different weighted training
set. \textbf{The idea:} in future iterations, focus on the \emph{misclassified} patterns by
increasing their relative weight; by iteratively using this approach and creating a weighted
combination of the various classifiers, it is possible to create a classifier with a lower overall
\emph{bias} (unlike bagging, which reduces the variance).

\textsc{AdaBoost}(data set $D$, base classifier $A$, maximum rounds $T$):
\begin{algorithmic}[1]
\State $t\gets 0$; for each instance $i$ initialise $W_1(i)=1/n$
\Repeat
  \State $t\gets t+1$
  \State determine the weighted error rate $\epsilon_t$ on $D$ when the base algorithm $A$ is
         applied to the data weighted with $W_t(\cdot)$
  \State $\alpha_t\gets \tfrac12\ln\big((1-\epsilon_t)/\epsilon_t\big)$
  \For{each \emph{misclassified} $X_i\in D$}
     \State $W_{t+1}(i)\gets W_t(i)e^{\alpha_t}$
  \EndFor
  \For{each \emph{correctly classified} $X_i\in D$}
     \State $W_{t+1}(i)\gets W_t(i)e^{-\alpha_t}$
  \EndFor
  \State normalise $W_{t+1}(i)\gets W_{t+1}(i)/\sum_j W_{t+1}(j)$
\Until{$t\ge T$ \textbf{or} $\epsilon_t=0$ \textbf{or} $\epsilon_t\ge 0.5$}
\State classify test instances with the ensemble using the component weights $\alpha_t$
\end{algorithmic}
$\epsilon_t$ is the fraction of incorrectly predicted training instances (computed after weighting
with $W_t$) by the model in the $t$-th iteration.

\textbf{Properties:} it produces a \emph{sequence} of classifiers with the same base learner, each
depending on the previous one and focusing on its errors. At testing, the results of the sequence
are combined. \textbf{Sensitivity to noise:} when the number of outliers is very large, the
emphasis placed on the hard examples can hurt performance.
\end{defbox}

\begin{center}
\small
\begin{tabular}{@{}p{7.6cm}p{8.0cm}@{}}
\toprule
\textbf{Random forests} & \textbf{Boosting} \\
\midrule
fast --- they can run \emph{in parallel} and search over a small set of features at each stage &
  exhaustive --- searches over the \emph{entire} feature set at each stage, and each stage depends
  on the previous one \\
as the trees are cheap to train, more of them can be grown in the same computational time, even on
  large and complicated data sets & must run \emph{sequentially}, which can be expensive \\
can produce surprisingly good results in comparison to the small parts of the problem space seen
  by each tree & for the \emph{same number} of trees, boosting outperforms random forests (as the
  forests search only a small subset at each stage) \\
\bottomrule
\end{tabular}
\end{center}

\subsubsection{Comparison of the classification methods}

\begin{center}
\small
\begin{tabular}{@{}p{2.9cm}p{2.3cm}p{2.2cm}p{2.3cm}p{4.6cm}@{}}
\toprule
\textbf{Method} & \textbf{Comprehen\-sibility} & \textbf{Accuracy} & \textbf{Speed}
  & \textbf{Note} \\
\midrule
Decision tree & high & medium & high & rectangular regions, pruning needed, ignores attribute
  correlations \\
Rules from C4.5 & highest & medium & high & rule pruning is more radical than tree pruning \\
Naive Bayes & medium & surprisingly good & very high & handles incomplete patterns; problem with
  zero probabilities \\
SVM & very low & high & low (training) & best for high-dimensional data (text); only 2 classes,
  only a real-valued space \\
Logistic regression & high & medium & high & probabilistic output, interpretable coefficients \\
Discriminant analysis & medium & good under normality & high & only $\mu_j$ and $S_j$ need to be
  estimated \\
Neural network / ELM & very low & high & ELM very high & ELM has no local minima and does not
  train the hidden layer \\
Random forest & low & very high & high (parallel) & reduces variance, does not worsen bias \\
Boosting / AdaBoost & low & very high & low (sequential) & reduces bias; sensitive to noise and
  outliers \\
\bottomrule
\end{tabular}
\end{center}

\subsection{Cluster analysis}
\label{ssec:clustering}

\subsubsection{Formulation of the task and distance measures}

\emph{Cluster analysis} divides the observed patterns into groups (clusters) of mutually similar
patterns. \textbf{Assumption:} we can measure the \emph{distance} between patterns. Each pattern is
characterised by $m$ numerical quantities.

\begin{defbox}[Distance measures]
For patterns $x_1=(x_{11},\dots,x_{1m})$ and $x_2=(x_{21},\dots,x_{2m})$:
\begin{enumerate}
\item \textbf{Hamming} (Manhattan, $L_1$):
      $d_H(x_1,x_2)=\sum_{i=1}^{m}|x_{1i}-x_{2i}|$
\item \textbf{Euclidean} ($L_2$):
      $d_E(x_1,x_2)=\sqrt{\sum_{i=1}^{m}(x_{1i}-x_{2i})^2}$
\item \textbf{Chebyshev} ($L_\infty$):
      $d_C(x_1,x_2)=\max_i |x_{1i}-x_{2i}|$
\item \textbf{Minkowski metric} ($L_q$) --- the previous three are its special cases:
      $d_q(x_1,x_2)=\Big(\sum_{i=1}^{m}|x_{1i}-x_{2i}|^{q}\Big)^{1/q}$
\end{enumerate}
We have $d_H=d_1$, $d_E=d_2$, $d_C=\lim_{q\to\infty}d_q$; for the pattern $(0,\dots,0)$ and a
pattern at unit distance all three metrics give the value 1.
\end{defbox}

The choice of the metric depends on the quantities involved and their range $\Rightarrow$ the
quantities have to be \textbf{normalised} (divided by their mean, standard deviation, range
$\max-\min$, \dots; see Chapter~\ref{sec:priprava}). The metrics above assume the \emph{same
variance} for all the quantities. For quantities with \emph{different} variance one uses

\begin{defbox}[Mahalanobis distance]
\[ d_M^2(x_1,x_2)=(x_1-x_2)^{\!\top}S^{-1}(x_1-x_2), \]
where $S$ is the covariance matrix. It takes both the different variances and the correlations
between the quantities into account (for $S=I$ it reduces to the squared Euclidean distance).
\end{defbox}

\begin{defbox}[The distance between two clusters $C_k$ and $C_l$]
\textbf{The method of the nearest neighbour} (single linkage) --- the minimum of the distances
between their elements: $D(C_k,C_l)=\min_{x_i,x_j} d(x_i,x_j)$, $x_i\in C_k$, $x_j\in C_l$.

\textbf{The method of the farthest neighbour} (complete linkage) --- the maximum:
$D(C_k,C_l)=\max_{x_i,x_j} d(x_i,x_j)$.

\textbf{The method of average distance} (average linkage) --- the average of the distances between
the patterns ($n_k$, $n_l$ being the numbers of patterns in the clusters):
\[ D(C_k,C_l)=\frac{1}{n_kn_l}\sum_{i=1}^{n_k}\sum_{j=1}^{n_l} d(x_i,x_j). \]

\textbf{The centroid method} --- the distance between the centres of the clusters:
$D(C_k,C_l)=d(c_k,c_l)$, where $c_k$, $c_l$ are the centres (centroids) of the clusters.
\end{defbox}

A \emph{centroid} is the centre of a cluster --- a \uv{prototype} representing the given cluster.
The same cluster can be represented by one or even more centroids at the same time, depending on
the shape of the cluster and the metric chosen.

\subsubsection{The $k$-means method}

\begin{defbox}[The $k$-means algorithm (Lloyd, 1982)]
\begin{algorithmic}[1]
\State divide the patterns \emph{randomly} into $k$ clusters $C_1,\dots,C_k$ (disjoint, their
       union being all patterns)
\State \label{alg:km2} determine the centroids of all the clusters in the recent division, e.g.\
       $c_k=\frac{1}{|C_k|}\sum_{x\in C_k}x$
\For{each pattern $x$}
   \State determine the distances $d(x,c_k)$ for $1\le k\le K$
   \State let $d(x,c_l)=\min_k d(x,c_k)$
   \If{$x$ does not belong to cluster $l$} \State move $x$ to cluster $l$ \EndIf
\EndFor
\If{there was a transfer} \State goto step~\ref{alg:km2} \Else \State \textbf{finish} \EndIf
\end{algorithmic}
Afterwards, the clusters are represented by their centroids.

\textbf{Variants:} at the initial division, pronounce the \emph{first $k$ patterns} to be
centroids (which eliminates step~\ref{alg:km2} in the first iteration); adjust the centroids
\emph{after every} transfer (inside the cycle rather than after it).
\end{defbox}

\textbf{Extensions to the Lloyd algorithm.} Alternating minimisation steps do not provide an
approximation guarantee, yet Lloyd's algorithm \emph{never increases} the cost of the initial
clustering $\Rightarrow$ it pays to provide the algorithm with a \emph{provably good
initialisation}:
\begin{itemize}
\item \textbf{$k$-means++} (Arthur, Vassilvitskii, 2007) --- sample $k$ of the $P$ patterns with
      probability proportional to their squared distance to the closest existing centroid:
      $\dfrac{d^2(x_i,c)}{\sum_{j=1}^{P}d^2(x_j,c)}$.
\item \textbf{LocalSearch++} (Lattanzi, Sohler, 2019) --- first use $k$-means++ to pick $k$
      initial centroids; afterwards, new patterns are sampled with an analogous probability and
      replace existing centroids if they decrease the overall squared distance between the
      patterns and their closest centroids.
\end{itemize}

\textbf{The $k$-medians method} uses the Manhattan distance as the objective function,
$d(x_i,c)=\|x_i-c\|_1$. It can be shown that the optimal representative $c$ is then the
\emph{median} of the data points along each dimension in the cluster. As the median is chosen
independently along each dimension, the resulting $m$-dimensional representative will typically
\emph{not} belong to the original data set. $k$-medians selects the cluster representatives in a
\emph{more robust} way than $k$-means, because the median is not as sensitive to outliers as the
mean.

\subsubsection{The $k$-medoids method, CLARA and CLARANS}

\begin{defbox}[The $k$-medoids algorithm]
For a data set $D$ of $n$ $d$-dimensional points $x_1,\dots,x_n$, the goal is to determine $k$
representatives $y_1,\dots,y_k$ \emph{out of the original database} $D$ ($k$ is specified by the
user) that minimise the objective function given by the sum of the distances between the data
points and their closest representatives:
\[ O=\sum_{i=1}^{n}\min_{j} d(x_i,y_j). \]
The time complexity of each iteration is $O(k\,n\,d)$; usually a small constant number of
iterations is necessary.

\textbf{Why select the representatives from $D$?} \emph{(1)}~$k$-means clusters may be distorted by
outliers --- the cluster representatives can be located in empty regions, unrepresentative of most
of the data points in that cluster, which may lead to a partial merging of different clusters
(this can be resolved at least partially with careful outlier handling). \emph{(2)}~For complex
data (e.g.\ a set of time series of varying length) it might be difficult to compute the optimal
central representative at all. The $k$-medoids algorithm can therefore be defined virtually for
\emph{any data type}, as long as an appropriate similarity or distance function can be provided.

\textbf{The procedure:} \emph{hill-climbing} --- the representative set $S$ is initialised to a
set of points from $D$ and is iteratively improved by exchanging a point from $S$ with a data point
from $D$. Possible strategies for the exchange:
\begin{enumerate}
\item try all $|S|\cdot|D|$ possibilities and select the best one --- \emph{extremely expensive};
\item use a randomly selected set of $r$ pairs $(x,y)$ ($x$ from $D$, $y$ from $S$) and use the
      best of the $r$ pairs for the exchange. This requires time proportional to $r\cdot|D|$ and is
      said to have converged when the objective function does not improve, or if the average
      improvement is below a user-specified threshold.
\end{enumerate}
$k$-medoids is \emph{slower} than $k$-means but applicable to various data types.
\end{defbox}

\begin{defbox}[Scalable variants: CLARA and CLARANS]
Scalable implementations of $k$-medoids; the base variant \emph{PAM} (Partitioning Around
Medoids) tests all $k(n-k)$ medoid--non-medoid exchanges, $O(kn^2d)$ per iteration --- too
expensive.
\begin{itemize}
\item \textbf{CLARA} (Clustering LARge Applications): apply PAM to a \emph{sample} only and
      assign the remaining points to the medoids found; repeat over independent samples,
      select the best. Risk: a sample may not include good medoids.
\item \textbf{CLARANS} (\dots\ based on RANdomized Search): works with the \emph{full} data
      and tries \emph{random} exchanges; a local optimum after \emph{MaxAttempt} failures,
      the whole process \emph{MaxLocal} times. Explores a greater diversity of solutions than
      CLARA.
\end{itemize}
\end{defbox}

\subsubsection{Hierarchical clustering}

\begin{defbox}[Agglomerative hierarchical clustering]
A \uv{bottom-up} approach.

\textbf{Initialisation:} \emph{(1)}~determine the mutual distances between all patterns;
\emph{(2)}~assign each pattern to an individual cluster.

\textbf{Main cycle:} while there are more clusters than one: \emph{(1)}~find the two mutually
closest clusters and merge them together; \emph{(2)}~for the new cluster, compute its distance from
the other clusters.

The result is displayed as a \textbf{dendrogram} --- it shows (from the left to the right) the
gradual merging of clusters; the height of a join corresponds to the distance at which it occurred.
\emph{The optimum number of clusters is not known in advance} --- it is chosen by cutting the
dendrogram at a suitable height.
\end{defbox}

\textbf{Divisive} hierarchical clustering (\uv{top-down}) proceeds the other way round: it starts
with one cluster containing all patterns and divides it gradually. The best-known example in the
context of networks is the Girvan--Newman algorithm (Chapter~\ref{sec:sna}).

There is also a scalable agglomerative variant, \textbf{CURE} (\emph{Clustering Using
REpresentatives}): sampling and hierarchical merging per partition; it is very fast and
discovers clusters of an \emph{arbitrary shape}.

\subsubsection{Vector quantisation: the LVQ algorithm}

\begin{defbox}[LVQ --- Learning Vector Quantisation]
Used for \textbf{supervised clustering} (the classes are denoted by $C$):
\begin{algorithmic}[1]
\State initialise all weight vectors $w_j(0)$, the learning rates $\mu(0)$, $k\gets 0$
\While{the stop condition is not satisfied}
  \For{each training pattern $x$}
     \State determine the index $j=q$ of the weight vector with the \emph{minimum} Euclidean
            distance, i.e.\ $\min_j\|x-w_j\|_2^2$
     \If{the class of $x$ agrees with the class of $w_q$}
        \State $w_q(k+1)\gets w_q(k)+\mu(k)\big(x-w_q(k)\big)$ \Comment{pull towards}
     \Else
        \State $w_q(k+1)\gets w_q(k)-\mu(k)\big(x-w_q(k)\big)$ \Comment{push away}
     \EndIf
  \EndFor
  \State $k\gets k+1$; decrease the learning rate, e.g.\ $\mu(k)=\mu(k-1)/(k+1)$ for $k>0$
\EndWhile
\end{algorithmic}
The weight vectors are usually initialised by the \emph{first $m$ vectors} of the training set
(which must include patterns from all the classes) --- this automatically \uv{calibrates} the
neurons with the classes.
\end{defbox}

\subsubsection{Grid-based and density-based methods}

Grid-based and density-based methods find clusters of an \emph{arbitrary shape} as
\textbf{connected dense regions} (the connected components of a connectivity graph).
\emph{Grid-based} methods (e.g.\ MAFIA) discretise the space into $p^d$ hypercubes and connect
\emph{dense} cells (at least $\tau$ points) that share a side. The density-based
\textbf{DBSCAN} builds on \emph{core points} instead of cells (at least $\tau$ points within
the radius $\varepsilon$): the clusters are the connected components of the graph of core
points at most $\varepsilon$ apart, \emph{border} points are assigned to the closest cluster,
and \emph{noise} points are reported as outliers.

\begin{center}
\small
\begin{tabular}{@{}p{2.6cm}p{2.3cm}p{2.4cm}p{2.2cm}p{5.0cm}@{}}
\toprule
\textbf{Method} & \textbf{Cluster shape} & \textbf{Number of clusters} & \textbf{Outliers}
  & \textbf{Note} \\
\midrule
$k$-means & spherical & given & sensitive & fast; depends on initialisation $\to$ $k$-means++ \\
$k$-medians & spherical & given & more robust & median per dimension; the representative is not
  from the data \\
$k$-medoids & spherical & given & robust & representatives \emph{from the data}; any data type;
  slower \\
CLARA / CLARANS & spherical & given & robust & scalable $k$-medoids (sampling / randomised
  search) \\
Hierarchical & depends on the linkage & read off the dendrogram & sensitive & no need to give $k$
  in advance; $O(n^2)$--$O(n^3)$ \\
CURE & arbitrary & given & robust & sampling $+$ hierarchical merging per partition \\
LVQ & spherical & given by the number of neurons & --- & \emph{supervised} clustering \\
Grid-based (MAFIA) & arbitrary & follows from the data & noise drops out & sensitive to the grid
  granularity \\
DBSCAN & arbitrary & follows from the data & \emph{detects them} & parameters $\varepsilon$,
  $\tau$; struggles with clusters of differing density \\
\bottomrule
\end{tabular}
\end{center}

\subsection{Summary for the examination}

\textbf{Association rules.}
\begin{itemize}
\item \textbf{Support} $=$ the fraction of transactions with both $i$ and $j$; \textbf{confidence}
      $=$ the fraction of transactions with both among the transactions with $i$; \textbf{lift}
      $=p(i\wedge j)/\big(p(i)p(j)\big)$. Lift $<1$ $\Rightarrow$ the rule is worse than chance and
      negating the conclusion may be better. Be able to compute all three from a small transaction
      table.
\item \textbf{Apriori} uses the \emph{downward closure property} (any subset of a frequent itemset
      is frequent) and searches \emph{breadth-first}; \textsc{Apriori-Gen} $=$ the join step
      (combinations matching in $k-2$ categories) $+$ the pruning step (a subcombination of length
      $k-1$ is missing). The rule space is $O(2^m)$ and MBA produces an \emph{excessive number of
      rules}.
\item \textbf{Taxonomies and virtual items} address the choice of the item level; rare items go
      higher in the taxonomy, frequent ones stay lower; virtual items carry information about the
      transaction (day, time, customer) --- beware of redundant rules.
\item \textbf{MS-Apriori}: the rare item problem $\to$ each item has its own $\mathrm{MIS}$, the
      minsup of a rule is the \emph{minimum} MIS of its items, plus the support difference
      constraint $\varphi$. \emph{Downward closure ceases to hold} $\Rightarrow$ the items are
      sorted in ascending order of MIS, the set of seeds $L$ is used (not $F_1$) and
      \texttt{tailCount} is kept for rule generation. Be able to give the counterexample to
      downward closure.
\item \textbf{CAR}: transactions labelled with a class, rule $X\to y$; mined \emph{directly in one
      step} via rule-items $(\mathit{condset},y)$; different minsups can be given to different
      classes (101~\% $=$ prohibit a class).
\item \textbf{Sequential patterns}: a sequence $=$ an ordered list of itemsets; \emph{size} $=$ the
      number of elements, \emph{length} $=$ the number of items; subsequence via nested inclusions.
      The \textbf{GSP} algorithm works analogously to Apriori. Note:
      $\langle\{3\}\{8\}\rangle$ and $\langle\{3,8\}\rangle$ do \emph{not} contain each other.
\item \textbf{FP-Growth}: 2 passes over the data, no candidate generation; the FP-tree shares
      \emph{prefixes} (item order by decreasing support --- a heuristic) and keeps linked lists per
      item; extraction bottom-up by \emph{divide and conquer} via \emph{prefix path sub-trees} and
      \emph{conditional FP-trees} (update counts $\to$ remove the item's nodes $\to$ remove
      infrequent items). Worst case: unique transactions $\Rightarrow$ the tree is larger than the
      data.
\end{itemize}

\textbf{Supervised learning.}
\begin{itemize}
\item \textbf{Decision trees}: TDIDT (divide and conquer), attribute choice by \emph{minimum
      weighted entropy} $H(A)=\sum_v\frac{|A(v)|}{|D|}H(A(v))$, equivalently maximum
      \emph{information gain}. Be able to go through the loan example (income 0.5747; account
      0.6667; gender 0.9183; unemployed 0.825) and to convert the tree into rules. Induction
      complexity $O(np\log p)$, classification $O(q\log p)$.
\item \textbf{ID3} (Quinlan 1986) --- greedy, information gain. \textbf{C4.5} (1993) --- missing
      and continuous data, pruning by validation (\emph{reduced-error pruning}, \emph{subtree
      raising}), \emph{rule post-pruning}, accuracy estimated by the lower bound
      $\mathit{acc}-1.96\sqrt{\mathit{acc}(1-\mathit{acc})/n}$, the \emph{information gain ratio}
      (which penalises excessive branching). \textbf{C5.0} --- commercial, boosting.
      \textbf{CART} --- \emph{binary} trees, the Gini index $G(v_i)=1-\sum_j p_j^2$ (lower $=$
      better), the criterion $\Phi=2P_LP_R\sum_j|P(C_j|t_L)-P(C_j|t_R)|$.
\item \textbf{Bayes}: $H_{MAP}=\arg\max_H P(E\mid H)P(H)$ (the denominator can be neglected); the
      \emph{naive} assumption of conditional independence $\Rightarrow$
      $\arg\max_{H_t}P(H_t)\prod_i P(E_i\mid H_t)$. Handles incomplete patterns; the problem of
      zero probabilities. Be able to compute the example (0.1042 vs.\ 0.0416).
\item \textbf{SVM}: the maximal margin $2/\|w\|$, the primal $\min\frac{\langle w\cdot
      w\rangle}{2}$ subject to $y_i(\langle w\cdot x_i\rangle+b)\ge 1$; the Kuhn--Tucker conditions
      are, for a convex problem with linear constraints, necessary \emph{and} sufficient;
      $\alpha_i>0$ only for the \emph{support vectors} (this follows from \emph{complementarity}).
      \textbf{The Wolfe dual} arises in two steps: set $\partial L_P/\partial w$ and
      $\partial L_P/\partial b$ to zero (giving $w=\sum_i y_i\alpha_i x_i$ and
      $\sum_i y_i\alpha_i=0$) and substitute back into $L_P$, which eliminates $w,b$; it is a
      \emph{constrained} problem --- $\max L_D$ subject to $\sum_i y_i\alpha_i=0$,
      $\alpha_i\ge 0$. Back again: $w$ from stationarity,
      $b=y_k-\langle w\cdot x_k\rangle$ from complementarity.
      \textbf{The kernel trick}:
      the dual needs only dot products $\Rightarrow$
      $K(x,z)=\langle\phi(x)\cdot\phi(z)\rangle$ and $\phi$ is never evaluated (Mercer's theorem).
      Only 2 classes, only a real-valued space, not interpretable.
\item \textbf{Logistic regression}: $\ln\frac{p}{1-p}=\beta^{\!\top}x$, hence
      $p=\sigma(z)=1/(1+e^{-z})$; learning by maximum likelihood / minimising cross-entropy;
      $\sigma'=\sigma(1-\sigma)$ $\Rightarrow$ the update rule
      $\Delta\beta_j=\eta\sum_i(y_i-\sigma(z_i))x_{ij}$.
\item \textbf{Discriminant analysis}: a $g_j(x)$ for each class, classify by the maximum; for two
      classes the sign of $g=g_1-g_2$ suffices; the optimum is $g_j\propto P(C_j\mid x)$; a normal
      distribution $\to$ a quadratic function, with identical covariance matrices \emph{linear}
      $\Rightarrow$ only $\mu_j$ and $S$ need to be estimated.
\item \textbf{ELM}: random parameters of the hidden neurons (even without knowing the data), output
      weights by \emph{one} pseudoinverse $\beta=H^{+}T$; regularisation $I/C$ on the diagonal.
      Fast, no local minima, works with any bounded piecewise continuous transfer function.
\item \textbf{Ensembles}: \emph{bagging} reduces the \emph{variance} ($\sigma^2\to\sigma^2/k$ for
      i.i.d.), the bootstrap contains $\approx 63.2\,\%$ distinct points ($1-1/e$), trees are the
      ideal components (low bias, high variance); it does not reduce the bias. \emph{Random
      forests} address the correlation of the components --- the variance is
      $\rho\sigma^2+(1-\rho)\sigma^2/k$ and the first term does not depend on $k$; randomness comes
      from the bootstrap \emph{and} from a random subset of $m\approx\sqrt{\#\text{features}}$
      features at each node. \emph{Boosting}/\textbf{AdaBoost} reduces the \emph{bias}, the weights
      of misclassified patterns grow, $\alpha_t=\frac12\ln\frac{1-\epsilon_t}{\epsilon_t}$, it runs
      \emph{sequentially} and is sensitive to noise. Be able to compare forests vs.\ boosting
      (parallel and cheap vs.\ exhaustive and more accurate for the same number of trees).
\end{itemize}

\textbf{Cluster analysis.}
\begin{itemize}
\item \textbf{Metrics}: Hamming/Manhattan ($L_1$), Euclidean ($L_2$), Chebyshev ($L_\infty$),
      Minkowski ($L_q$, of which the others are special cases); \textbf{Mahalanobis}
      $d_M^2=(x_1-x_2)^{\!\top}S^{-1}(x_1-x_2)$ for quantities with different variance. Normalise
      before clustering.
\item \textbf{Cluster distance}: nearest neighbour (min), farthest neighbour (max), average
      distance, centroid method.
\item \textbf{$k$-means} (Lloyd 1982): random division $\to$ centroids $\to$ transfers $\to$ repeat
      while transfers occur. It never increases the cost but gives no approximation guarantee
      $\Rightarrow$ \textbf{$k$-means++} (probability $\propto$ the squared distance to the closest
      centroid), \textbf{LocalSearch++}.
\item \textbf{$k$-medians} (Manhattan, median per dimension, more robust) vs.\
      \textbf{$k$-medoids} (representatives \emph{from the data}, hill-climbing with exchanges,
      $O(knd)$ per iteration, any data type). \textbf{PAM} tests all $k(n-k)$ exchanges
      ($O(kn^2d)$); \textbf{CLARA} applies PAM to samples (risk of a bad sample),
      \textbf{CLARANS} tries random exchanges over the full data (\emph{MaxAttempt},
      \emph{MaxLocal}).
\item \textbf{Hierarchical clustering} bottom-up: every pattern its own cluster $\to$ merge the two
      closest; the output is a \textbf{dendrogram} and the number of clusters is chosen only by
      cutting it. \textbf{CURE}: a scalable variant (sampling, hierarchical merging per
      partition); fast, arbitrary cluster shapes.
\item \textbf{LVQ}: \emph{supervised} clustering --- the winning weight vector is pulled towards the
      pattern if the classes agree and pushed away if they do not; $\mu$ decreases as
      $\mu(k-1)/(k+1)$.
\item \textbf{Grid-based and density-based methods} (MAFIA, \textbf{DBSCAN}): clusters of
      arbitrary shape as connected dense regions; DBSCAN --- core points ($\ge\tau$ points
      within radius $\varepsilon$), clusters $=$ the connected components of the core-point
      graph, noise $=$ outliers.
\end{itemize}

%=====================================================================
\newpage
\section{Characteristic and discriminant rules and their interestingness}
\label{sec:pravidla}
%=====================================================================

The topic sentence asks for \emph{methods for the extraction of characteristic
discriminant rules and measuring their interestingness}. The chapter is built primarily
on the slides of the course \emph{Internet and classification methods} (Holeňa), above
all its 5th lecture \emph{When is a~classification comprehensible to the user?}: why
rules (Section~\ref{ssec:srozumitelnost}), extraction of rules from classification trees
(Section~\ref{ssec:extrakce}) and by evolution (Section~\ref{ssec:evoluce}),
observational rules (Section~\ref{ssec:observacni}), the statistical view
(Section~\ref{ssec:stat-pohled}), measuring interestingness
(Section~\ref{ssec:zajimavost}) and applications
(Section~\ref{ssec:aplikace-pravidel}). The notions of a~\emph{characteristic} and
a~\emph{discriminant} rule with the t-weight and d-weight measures
(Section~\ref{ssec:char-diskr}), together with \emph{attribute-oriented induction}
(Section~\ref{ssec:aoi}) which produces those rules, are added following Han--Kamber,
because the topic names them explicitly.

\begin{poznbox}[Sources of the chapter]
\small
Of the six lectures of the course \emph{Internet and classification methods}
(\texttt{ikm1}--\texttt{ikm6}), the key one for this chapter is \texttt{ikm5}
(comprehensibility of classification: rules, their extraction and the observational
calculus); the chapter further uses \texttt{ikm1} (applications: spam, recommending,
network threats), \texttt{ikm2} (quality measures of binary classification) and
\texttt{ikm3} (discriminant analysis). The material of \texttt{ikm4} (SVM: margin
maximisation, support vectors, the kernel trick) and \texttt{ikm6} (teams of
classifiers: bagging, boosting, random forests) coincides with
Chapter~\ref{sec:metody}, and the confusion matrix and ROC/AUC of \texttt{ikm2} with
Chapter~\ref{sec:reprez} --- they are treated there. The slides
\texttt{nmr1}--\texttt{nmr5} from the same folder belong to a~different (optional)
course by Holeňa and do not cover the subject matter of this chapter. Sections~\ref{ssec:aoi} and~\ref{ssec:char-diskr} follow Han--Kamber; regression and
discriminant analysis are also backed by the NDBI023 slides (\emph{Data Analysis}).
\end{poznbox}

\subsection{Rules as a comprehensible form of classification}
\label{ssec:srozumitelnost}

Mathematically, a~classifier is a~mapping $f\colon X\to\{C_1,\dots,C_m\}$ from the
feature space into a~finite set of classes; more precisely $f\colon
X\times\Theta\to\{C_1,\dots,C_m\}$, where learning means choosing the parameters
$\theta\in\Theta$ according to the experience with the quality of classification under
previous choices (Section~\ref{ssec:supervised}). The concrete learned prescription ---
say, the sign of $\langle w,x\rangle+b$ for an SVM or the composed nonlinear function of
a~neural network --- is, however, hardly comprehensible to an ordinary user. If a~binary
classifier $f\colon X\to\{1,0\}$ describes the truth of some implication
$\varphi\Rightarrow\psi$ or equivalence $\varphi\Leftrightarrow\psi$ between properties
of the object, it suffices to know this \emph{rule} instead of the whole $f$: a~rule of
logic is more comprehensible than a~mathematical formula.

\begin{defbox}[Boolean and fuzzy rules]
\textbf{Boolean} rules work with truth values $t(\varphi)\in\{0,1\}$: the
\emph{equivalence} $\varphi\Leftrightarrow\psi$ is true iff $t(\varphi)=t(\psi)$; the
\emph{implication} $\varphi\Rightarrow\psi$ is true iff $t(\varphi)=0$ or
$t(\varphi)=t(\psi)=1$. Rules are typically applied to objects on which the antecedent
holds ($t(\varphi)=1$) --- there the implication coincides with the equivalence.

\textbf{Fuzzy} rules admit degrees of truth $t(\varphi)\in[0,1]$. The fuzzy implication
is the \emph{residuum} of a~t-norm $T$:
\[ t(\varphi\Rightarrow\psi)=\sup\{\,c\in[0,1]: T(t(\varphi),c)\le t(\psi)\,\}, \]
and the fuzzy equivalence is $t(\varphi\Leftrightarrow\psi)
=T\big(t(\varphi\Rightarrow\psi),\,t(\psi\Rightarrow\varphi)\big)$. Important special
cases: if $t(\varphi)\le t(\psi)$, then $t(\varphi\Rightarrow\psi)=1$; if
$t(\varphi)=t(\psi)$, then $t(\varphi\Leftrightarrow\psi)=1$.
\end{defbox}

\textbf{Intuition.} A~fuzzy degree of truth is \emph{not} the fraction of objects on
which the rule holds (that is the confidence of an association rule, a~different layer)
but the degree to which \emph{one particular object} satisfies a~vague property. The rule
``whoever buys \emph{many} rolls buys \emph{a~lot of} butter'' is evaluated for each
customer separately. For Petr, 8~rolls are ``many'' to degree $0.7$ and 250\,g of butter
is ``a~lot'' to degree $0.9$; the antecedent $0.7 \le 0.9$, so the rule holds for him
\emph{fully}, $t = 1$. His butter could drop to $0.7$ and the rule would still hold
fully --- that slack is not added anywhere, truth is capped at one. The residuum handles
the opposite case: Jana has rolls ``many'' to $0.9$ but butter only to $0.4$, so she
violates the rule and the question is \emph{by how much}. In Boolean logic the rule would
simply be false; fuzzy logic says ``partly'', and the different t-norms are differently
strict ways of penalising the gap between antecedent and consequent (for
$t(\varphi) > t(\psi)$): minimum (Gödel) gives $t(\psi) = 0.4$, product (Goguen) gives
$t(\psi)/t(\varphi) \approx 0.44$, Łukasiewicz gives $1 - t(\varphi) + t(\psi) = 0.5$.
Only averaging these values over all customers yields the fuzzy counterpart of rule
confidence.

Comprehensibility is the reason why rules are used wherever the result of the
classification is judged by a~human: a~malware analyst wants to know \emph{why}
a~program is suspicious, and recommender systems increase the user's trust in
a~recommendation by explaining it (Section~\ref{ssec:aplikace-pravidel}).

\subsection{Attribute-oriented induction (AOI)}
\label{ssec:aoi}

\mimo{} This section follows Han--Kamber; the topic speaks of characteristic and discriminant
rules, and AOI is the method that produces them.

\textbf{Motivation.} Data in a~database sit at \emph{too low a~level of abstraction}. A~single
tuple \uv{John Newman, 23, Brno, physics} says nothing about a~class --- describing a~class
requires \emph{concepts}, not individual values. \textbf{Attribute-oriented induction} (AOI)
therefore \emph{lifts} the data to a~higher conceptual level and only then looks for
regularities.

\begin{defbox}[Concept hierarchy]
Background domain knowledge supplied by the user (or derived from the data) telling us how the
values of an attribute are to be generalised. For example
\[ \text{Brno} \to \text{South Moravia} \to \text{Czechia} \to \text{Europe} \to \mathrm{ANY} \]
The top \textbf{ANY} means \uv{the value does not matter} --- generalising this far effectively
removes the attribute. For numeric attributes the hierarchy consists of intervals
(age $21 \to 20$--$25 \to$ young) obtained by discretisation (Chapter~\ref{sec:priprava}).
\end{defbox}

\begin{defbox}[The two basic AOI operations]
\textbf{1. Attribute removal} --- applied when an attribute has \emph{many distinct values} and
\emph{no concept hierarchy is available} (name, birth number, telephone). Such an attribute
cannot be generalised, and keeping it would mean that no two tuples ever merge.

\textbf{2. Attribute generalization} --- if a~hierarchy exists, the values are replaced by the
concepts one level up. This is repeated until the number of distinct values drops below
a~chosen threshold.

After each step, \textbf{identical generalised tuples are merged} and their number is
accumulated into a~\texttt{count} attribute (possibly along with \texttt{sum}, \texttt{avg}).
\end{defbox}

\begin{defbox}[When to stop generalising --- two thresholds]
\textbf{Attribute generalization threshold}: each attribute may have at most $k$ distinct
values, typically 2--8. Generalisation continues until it falls below the threshold.

\textbf{Generalized relation threshold}: the resulting relation may have at most $k$ tuples,
typically 10--30. If it has more, generalisation continues (the attribute with the most values
is lifted one level).

The choice of thresholds is essential and it is a~\emph{trade-off}: thresholds that are too low
lead to \textbf{over-generalisation} --- the rule \uv{all clients are people} is true and
worthless; thresholds that are too high leave the result unreadable.
\end{defbox}

\noindent\textbf{The AOI algorithm:}
\begin{algorithmic}[1]
\Require relation $W$ (the result of a~query on the target class), concept hierarchies,
  thresholds
\Ensure the generalised (\emph{prime}) relation $P$
\ForAll{attributes $A$ of the relation $W$}
  \State determine the number of distinct values of $A$
  \If{$A$ has no hierarchy and many values}
    \State remove $A$ \Comment{attribute removal}
  \Else
    \While{number of distinct values of $A$ $>$ attribute threshold}
      \State replace the values of $A$ by the concepts one level up
        \Comment{attribute generalization}
    \EndWhile
  \EndIf
\EndFor
\State merge identical tuples, accumulate \texttt{count}
\While{$|P| >$ relation threshold}
  \State generalise further the attribute with the most distinct values and merge again
\EndWhile
\State \Return $P$
\end{algorithmic}

\medskip
\noindent\textbf{Example.} The raw student table has the attributes
\emph{name, birth number, gender, field, country, age, grade}.
\emph{Name} and \emph{birth number} are removed (many values, no hierarchy);
\emph{field} is generalised (physics, chemistry, biology $\to$ science),
\emph{country} (Czechia, Slovakia $\to$ domestic; otherwise foreign),
\emph{age} (21, 22, 23 $\to$ 20--25) and \emph{grade} (1, 2 $\to$ excellent).
Hundreds of rows collapse to a~few:

\begin{center}\small
\begin{tabular}{lllll r}
\toprule
\textbf{gender} & \textbf{field} & \textbf{country} & \textbf{age} & \textbf{grade}
  & \textbf{count} \\
\midrule
M & science   & domestic & 20--25 & excellent & 84 \\
F & science   & foreign  & 25--30 & excellent & 36 \\
M & humanities & domestic & 20--25 & good      & 30 \\
\bottomrule
\end{tabular}
\end{center}

\begin{poznbox}[Connections that get asked about]
\textbf{AOI vs.\ OLAP.} Generalisation corresponds to the \emph{roll-up} aggregation operation
over a~data-warehouse cube (the opposite direction being \emph{drill-down}). The difference:
OLAP works over a~\emph{precomputed} cube and is driven manually by the user, whereas AOI runs
\emph{automatically} over a~relation and decides itself which attributes to remove and which to
generalise.

\textbf{Characterisation vs.\ discrimination.} Run on the \emph{target class alone}, AOI yields
a~\emph{characterisation}; run on the \emph{target and the contrasting classes at once} (at the
same level of generalisation, so that they are comparable), it yields a~\emph{discrimination}.
The t-weight and d-weight are then computed over the resulting relation
(Sect.~\ref{ssec:char-diskr}).

\textbf{Complexity.} It is essentially a~single pass over the data with aggregation, hence
linear in the number of tuples --- which is why AOI is usable on large databases.

\textbf{Forms of presentation:} the generalised relation (the table above), a~cross-tab,
quantitative rules with weights, or charts (pie, bar).
\end{poznbox}

\subsection{Characteristic and discriminant rules: t-weight and d-weight}
\label{ssec:char-diskr}

\mimo{} This section follows Han--Kamber --- the topic names both kinds of rules
explicitly.

Both kinds of rules answer two different questions about the \emph{target class} (e.g.\
clients with a~mortgage). \textbf{Characterisation:} \uv{what does a~typical member of
the class look like?} \textbf{Discrimination:} \uv{how does it differ from the
\emph{contrasting class}?} (clients without a~mortgage). Both are read off the same
generalised relation produced by AOI (Sect.~\ref{ssec:aoi}): rows $q_1,\dots,q_N$ with
\texttt{count}s, for discrimination separately for the target and the contrasting class.
Each row gets one number for each of the two questions.

\begin{defbox}[t-weight --- typicality (characteristic rule)]
\textbf{Purpose:} measures how large a~part of the target class the row $q_a$ describes.
\textbf{Formula:}
\[ t\text{-}\mathrm{weight}(q_a)=\frac{\mathrm{count}(q_a)}
   {\sum_{i=1}^{N}\mathrm{count}(q_i)}, \]
i.e.\ the number of target-class objects in the row divided by the size of the whole
target class; the sum over the rows is 100~\%. \textbf{Meaning:} \uv{45~\% of clients
with a~mortgage are 30--40 years old, high-income, from Prague.} It appears in the
\emph{characteristic rule}, which goes \emph{from} the class to the properties
(a~necessary condition):
\[ \mathrm{target}(X)\Rightarrow \mathrm{cond}_1(X)\,[t{:}\,w_1]\ \vee\ \cdots\ \vee\
   \mathrm{cond}_N(X)\,[t{:}\,w_N]. \]
It says nothing about whether the property \emph{distinguishes} the class from the others.
\end{defbox}

\begin{defbox}[d-weight --- discriminating power (discriminant rule)]
\textbf{Purpose:} measures how reliably the row $q_a$ points to the target class rather
than to a~contrasting one. \textbf{Formula:}
\[ d\text{-}\mathrm{weight}(q_a)=\frac{\mathrm{count}(q_a\in C_t)}
   {\mathrm{count}(q_a\in C_t)+\sum_{j=1}^{m}\mathrm{count}(q_a\in C_j)}, \]
i.e.\ of all objects (target and contrasting) that satisfy the row, the fraction from the
target class --- the precision of the row as a~classifier. \textbf{Meaning:} \uv{90~\% of
people aged 30--40 with a~high income from Prague have a~mortgage.} It appears in the
\emph{discriminant rule}, which goes from the properties \emph{into} the class
(a~sufficient condition):
\[ \mathrm{target}(X)\Leftarrow \mathrm{cond}(X)\,[d{:}\,w]. \]
Always compare it with the \emph{a-priori proportion} of the target class: a~d-weight
equal to the a-priori proportion means zero discriminating power, a~d-weight below it
rather excludes the class. The \emph{quantitative description rule} states both weights
at once ($\Leftrightarrow$, $[t{:}\,w,\ d{:}\,w']$).
\end{defbox}

\begin{defbox}[Example --- clients with a mortgage (200) vs.\ without one (800)]
The generalised relation:
\begin{center}
\small
\begin{tabular}{@{}lllrrrr@{}}
\toprule
Age & Income & Region & \multicolumn{2}{c}{count} & $t$-weight & $d$-weight \\
 & & & target & contrast & (target) & \\
\midrule
30--40 & high & Prague & 90 & 10 & \textbf{45~\%} & \textbf{90~\%} \\
30--40 & middle & Prague & 50 & 150 & 25~\% & 25~\% \\
40--50 & high & outside Prague & 40 & 160 & 20~\% & 20~\% \\
20--30 & low & outside Prague & 20 & 480 & 10~\% & 4~\% \\
\bottomrule
\end{tabular}
\end{center}
First row: $t=90/200=45\,\%$ (almost half of all mortgages are here),
$d=90/(90+10)=90\,\%$ (whoever falls here has a~mortgage with 90~\%).
\textbf{The trap:} the 3rd row has a~decent $t=20\,\%$, but $d=40/200=20\,\%$ is exactly
the a-priori proportion of mortgages ($200/1000$) --- typical, yet it distinguishes
nothing; the 4th row with $d=4\,\%$ rather excludes a~mortgage.
\end{defbox}

\subsection{Extraction of rules from classification trees}
\label{ssec:extrakce}

The most direct method of rule extraction starts from a~classification tree: the nodes
of a~binary tree test Boolean conditions (typically $x_i\le c$, i.e.\ cuts perpendicular
to the coordinate axes), the input falls left or right according to the outcome of the
test, and every leaf corresponds to a~single class.

\begin{defbox}[Rules from a tree]
\begin{itemize}
\item To every root $\to$ leaf path corresponds a~\emph{conjunction} of the conditions
      tested along the path; the class of the leaf provides the consequent:
      $\big(\mathrm{cond}_1\wedge\cdots\wedge\mathrm{cond}_k\big)\Rightarrow C_j$.
\item Paths to leaves of the \emph{same} class are joined by a~\emph{disjunction} ---
      the class is the union of the rectangular regions delimited by the individual
      paths.
\end{itemize}
\end{defbox}

\textbf{Learning the tree} is an optimisation of the split in a~node with respect to
the choice of the variable and the splitting point (subspaces of continuous inputs,
subsets of discrete ones); the two most frequently optimised criteria are the
\emph{information gain}
$IG(\text{parent};L,R)=H_{\text{parent}}-\big(\tfrac{n_L}{n}H_L+\tfrac{n_R}{n}H_R\big)$
and the \emph{Gini impurity index} $1-\sum_j p_j^2$; in detail (TDIDT, ID3, C4.5, CART)
see Section~\ref{ssec:supervised}. New inputs are classified by passing the node tests;
in the leaves, the distribution of the new inputs is assumed to match the distribution
of the training data.

\textbf{Pruning} reduces overfitting and at the same time \emph{shortens} the extracted
rules (cf.\ rule post-pruning in C4.5, Section~\ref{ssec:supervised}).
\textbf{Generalised splits} --- a~general hyperplane or a~quadratic surface instead of
an axis-perpendicular cut --- improve the approximation of the class boundary but
reduce the comprehensibility of the rules obtained as conjunctions of conditions.

\subsection{Evolutionary construction of rules}
\label{ssec:evoluce}

Rules can also be constructed by a~genetic algorithm --- the computer-science
application of the biological evolution of species: crossover and mutation of a~string
(genome) encoding the rule reinforce the desired properties, here the classification
quality measures (accuracy AC, precision PR, TPr, TNr or AUC ---
Section~\ref{ssec:zajimavost}). The approach is suitable when both the antecedent and
the consequent are conjunctions of a~small number of atoms. All evolutionary algorithms
are \emph{population-based}; according to what constitutes an individual of the
population, two approaches are distinguished:

\begin{center}
\small
\begin{tabular}{@{}lp{5.2cm}p{6.8cm}@{}}
\toprule
\textbf{Approach} & \textbf{Individual} & \textbf{Properties} \\
\midrule
Michigan & a~single rule; the whole population is sought & simple, but the population
  needlessly converges to similar rules \\
Pittsburgh & the entire sought set of rules & the result is directly a~consistent set;
  the evolution of a~population of sets is more complex and computationally demanding \\
\bottomrule
\end{tabular}
\end{center}

\subsection{Observational rules: confidence, support and hypothesis tests}
\label{ssec:observacni}

Classical logic can say about data only \uv{$\varphi\Rightarrow\psi$ holds for
\emph{all} objects} (almost never) or \uv{it holds for at least \emph{one}} (almost
nothing). The \emph{observational calculus} (Hájek--Havránek 1978, the basis of the
Czech GUHA method) therefore replaces the quantifiers $\forall$ and $\exists$ by
a~\emph{generalised quantifier} $\sim$ meaning \uv{holds for \emph{enough} objects},
and defines \uv{enough} statistically. An observational rule $\varphi\sim\psi$ is thus
the statement \uv{in the data, $\varphi$ and $\psi$ occur together sufficiently often},
and the individual quantifiers differ only in how they measure \uv{sufficiently}.

\begin{defbox}[Four-fold table --- the only input of all quantifiers]
For a~pair $\varphi,\psi$ four numbers are counted from the data:
\begin{center}
\begin{tabular}{@{}lcc@{}}
\toprule
 & $\psi$ & $\neg\psi$ \\
\midrule
$\varphi$ & $a$ & $b$ \\
$\neg\varphi$ & $c$ & $d$ \\
\bottomrule
\end{tabular}
\end{center}
$a$ = objects satisfying both, $b$ = only $\varphi$, $c$ = only $\psi$, $d$ = neither,
$n=a+b+c+d$. Every generalised quantifier is just a~condition on $a,b,c,d$.
\end{defbox}

\begin{defbox}[Type 1 --- confidence and support (association rules)]
\textbf{Purpose:} \uv{if $\varphi$, then often enough also $\psi$, and this concerns
a~large enough part of the data.} \textbf{Formula:}
\[ \mathrm{conf}=\frac{a}{a+b}\ \ (\text{estimate of } P(\psi\mid\varphi)),\qquad
   \mathrm{sup}=\frac{a}{n}, \]
\[ (\varphi\sim\psi)=1 \iff \mathrm{conf}\ge c_0\ \wedge\ \mathrm{sup}\ge s_0 . \]
\textbf{Meaning:} confidence is the \emph{reliability} of the rule (of the baskets with
rolls, how many also contain butter), support is its \emph{significance} (how many
baskets it concerns at all). The thresholds $c_0, s_0$ are chosen by the user. Weakness:
a~confidence of 60~\% means nothing if 60~\% of \emph{all} customers buy butter --- the
confidence must be compared with $\tfrac{a+c}{n}$, which leads to type~2.
\end{defbox}

\begin{defbox}[Type 2 --- hypothesis test]
\textbf{Purpose:} instead of a~fixed threshold we ask whether the joint occurrence of
$\varphi$ and $\psi$ could be \emph{chance}. \textbf{Formula:} $H_0$: $\varphi$ and
$\psi$ are independent, $P(\psi\mid\varphi)=P(\psi)$; $H_1$: positive dependence,
$P(\psi\mid\varphi)>P(\psi)$. The one-sided Fisher test
(Section~\ref{ssec:kontingencni}) computes the $p$-value = the probability that, under
independence and with the given marginal sums of the table, $a$ would come out at least
as large as observed;
\[ (\varphi\sim\psi)=1 \iff p\le\alpha \quad(\text{e.g.\ } \alpha=0.05). \]
\textbf{Meaning:} the rule holds when $\varphi$ and $\psi$ occur together
\emph{demonstrably} more often than chance would explain. The test automatically
accounts for both the strength of the dependence and the number of objects: with
a~small $n$ even a~high confidence is not enough.
\end{defbox}

\textbf{Example.} $n=1000$ baskets, rolls in 100 of them, butter in 400; together
$a=60$, hence $b=40$, $c=340$, $d=560$. Confidence $60/100=60\,\%$, support
$60/1000=6\,\%$; with thresholds $c_0=0.5$, $s_0=0.05$ the rule holds. Overall $40\,\%$
of customers buy butter, among roll buyers $60\,\%$ --- under independence we would
expect $a\approx40$, we observed 60, the Fisher test gives $p\ll0.05$ and the rule holds
by type~2 as well.

\subsection{The statistical view: regression and discriminant analysis}
\label{ssec:stat-pohled}

A~characteristic and discriminant description of classes need not take the form of
logical rules --- classical statistics solves the same two tasks by \emph{regression}
and \emph{discriminant analysis}. The context: classification methods are divided
according to whether they primarily \emph{seek} the boundary between the classes
(neural networks, decision trees, SVM --- modern, computationally demanding), or do
\emph{not} seek it and start from the experience of similar cases or from the class
probabilities ($k$-NN, Bayesian classification, discriminant analysis --- traditional,
computationally cheap).

\textbf{Regression analysis} determines the relationship of a~variable $Y$ to one or
more variables $X_1,\dots,X_n$: linear regression of one variable by the least squares
method, multi-dimensional linear regression, and logistic regression for
a~categorical output --- in detail in Sections~\ref{ssec:regrese}
and~\ref{ssec:logreg}. A~regression model is the quantitative counterpart of
a~\emph{characteristic} description: it says which output values are typical for the
given inputs.

\textbf{Discriminant analysis} (DA) classifies patterns into \emph{known} classes but
does not seek the boundary: to every class $C_j$ it \emph{fits} a~probability
distribution from a~pre-chosen family of densities
$F=\{f(\cdot\mid\theta):\theta\in\Theta\}$ --- in Fisher's original DA the multivariate
normal distribution. Learning is the maximum likelihood estimation of the parameters
$\theta_j$ from the training data of the class $C_j$; a~new pattern $x$ is classified
by $\arg\max_j p_j\,f(x\mid\theta_j)$, where $p_j$ is the prior probability of the
class.

\begin{defbox}[Linear vs.\ quadratic discriminant analysis]
\textbf{LDA}: the covariance matrix $\Sigma$ is given in advance and shared by all
classes; the discriminant function is
\[ g_j(x)=\ln p_j-\tfrac12\,(x-\mu_j)^{\!\top}\Sigma^{-1}(x-\mu_j), \]
only $\mu_j$ is estimated ($q$ parameters per class for $x\in\mathbb{R}^q$); the set
$\{x: g_j(x)=g_k(x)\}$ between two classes is a~\emph{hyperplane}.

\textbf{QDA}: in addition, each class has its own covariance matrix $\Sigma_j$ (another
$\tfrac{q(q+1)}{2}$ parameters per class); the boundary between the classes is
a~\emph{quadratic surface}.
\end{defbox}

For two classes, instead of $g_1,g_2$ one can seek a~single function $g=g_1-g_2$ and
classify by its \emph{sign} --- the condition $g(x)>0$ is the quantitative counterpart
of a~\emph{discriminant rule} (a~sufficient condition for membership of the first
class). For the derivation of the linear and quadratic discriminant functions from
normal distributions see Section~\ref{ssec:da}.

\begin{center}
\small
\begin{tabular}{@{}lll@{}}
\toprule
\textbf{Task} & \textbf{Rule form} & \textbf{Statistical form} \\
\midrule
characterisation & characteristic rule ($\Rightarrow$, t-weight) & regression analysis \\
discrimination & discriminant rule ($\Leftarrow$, d-weight) & discriminant analysis \\
\bottomrule
\end{tabular}
\end{center}

\subsection{Measuring the interestingness of rules}
\label{ssec:zajimavost}

Extraction typically yields more rules than the user can bear --- the rules have to be
measured, sorted and filtered. A~rule $\varphi\Rightarrow\psi$ is itself a~binary
classifier, so the quality measures of binary classification apply to it
(Section~\ref{ssec:eval}): accuracy AC, precision (predictive value) PR, sensitivity
(recall) TPr, specificity TNr $=1-{}$FPr, the F-measure
$FM=\frac{2\,PR\cdot TPr}{PR+TPr}$ and AUC. For observational rules they correspond to
the \emph{confidence} $\tfrac{a}{a+b}$ (the precision of the rule) and the
\emph{support} $\tfrac{a}{n}$; for characteristic and discriminant rules, the d-weight
is the precision of the rule and the t-weight its recall with respect to the target
class.

The roles of the measures in working with rules:
\begin{itemize}
\item the \emph{fitness} in the evolutionary construction of rules
      (Section~\ref{ssec:evoluce});
\item \emph{sorting}: the rules for malware families are ordered by precision
      (Section~\ref{ssec:aplikace-pravidel});
\item \emph{filtering}: removal of redundant and imprecise rules;
\item \emph{thresholding}: an observational rule is \uv{interesting} if its confidence
      and support exceed the chosen thresholds, or if the test rejects independence
      (Section~\ref{ssec:observacni}) --- every choice of a~measure and a~threshold is
      a~different generalised quantifier, and hence a~different set of rules found in
      the data.
\end{itemize}

\textbf{Different costs of different errors.} Traditional accuracy and error rate
assume that every error matters equally --- unrealistically: an undelivered ham
matters more than an unrecognised spam, an unrecognised network intrusion more than
a~false alarm. The generalisation of the error rate is the \emph{mean cost}
$\frac1n\sum_{i,j}c(i,j)\,n_{ij}$, where $c(i,j)$ is the cost of classifying class $i$
as $j$ (the traditional error rate: $c(i,i)=0$, otherwise 1); the interestingness of
a~rule is then measured by the cost of the errors it commits.

\mimoq{subjective measures following Han--Kamber} Besides the objective measures,
\emph{subjective} ones decide: comprehensibility (the whole point of rule extraction),
simplicity (the length of the conjunction, the number of rules), unexpectedness
(a~contradiction with the user's expectation) and actionability (one can act on the
rule). In practice: a~rule is interesting if it is at the same time comprehensible,
valid on new data, potentially useful and novel.

\subsection{Applications: spam, malware and network threats, recommending}
\label{ssec:aplikace-pravidel}

The three internet applications of classification methods from the 1st lecture of the
course illustrate the role of rules:

\textbf{Spam filtering} is a~classification \{combinations of feature values\} $\to$
\{spam, ham\}: features from the content of the message (words, combinations of words,
images) and from the meta-information (domain, encoding, a~forged sender address;
anomalies of the SMTP session; the same message sent to many recipients). If low FP and
FN are not attainable simultaneously, a~third class between spam and ham helps ---
quarantine/uncertainty (e.g.\ SpamBayes). Rule-based filters work mainly with the
meta-information; a~plain Boolean disjunction of the rules leads to high false
positivity (\uv{graphical body of the message $\Rightarrow$ automatically spam}), the
Łukasiewicz fuzzy disjunction of rules with the same consequent works better. And FP
(an undelivered ham) matters more than FN.

\textbf{Malware and network threat detection.} Suspicious programs are in the end
evaluated by an analyst --- who wants to understand \emph{why} they are suspicious,
hence rules. Many of them arise, so redundant and imprecise ones are removed; typical
is a~hierarchical classification with two layers of rules (the first precise for
normal software and sensitive for malware, the second sorting malware families,
ordered by precision). For network threats (R2L, U2R, DoS and probe attacks), the
counterpart of rules are the \emph{signatures} of known attacks --- patterns with
a~certain (one hundred per cent) conclusion.

\textbf{Recommender systems.} In collaborative filtering by the user's behaviour, rules
of the form \uv{the user visited the pages $S_1,\dots,S_n$ $\Rightarrow$ recommend the
object $O$ (more generally: offer the next page $S_d$)} are used. They are obtained by
a~statistical analysis of clickstream data --- probability estimates and tests of the
validity of models (hypotheses), from whose results the rules are constructed
heuristically; they are thus observational rules of both types from
Section~\ref{ssec:observacni}. Recommender systems often explain \emph{why} they
recommend an object --- the explanation increases the comprehensibility of the
recommendation for the user.

\subsection{Summary for the examination}

\begin{itemize}
\item \textbf{Why rules}: a~classifier $f\colon X\to\{C_1,\dots,C_m\}$ is
      a~mathematical prescription, hardly comprehensible; a~rule of logic
      ($\varphi\Rightarrow\psi$, $\varphi\Leftrightarrow\psi$) is more comprehensible.
      The Boolean implication is true at $t(\varphi)=0$ or $t(\varphi)=t(\psi)=1$; the
      fuzzy implication $=$ the residuum of a~t-norm, $t(\varphi)\le
      t(\psi)\Rightarrow$ the implication is true to degree 1.
\item \textbf{AOI --- attribute-oriented induction}: lifts the data to a~higher conceptual
      level following \emph{concept hierarchies} and only then looks for regularities. Two
      operations: attribute \emph{removal} (many values and no hierarchy --- name, birth
      number) and attribute \emph{generalization} (values $\to$ concepts one level up),
      with identical tuples merged after each step and a~\texttt{count} accumulated. When
      to stop: a~threshold on the number of distinct values of an \emph{attribute} (2--8)
      and a~threshold on the number of tuples of the \emph{relation} (10--30); thresholds
      that are too low lead to \textbf{over-generalisation} (\uv{all clients are people}).
      It corresponds to the \emph{roll-up} operation in OLAP but runs automatically over
      a~relation. Run on the target class it yields a~\emph{characterisation}, run on the
      target and the contrasting classes a~\emph{discrimination}.
\item \textbf{Characteristic vs.\ discriminant rule} (Han--Kamber): characterisation
      describes the target class ($\Rightarrow$, a~necessary condition,
      \textbf{t-weight} $=\frac{\mathrm{count}(q_a)}{\sum_i\mathrm{count}(q_i)}$);
      discrimination distinguishes it from the contrasting ones ($\Leftarrow$,
      a~sufficient condition, \textbf{d-weight} $=$ the fraction of the covered objects
      belonging to the target class). Be able to compute both from a~table; always
      compare the d-weight with the \emph{a-priori} class proportion.
\item \textbf{Extraction from a tree}: a~root $\to$ leaf path $=$ a~conjunction of
      conditions, paths to leaves of the same class $=$ a~disjunction; learning by
      splitting according to the information gain or the Gini index; pruning shortens
      the rules; generalised (oblique) splits reduce comprehensibility.
\item \textbf{Evolution of rules}: a~genetic algorithm, fitness $=$ AC/PR/TPr/TNr/AUC;
      the Michigan approach (individual $=$ a~rule) vs.\ the Pittsburgh approach
      (individual $=$ a~set of rules).
\item \textbf{Observational rules} (Hájek--Havránek 1978, GUHA): $\varphi\sim\psi$ with
      a~generalised quantifier over a~four-fold table $(a,b,c,d)$; (1)~probability
      estimates --- association rules, $\frac{a}{a+b}\ge c_0$ (confidence) $\wedge$
      $\frac{a}{n}\ge s_0$ (support); (2)~hypothesis tests --- the one-sided Fisher
      test rejects independence in favour of positive dependence.
\item \textbf{Statistically}: regression analysis $=$ a~characteristic description of
      the relationship of $Y$ to $X_1,\dots,X_n$ (least squares, multi-dimensional,
      logistic); discriminant analysis $=$ classification into known classes by
      fitting distributions: ML estimation of $\theta_j$, prediction
      $\arg\max_j p_j f(x\mid\theta_j)$; \textbf{LDA} (shared $\Sigma$, hyperplane
      boundary, $q$ parameters) vs.\ \textbf{QDA} (per-class $\Sigma_j$, quadratic
      boundary, $+\frac{q(q+1)}{2}$ parameters); two classes $\to$ the sign of
      $g=g_1-g_2$.
\item \textbf{Interestingness}: a~rule $=$ a~binary classifier $\to$ AC, PR
      ($=$ confidence), TPr (recall), TNr, F-measure, AUC; roles: fitness, sorting
      (malware by precision), filtering out redundant rules, thresholding (confidence
      and support $=$ a~generalised quantifier); different costs of errors (the mean
      cost). Subjectively: comprehensibility, simplicity, unexpectedness,
      actionability.
\item \textbf{Applications}: spam (meta-information, the 3rd class quarantine, the
      Łukasiewicz disjunction against FP), malware (a~two-layer hierarchy of rules),
      network threats (signatures), recommending (rules from clickstream data).
\end{itemize}
%=====================================================================
\newpage
\section{Representation, evaluation and visualisation of knowledge}
\label{sec:reprez}
%=====================================================================

The chapter has three parts, matching the three words in the sentence of the topic: \emph{in what
form} we express the knowledge, \emph{how we verify} that it is good, and \emph{how we display} it.
The middle part --- the evaluation of classifiers --- is the only one the slides cover in full,
and it is at the same time the one most frequently asked about in the examination.

\subsection{Forms of knowledge representation}

\mimoq{the table as an overview; the individual methods in it are backed by the slides}
\begin{center}
\small
\begin{tabular}{@{}p{3.4cm}p{4.4cm}p{7.0cm}@{}}
\toprule
\textbf{Form} & \textbf{Typical source} & \textbf{Properties} \\
\midrule
IF-THEN rules & association rules, C4.5 & the most comprehensible; local (each rule is read on its
  own); a risk of conflicts and over-production \\
Decision tree & ID3/C4.5/CART & a global model, easily readable up to depth 4--5; unreadable
  beyond \\
Decision table & discretised data & a complete enumeration of the combinations; readable only for
  few attributes \\
Mathematical function & regression, SVM, NN & compact and accurate; hard to interpret (except a
  linear model with normalised weights) \\
Prototypes / centroids & $k$-means, $k$-medoids, LVQ & the \uv{typical representative} of a group;
  very intuitive for business \\
Probabilistic model & naive Bayes, Bayesian network & explicit uncertainty; the network graph is
  readable \\
Graphs and networks & social network analysis & relations between entities; visually strong \\
Instances (case memory) & $k$-NN, CBR & the \uv{model} is the data itself; explanation by reference
  to similar cases \\
\bottomrule
\end{tabular}
\end{center}

\begin{defbox}[Interpretability and explainability]
\mimoq{explainability and the post-hoc methods; interpretability as a criterion for evaluating a
model is in the slides --- see Section~\ref{ssec:eval}}
\emph{Interpretability} --- a human's ability to understand how the model arrived at its output.
\emph{Explainability} --- the ability to give a comprehensible justification of a concrete
decision.

We distinguish \emph{inherently interpretable models} (linear models, decision trees, rule sets,
prototypes) and \emph{post-hoc explanations of black boxes} --- attribute importance (permutation,
SHAP), local surrogate models (LIME), partial dependence plots (PDP), rule extraction from a
neural network, sensitivity analysis, counterfactual explanations (\uv{had your income been 5000
higher, the loan would have been approved}).

Interpretability is not merely a convenience: in regulated domains (loans, insurance, medicine) it
is a legal requirement (GDPR Art.~22 --- the right to an explanation of an automated decision; the
AI Act for high-risk systems).
\end{defbox}

\subsection{Evaluation of models}
\label{ssec:eval}

\subsubsection{Criteria for evaluating classification methods}

\begin{itemize}
\item \textbf{Predictive accuracy} --- the fraction of correctly classified patterns:
      \[ \mathit{Accuracy}=\frac{\text{number of correct classifications}}
         {\text{number of all test cases}}. \]
\item \textbf{Efficiency} --- the time to construct the model and the time to use it.
\item \textbf{Robustness} --- handling noise and missing values.
\item \textbf{Scalability} --- efficiency for disk-resident databases.
\item \textbf{Interpretability} --- the clarity and insight provided by the model.
\item \textbf{Compactness of the model} --- the size of the tree, the number of rules, \dots
\end{itemize}

\subsubsection{Methods of error estimation}

\begin{defbox}[Holdout, cross-validation, leave-one-out, validation set]
\textbf{Holdout set}: the available set $D$ is divided into two disjoint subsets --- the
\emph{training} set $D_{train}$ (for learning the model) and the \emph{test} set $D_{test}$ (for
testing it). \textbf{The essential rule:} the training set must not be used in testing and the test
set must not be used in learning --- only an unseen test set provides an \emph{unbiased} estimate
of accuracy. Used mainly for a \emph{large} $D$.

\textbf{$k$-fold cross-validation}: the data is partitioned uniformly into $k$ equal-size disjoint
subsets. Each subset is used as a test set and the remaining $k-1$ are combined as the training set.
The procedure is run $k$ times, yielding $k$ accuracies $\mathit{Accuracy}_i$; the final estimate
is their average
\[ \mathit{Accuracy}_{CV}=\frac1k\sum_{i=1}^{k}\mathit{Accuracy}_i. \]
10-fold and 5-fold cross-validation are common; the method is used rather for \emph{smaller} data.

\textbf{Leave-one-out}: a special case of cross-validation where each fold has only a single test
example and all the rest is used in training --- for $m$ examples this is $m$-fold
cross-validation. Used for \emph{very small} data.

\textbf{Validation set}: the data is divided into \emph{three} subsets --- a training, a validation
and a test set. The validation set is used for estimating the \emph{parameters} of learning
algorithms: the values that give the best accuracy on it are used as the final parameter values.
Cross-validation can be used for parameter estimation as well.
\end{defbox}

\begin{defbox}[Confidence interval for cross-validation]
The approximate $N\%$ confidence interval:
\[ \mathit{Accuracy}_{CV}\pm t_{N,k-1}\cdot s_{\mathit{Accuracy}},\qquad
   s_{\mathit{Accuracy}}=\sqrt{\frac{1}{k(k-1)}\sum_{i=1}^{k}
   \big(\mathit{Accuracy}_i-\mathit{Accuracy}_{CV}\big)^2}. \]
Values of $t_{N,k-1}$ for two-sided confidence intervals:
\begin{center}
\small
\begin{tabular}{@{}lrrr@{}}
\toprule
& \multicolumn{3}{c}{confidence level $N$} \\
$k$ & 90~\% & 95~\% & 99~\% \\
\midrule
2 & 2.92 & 4.30 & 9.92 \\
10 & 1.81 & 2.23 & 3.17 \\
30 & 1.70 & 2.04 & 2.75 \\
$\infty$ & 1.64 & 1.96 & 2.58 \\
\bottomrule
\end{tabular}
\end{center}
Often $t=1.96$ is used (for $k\to\infty$ and $N=95\,\%$).
\end{defbox}

\subsubsection{The confusion matrix and the derived measures}

Accuracy is only \emph{one} measure ($\mathit{error}=1-\mathit{accuracy}$) and in some applications
it is not suitable. In text mining we may be interested only in the documents of a particular
topic, which are a small portion of a big collection. In classification involving strongly
\emph{imbalanced} data (network intrusion, financial fraud detection) we are interested only in the
minority class: with 1~\% intrusions we achieve 99~\% accuracy \emph{by doing nothing} --- and
detect no intrusion at all. The class of interest is called the \emph{positive} class, the rest
\emph{negative}.

\begin{defbox}[Confusion matrix, precision, recall, $F_1$]
\begin{center}
\small
\begin{tabular}{@{}llcc@{}}
\toprule
& & \multicolumn{2}{c}{\textbf{classified as}} \\
& & positive & negative \\
\midrule
\textbf{actual} & positive & TP & FN \\
                & negative & FP & TN \\
\bottomrule
\end{tabular}
\end{center}
$TP$ --- the number of correct classifications of the positive examples (true positive); $FN$ ---
incorrect classifications of positive examples (false negative); $FP$ --- incorrect classifications
of negative examples (false positive); $TN$ --- correct classifications of negative examples.

\[ p=\mathit{precision}=\frac{TP}{TP+FP},\qquad r=\mathit{recall}=\frac{TP}{TP+FN}. \]
\emph{Precision} $p$ is the number of correctly classified positive examples divided by the total
number of examples classified as positive. \emph{Recall} $r$ is the number of correctly classified
positive examples divided by the total number of actual positive examples in the test set.
\textbf{Both measures evaluate the classification only on the positive class.}

The \textbf{$F_1$ measure} combines both into one --- it is their \emph{harmonic mean}:
\[ F_1=\frac{2pr}{p+r}=\frac{2}{\frac1p+\frac1r}. \]
The harmonic mean of two numbers tends to be closer to the \emph{smaller} of the two, so $F_1$ is
large only if \emph{both} measures are large.
\end{defbox}

\emph{An illustration of the trap:} a classifier that correctly labels a single positive example and
labels no negative one wrongly has $p=100\,\%$ and $r=1\,\%$. High precision alone therefore says
nothing about usefulness.

\subsubsection{Scoring, ranking and lift analysis}

\emph{Scoring} is related to classification: we are interested in a single (positive) class ---
e.g.\ the buyers class in a marketing database. Instead of assigning each test instance a definite
class, scoring assigns a \emph{probability estimate} (PE) that it belongs to the positive class.
Classification systems can be used for scoring, but they must produce a probability estimate ---
e.g.\ in decision trees the confidence value at each leaf node can be used as the score.

After scoring, the examples are \emph{ranked} by their PE and divided into $n$ (say 10) bins; a
\textbf{lift curve} is drawn according to how many positive examples are in each bin --- this is
\emph{lift analysis}.

\begin{defbox}[Example --- lift analysis for a targeted campaign]
We want to send promotion materials to potential customers (a watch). Each package costs \$0.50 to
send; if a watch is sold, we make \$5 profit. How many packages should we send, and to whom?

The test set has 10\,000 instances, out of which 500 are positive. After scoring and ranking into
10 bins (1000 instances each):

\begin{center}
\small
\begin{tabular}{@{}lrrrrrrrrrr@{}}
\toprule
Bin & 1 & 2 & 3 & 4 & 5 & 6 & 7 & 8 & 9 & 10 \\
\midrule
Positives & 210 & 120 & 60 & 40 & 22 & 18 & 12 & 7 & 6 & 5 \\
\% of total & 42 & 24 & 12 & 8 & 4.4 & 3.6 & 2.4 & 1.4 & 1.2 & 1.0 \\
Cumulative & 42 & 66 & 78 & 86 & 90.4 & 94 & 96.4 & 97.8 & 99 & 100 \\
\bottomrule
\end{tabular}
\end{center}

The first bin contains 42~\% of all positive cases, the first three 78~\%. Random mailing would give
10~\% in each bin --- the gap between the lift curve and this diagonal is the model's contribution.
The decision \uv{how many packages} then follows from the economics: mailing to the first bin earns
$210\cdot\$5-1000\cdot\$0.50=\$550$, to the tenth $5\cdot\$5-1000\cdot\$0.50=-\$475$.
\end{defbox}

\subsubsection{The ROC curve and AUC}

\begin{defbox}[The ROC curve (Receiver Operating Characteristic)]
A plot of the \emph{true positive rate} (TPR) against the \emph{false positive rate} (FPR):
\[ \mathit{TPR}=\frac{TP}{TP+FN},\qquad \mathit{FPR}=\frac{FP}{FP+TN},\qquad
   \mathit{TNR}=\frac{TN}{FP+TN}. \]
$\mathit{TPR}$ is the fraction of actual positive cases that are correctly classified --- it is
essentially the \emph{recall} of the positive class, also called \emph{sensitivity}.
$\mathit{FPR}$ is the fraction of actual negative cases classified as positive. $\mathit{TNR}$
(the recall of the negative class) is called \emph{specificity}, and
$\mathit{FPR}=1-\text{specificity}$.

Each ROC curve starts at $(0,0)$ (every case classified as negative) and ends at $(1,1)$ (every
case classified as positive). \textbf{The main diagonal} represents random guessing --- predicting
each case to be positive with a fixed probability, when $\mathit{TPR}=\mathit{FPR}$.
\end{defbox}

\textbf{Drawing a ROC curve} given the classification as a ranking of test cases:
\begin{enumerate}
\item Obtain the ranking by sorting the test cases in \emph{decreasing} order of the classifier's
      output values (e.g.\ posterior probabilities).
\item Partition the rank list into two subsets \emph{at every} test case, and regard every test case
      in the upper part as a positive case and every one in the lower part as negative.
\item For each such partition, compute a pair $(\mathit{FPR},\mathit{TPR})$. When the upper part is
      empty we obtain the point $(0,0)$; when the lower part is empty, the point $(1,1)$.
\item Connect the adjacent points.
\end{enumerate}

\begin{defbox}[AUC --- the area under the ROC curve]
Two classifiers may \emph{cross} in the ROC space: in the example from the slides, $C_1$ is better
for $\mathit{FPR}<0.43$ and $C_2$ is better for $\mathit{FPR}>0.43$ --- so there is no unambiguous
answer to \uv{which one is better}. Overall, the \emph{area under the ROC curve} (AUC) can be used:
if $\mathrm{AUC}(C_i)>\mathrm{AUC}(C_j)$, we say that $C_i$ is better than $C_j$. A perfect
classifier has $\mathrm{AUC}=1$, random guessing $\mathrm{AUC}=0.5$.
\end{defbox}

\subsection{Visualisation techniques}

\mimoq{the table as an overview; the box plot, dendrogram and ROC/lift curve themselves are in
the slides}
\begin{center}
\small
\begin{tabular}{@{}p{3.9cm}p{5.0cm}p{6.0cm}@{}}
\toprule
\textbf{Technique} & \textbf{Used for} & \textbf{Limitation / note} \\
\midrule
Histogram & the distribution of one numeric quantity & sensitive to the choice of bin width \\
Box plot & median, quartiles, outliers; comparison of groups & does not show bimodality (the
  violin plot does) \\
Scatter plot & the relation of two numeric quantities & with many points, overplotting $\to$
  transparency, hexbin \\
Scatter plot matrix (SPLOM) & all pairs of attributes at once & usable up to roughly 10
  attributes \\
Parallel coordinates & high-dimensional data --- an object is a polyline across vertical attribute
  axes & strongly dependent on the axis order; scaling required; crossing lines $=$ negative
  correlation \\
Heatmap & correlation matrix, confusion matrix & the choice of the colour scale (diverging for
  correlations) \\
Dendrogram & the result of hierarchical clustering & the leaf order is not unique \\
Projections (PCA, MDS, t-SNE, UMAP) & a 2-D view of high-dimensional data & t-SNE/UMAP distort
  global distances --- the \emph{sizes} of clusters and the gaps between them must not be
  interpreted \\
ROC / lift curve & the quality of a classifier & see Section~\ref{ssec:eval} \\
Mosaic plot & a contingency table of categorical attributes & a visual alternative to $\chi^2$ \\
Network visualisation & social networks, relation graphs & a \uv{hairball} for large graphs $\to$
  aggregate the communities \\
\bottomrule
\end{tabular}
\end{center}

\textbf{Three roles of visualisation in KDD:} \emph{(1)}~\emph{exploratory analysis} (the data
understanding phase) --- it reveals distributions, outliers, missing values, unexpected relations;
Anscombe's quartet shows four data sets with identical means, variances and correlations but
completely different structure, which only a plot reveals; \emph{(2)}~\emph{model visualisation} ---
a tree, a dendrogram, a network, decision boundaries; \emph{(3)}~\emph{presentation of the results}
to a manager --- aggregation, simplicity, one chart $=$ one message.

\textbf{Principles of good visualisation:} a high data-to-ink ratio (Tufte); do not distort (a zero
axis for bar charts, linear axes or clearly labelled logarithmic ones); choose the colour according
to the type of data (categories $\to$ a qualitative palette, order $\to$ sequential, deviation from
a centre $\to$ diverging) and with regard to colour blindness; always label and give units; avoid
3-D effects and pie charts with many slices.

\subsection{Evaluation of knowledge from the user's point of view}

The notion of \uv{evaluating the acquired knowledge} has three different levels that must be kept
apart:
\begin{itemize}
\item \emph{statistical evaluation of the model} --- accuracy, confusion matrix, ROC/AUC, error
      estimation by cross-validation (Section~\ref{ssec:eval});
\item \emph{evaluation of the interestingness of individual patterns} --- support, confidence, lift,
      $\chi^2$, t-weight and d-weight (Section~\ref{ssec:ar} and Chapter~\ref{sec:pravidla});
\item \emph{evaluation from the user's point of view} --- whether the knowledge found is
      comprehensible, valid, useful and novel.
\end{itemize}
The first two levels are objective and measurable, the third is necessarily subjective --- the same
pattern may be a discovery for one user and common knowledge for another. Part of the results may be
uninteresting or obvious from the user's point of view, part can be used directly and part has to be
presented more comprehensibly. It is useful to summarise the results in an \emph{analytical
report}; the output may also be the execution of a reasonable action.

\subsection{Summary for the examination}

\begin{itemize}
\item \textbf{Forms of knowledge representation}: IF-THEN rules (the most comprehensible, local),
      decision tree (global, readable up to depth 4--5), decision table, mathematical function
      (compact, not interpretable), prototypes/centroids, probabilistic model, graphs and networks,
      instances ($k$-NN, CBR). Be able to say, for each form, which method produces it and what its
      comprehensibility costs.
\item \textbf{Criteria for evaluating methods}: predictive accuracy, efficiency (time to construct
      and to use), robustness, scalability, interpretability, compactness of the model.
\item \textbf{Error estimation}: \emph{holdout} (large data; the training and test sets must not be
      mixed), \emph{$k$-fold cross-validation} (smaller data; 5 and 10 are common; accuracy as the
      average of $k$ accuracies $+$ a confidence interval with $t_{N,k-1}$),
      \emph{leave-one-out} (very small data, $m$-fold CV), \emph{validation set} (a third part of
      the data, on which the \emph{parameters} of the algorithm are tuned).
\item \textbf{Confusion matrix}: $TP$, $FN$, $FP$, $TN$; $p=\frac{TP}{TP+FP}$,
      $r=\frac{TP}{TP+FN}$, $F_1=\frac{2pr}{p+r}$ (the harmonic mean --- closer to the smaller of
      the two). Be able to explain \emph{why} accuracy is not enough for imbalanced data (1~\%
      intrusions $\Rightarrow$ 99~\% accuracy by \uv{doing nothing}).
\item \textbf{ROC}: $\mathit{TPR}=\frac{TP}{TP+FN}$ (sensitivity $=$ the recall of the positive
      class) against $\mathit{FPR}=\frac{FP}{FP+TN}$; $\mathit{TNR}$ $=$ specificity
      $=1-\mathit{FPR}$. The curve runs from $(0,0)$ to $(1,1)$, the diagonal $=$ random guessing.
      Be able to describe the construction from a ranking (partition at every case). \textbf{AUC}:
      1 for a perfect and 0.5 for a random classifier; needed because the ROC curves of two
      classifiers may \emph{cross}.
\item \textbf{Scoring and lift}: instead of a class, a probability estimate is assigned (for a tree,
      the confidence at the leaf), the examples are ranked and divided into bins; the lift curve is
      compared against the diagonal of random selection. Be able to carry out the economic
      reasoning (how many packages it pays to send).
\item \textbf{Three levels of \uv{evaluation}}: the statistical quality of the model, the
      interestingness of the patterns, and the user's evaluation. The first two are objective, the
      third subjective. \textbf{Four criteria of interesting knowledge}: comprehensibility,
      validity, usefulness, novelty.
\item \textbf{Visualisation}: the technique according to the task (histogram, box plot, scatter,
      SPLOM, parallel coordinates, heatmap, dendrogram, projections, ROC/lift); three roles in KDD
      (exploration --- Anscombe's quartet, model visualisation, presentation); the principles
      (Tufte, do not distort axes, colour by data type). Know the limitation: t-SNE and UMAP distort
      global distances.
\end{itemize}

%=====================================================================
\newpage
\section{Models for social network analysis, centrality and community detection}
\label{sec:sna}
%=====================================================================

The sentence of the topic names three things: \emph{models} for social network analysis,
\emph{centrality measures}, and \emph{community detection}. The chapter is structured accordingly:
first the representation of networks and the measures at the actor level (centrality) and the
network level (cohesion), then the models --- of the properties of real networks, of influence
spreading, of homophily, segregation and structural balance, plus link analysis (HITS, PageRank)
--- and finally community detection and link prediction.

\subsection{Delimitation of the field and motivation}

\begin{defbox}[Social network analysis]
\emph{Social network analysis} (SNA) is an interdisciplinary field with its origin in three areas:
\begin{itemize}
\item \emph{social science} --- questions and problems that explain social behaviour through an
      analysis of societies;
\item \emph{network analysis} --- the formulation and solution of problems that have a network
      structure;
\item \emph{graph theory} --- the mathematical apparatus for representing and analysing graphs.
\end{itemize}
A fundamental shift of perspective: SNA focuses on the \textbf{relationships} between social
entities (individuals, groups, institutions) \emph{rather than on the entities themselves} and their
attributes. Studying society from a network perspective considers individuals as \emph{embedded} in
a network of relations and seeks explanations for social behaviour in the structure of these
networks rather than in the individuals alone.
\end{defbox}

\textbf{Link analysis} as motivation: the objectives are to find relationships among the data and to
visualise the links. Application areas: telecommunications; criminalistics and law (groups of
criminals are mutually interlinked, and an analysis of these relationships can help uncover them);
marketing (relationships among customers).

\textbf{Six degrees of separation} (Milgram, 1967): an experiment on the small-world phenomenon.
Random people from Nebraska were asked to send a letter (via intermediaries) to a stockbroker in
Boston; they could only send it to someone they knew on a first-name basis. Among the letters that
found the target, the average number of links was \emph{six}. However, not many of the letters
arrived. The experiment presaged some later discoveries --- hubs, the tendency to go by geography or
profession. The \emph{Erdős numbers} confirm this: Paul Erdős wrote more than 1500 papers with 507
co-authors, and the average Erdős number among mathematicians is 4.65, the maximum 13.

\textbf{Scale-free (SF) networks} as a second motivation: some nodes have extremely many connections
(\emph{hubs}), most nodes have just a few. The consequence is \emph{resilience to a random
disturbance} together with \emph{vulnerability to coordinated attacks}: 80~\% of randomly chosen
nodes can fail without leading to the fragmentation of the cluster, but the elimination of 5--15~\%
of all the \emph{hubs} can cause the entire system to collapse. Applications: protection against
computer viruses spread through the internet (the eradication of internet viruses is essentially
impossible in an SF architecture), medicine (vaccination campaigns focused on hubs; new drugs
targeting the key molecules), business (cascaded financial crashes; marketing --- the spread of
\uv{contagion}, influencer marketing).

Examples of SF networks: \emph{social} (scientific collaboration --- researchers co-authoring
papers; Hollywood --- actors playing roles in the same movie), \emph{biological} (cell metabolism,
protein regulatory networks), \emph{socio-technical} (the internet --- routers and connections; the
World Wide Web --- web pages and URLs).

\subsection{Representing networks}

\begin{defbox}[Basic representations]
A network is a graph $G=(V,E)$: the \emph{vertices} (nodes) are actors, the \emph{edges} (links) are
the relations between them. Three equivalent notations:
\begin{itemize}
\item \textbf{a drawing of the graph};
\item \textbf{an edge list} --- pairs of vertices, plus a column of weights for weighted graphs;
\item \textbf{an adjacency matrix} $A=(a_{ij})$, where $a_{ij}=1$ if vertex $i$ is linked to vertex
      $j$ and 0 otherwise; for weighted graphs it contains the weights.
\end{itemize}
For \textbf{directed} graphs (who contacts whom) the adjacency matrix \emph{does not have to} be
symmetric; for \textbf{undirected} ones (who knows whom) it is symmetric. The same edge list has a
different interpretation in the two cases.
\end{defbox}

\textbf{Edge weights} can express: the frequency of interaction within the period of observation;
the number of items exchanged; individual perceptions of the strength of a relationship; costs in
communication or exchange (e.g.\ distance); combinations of the above. Edges can represent
interactions, flows of information or goods, similarities and affiliations, or social relations.

\textbf{Ego networks and \uv{whole} networks.} An \emph{ego network} is the personal network of one
actor (the \emph{ego}) and their neighbours (the \emph{alters}); the ego is \emph{not} needed for
further analysis, as all alters are by definition connected to the ego. A \uv{whole} network also
contains \emph{isolates}. \textbf{No studied network is \uv{whole} in practice} --- it is usually a
partial picture of the real-life network (the \emph{boundary specification problem}).

\subsection{Tie strength}

Social media connect users to each other extremely easily --- a user might have thousands of
\uv{friends} online --- but the connections are not of the same strength and play different roles in
community formation and information diffusion. \emph{Strong ties} are close friends, \emph{weak
ties} acquaintances. \textbf{The strength of weak ties} (Granovetter, 1973): occasional encounters
with distant acquaintances can provide important information about new opportunities (e.g.\ in a job
search).

\begin{defbox}[Homophily, transitivity, bridging]
\textbf{Homophily} --- the tendency to relate to people with similar characteristics (status,
beliefs, \dots). It leads to homogeneous groups (clusters) where forming relations is easier;
\emph{extreme} homogenisation, however, can act counter to innovation and idea generation, so
\emph{heterophily} is desirable in some contexts. Homophilous ties can be strong or weak.

\textbf{Transitivity} --- \uv{if there is a tie between $A$ and $B$ and a tie between $B$ and $C$,
then $A$ and $C$ will also be connected in a transitive network.} Strong ties are more often
transitive than weak ties, so transitivity \emph{indicates} the existence of strong ties (though it
is neither a necessary nor a sufficient condition). Transitivity and homophily together lead to the
formation of \emph{cliques} (fully connected clusters).

\textbf{Bridging} --- a \emph{bridge} is a node or an edge that connects different groups; usually
these are weak ties, but not every weak tie is a bridge. Bridges facilitate inter-group
communication, increase social cohesion and help spur innovation.
\end{defbox}

\textbf{Learning tie strength from the network topology.} Bridges connecting two different
communities are weak ties; an edge is a \emph{bridge} if its removal results in the disconnection of
its terminal nodes. Bridges are relatively rare in real-life networks $\Rightarrow$ the definition
is \emph{relaxed} to a \textbf{shortcut bridge}: one checks whether the distance between the two
terminal nodes \emph{increases} (to more than 2) if the edge is removed. The larger the distance, the
\emph{weaker} the tie. E.g.\ $d(2,5)=4$ if $(2,5)$ is removed, but $d(5,6)=2$ if $(5,6)$ is removed
$\Rightarrow$ $(5,6)$ is a stronger tie than $(2,5)$.

\begin{defbox}[Neighbourhood overlap]
For two nodes $i,j$ connected by a tie --- the larger the overlap, the \emph{stronger} the tie:
\[ O(i,j)=\frac{\text{number of shared friends of both } i \text{ and } j}
   {\text{number of friends adjacent to at least }i\text{ or }j}
   =\frac{|N_i\cap N_j|}{|N_i\cup N_j|-2}, \]
where the $-2$ in the denominator excludes $i$ and $j$ themselves.

\emph{Examples:} $O(2,5)=0$; $O(5,6)=\frac{|\{4\}|}{|\{2,4,5,6,7,10\}|-2}=\frac14$; for the graph
$G=(\{1,2\},\{(1,2)\})$ we get $O(1,2)=0$.
\end{defbox}

\textbf{Learning tie strength from profiles and interactions.} On Twitter (one can follow others
without their confirmation), the real friendship network is determined by the \emph{frequency} with
which two users talk to each other rather than by the follower--followee network --- and it is this
real friendship network that is more influential in driving Twitter usage. On Facebook, tie strength
can be predicted from various information (friend-initiated posts, messages exchanged in wall posts,
the number of mutual friends). Numeric link strength can be learned e.g.\ by maximum likelihood
estimation: user profile similarity determines the strength and the link strength in turn determines
user interaction $\Rightarrow$ maximise the likelihood based on the observed profiles and
interactions.

\subsection{Key players: centrality measures}

\begin{defbox}[Degree centrality]
The node's in- and out-degree $=$ the number of links that lead into / out of the node. In undirected
graphs the in- and out-degrees are identical.

It measures the node's \emph{connectedness}, \emph{influence} and/or \emph{popularity}; it helps
assess which nodes are central with respect to spreading information and influencing others in their
\uv{neighbourhood}.
\end{defbox}

\begin{defbox}[Paths and shortest paths]
A \emph{path} between two nodes is any sequence of non-repeating nodes that connects them. The
\emph{shortest path} connects the two nodes with the shortest number of edges (the \emph{distance}
between the nodes). Shorter paths are preferred for quick communication or exchange --- but not
always, for instance when spreading disease.

\textbf{Shortest path algorithms.} \emph{Single source:} general graphs (directed, negative edge
weights allowed) --- the Bellman--Ford algorithm, $O(|V||E|)$; graphs without negative edge weights
--- Dijkstra's algorithm, $O(|V|^2)$; to find the shortest path between only two vertices the
algorithm can stop once it is found. \emph{All pairs:} negative edge weights allowed --- the
Floyd--Warshall algorithm, $O(|V|^3)$; sparser graphs with negative weights --- Johnson's algorithm,
$O(|E||V|+|V|^2\log_2|V|)$; without negative weights --- Dijkstra, $O(|V|^3)$.
\end{defbox}

\begin{defbox}[Betweenness centrality]
The number of shortest paths passing through a node, divided by all shortest paths in the network:
for a given node $v$, calculate for each pair $s,t$ the number of shortest paths between $s$ and $t$
that pass through $v$, divide by all shortest paths between $s$ and $t$, and sum the computed values
over all node pairs; it is sometimes normalised so that the highest value is 1.

It indicates the nodes that are on \emph{communication paths} between other nodes, and helps
determine the points where the network might \emph{break apart}.
\end{defbox}

\begin{defbox}[Closeness centrality]
A measure of \emph{reach} --- the speed with which information can reach other nodes from a given
starting node: calculate the \emph{mean} length of all shortest paths from the node to all other
nodes in the network (how many hops it takes on average to reach every other node) and take the
reciprocal. Higher values indicate higher closeness (they are \uv{better}). Useful when the speed of
dissemination is of main concern.
\end{defbox}

\begin{defbox}[Eigenvector centrality]
It is proportional to the sum of the eigenvector centralities of all nodes directly connected to the
node considered: a node with a high eigenvector centrality is connected to other nodes with a high
eigenvector centrality. The metric is based on the idea that the power and status of an actor are
\emph{recursively} defined by the power and status of their neighbours; it measures how well an actor
is connected to other \emph{well-connected} actors.

Formally: for $G=(V,E)$ with adjacency matrix $A=(a_{i,j})$,
\[ x_i=\frac1\lambda\sum_{j\in M(i)}x_j=\frac1\lambda\sum_{j\in V}a_{i,j}x_j
   \quad\Longleftrightarrow\quad Ax=\lambda x, \]
where $M(i)$ is the set of neighbours of $i$. There are many eigenvalues $\lambda$ for which a
non-zero eigenvector solution exists; the additional \emph{requirement that all the entries be
non-negative} implies, however, that only the \emph{greatest} eigenvalue results in the desired
centrality measure. The eigenvector is defined only up to a common factor, so normalisation is
needed (e.g.\ the sum over all vertices being 1 or $|V|$).

It is closely related to Google's PageRank (links from highly linked-to pages count more,
Section~\ref{ssec:pagerank}).
\end{defbox}

\begin{center}
\small
\begin{tabular}{@{}p{2.7cm}p{12.6cm}@{}}
\toprule
\textbf{Measure} & \textbf{Interpretation in social networks} \\
\midrule
Degree & How many people can this person reach \emph{directly}? --- \emph{Network of music
  collaborations:} how many people has he or she collaborated with? \\
Betweenness & How likely is this person to be on the most direct route in the network? ---
  \emph{Network of spies:} through whom is most of the information likely to flow? \\
Closeness & How fast can this person reach everyone in the network? --- \emph{Network of sexual
  relations:} how fast will an STD spread to the rest of the network? \\
Eigenvector & How well is this person connected to other well-connected people? ---
  \emph{Network of citations:} who is the author cited most by other well-cited authors? \\
\bottomrule
\end{tabular}
\end{center}

\textbf{Sets of key players.} Individual measures are not enough. In the example from the slides,
node 10 is the most central according to degree centrality, but nodes 3 and 5 \emph{together} reach
more nodes; moreover, the tie between 3 and 5 is \emph{critical} --- if severed, the network breaks
into two isolated sub-networks. Other things being equal, nodes 3 and 5 \emph{together} are
therefore more important to this network than node 10 alone. \emph{Thinking about sets of key players
can help.}

\begin{defbox}[Structural equivalence]
It expresses the \emph{similarity} between actors in a social network; it is based on the neighbours
they share (equivalently, on the number of identical ties they have). Two actors are
\emph{structurally equivalent} if their positions can be swapped \emph{without modifying the overall
structure} of the network. Approximate structural equivalence can be used for hierarchical
clustering of social networks.
\end{defbox}

\subsection{Cohesion: measures at the network level}

\begin{defbox}[Reciprocity, density, clustering coefficient]
\textbf{Reciprocity} --- the ratio of \emph{reciprocated} relations (an edge in both
directions) over all relations in the network; \textbf{it makes sense only in directed
graphs}.

\textbf{Density} --- the ratio of the number of edges over the number of \emph{possible} edges:
\[ \text{density}=\frac{|E|}{|V|(|V|-1)/2}\quad\text{(undirected graph)},\qquad
   \frac{|E|}{|V|(|V|-1)}\quad\text{(directed)}. \]
A \emph{clique} (a perfectly connected network) has density 1; directed graphs have half the
density, as there are twice as many possible edges.

\textbf{Clustering coefficient} of node $i$ --- the density of its \emph{neighbourhood}
$N_i=\{v_j: e_{ij}\in E\ \vee\ e_{ji}\in E\}$, $k_i=|N_i|$:
\[ C_i=\frac{2\,|\{e_{jk}: v_j,v_k\in N_i,\ e_{jk}\in E\}|}{k_i(k_i-1)}
   \ \text{(undirected)},\qquad
   C_i=\frac{|\{e_{jk}\}|}{k_i(k_i-1)}\ \text{(directed)}. \]
The network coefficient $C=\frac{1}{|V|}\sum_i C_i$ indicates the presence of
(sub)communities.
\end{defbox}

\begin{defbox}[Diameter, average distance, eccentricity, small worlds]
The network's \textbf{diameter} $=$ the longest shortest path between two nodes; the
\textbf{average distance} $=$ the average of all shortest paths; the \textbf{eccentricity}
of a vertex $=$ its greatest geodesic distance from another vertex.

A \textbf{small world} --- a network with a \emph{significantly high clustering coefficient}
and a \emph{short average path length}: clusters arise from the transitivity of strong ties,
short paths are provided by weak ties (bridges) between the clusters. Very common in social
networks.
\end{defbox}

\subsection{A model of real-world complex networks and their properties}

\begin{defbox}[Six characteristic properties of real-world complex networks]
Real-world complex networks are \emph{non-random and non-regular} graphs with unique features.
Examples: social networks, information or knowledge networks, technological networks, biological
networks. The properties:
\begin{enumerate}
\item \textbf{The small-world effect} --- Stanley Milgram was the first to point out the existence of
      small-world effects in real social networks (the experiment of the 1960s, about 300
      participants; the median path length of the successful paths was 6). It rests on the concept of
      \uv{six degrees of separation}: everyone is, on average, six steps away from any other person.
      There are important consequences for the speed at which information and diseases spread.
\item \textbf{Transitivity or clustering} --- in real-world, especially social networks there is a
      high probability of finding \emph{complete sub-networks} within larger networks (everyone knows
      everyone).
\item \textbf{Power-law degree distributions} --- the \emph{degree distribution} is the probability
      distribution of the degrees of nodes over the whole network (a histogram of the fraction of
      nodes with degree $k$). Real networks have degree distributions quite different from those of
      random graphs.
\item \textbf{Network resilience} --- the impact on the connectivity of the network when one or more
      nodes are removed. Most networks are \emph{robust} against random vertex removal, but
      considerably less robust to the \emph{targeted} removal of the highest-degree vertices. Also,
      when \emph{gatekeepers} are deleted, there are significant changes in the network's
      communication ability, since some nodes become disconnected. The removal of a single node is
      usually not cause for alarm --- it is more appropriate to test the resilience by removing a
      certain \emph{percentage} of nodes.
\item \textbf{Mixing patterns} --- in networks where different types of nodes coexist, it is common to
      observe a certain \emph{selectivity} in the establishment of connections. Selective linking is
      called \emph{assortative mixing}, or homophily; a classic example is mixing by race. Real
      networks show higher tendencies for assortative mixing.
\item \textbf{Community structure} --- most real social networks show community structure, which
      arises as a consequence of both \emph{global and local heterogeneity} of the edges'
      distribution in the graph: we often find high concentrations of edges within some areas of the
      graph (\emph{communities}) and low concentrations between those regions.
\end{enumerate}
\end{defbox}

\begin{defbox}[Scale-free networks and preferential attachment]
\textbf{Scale-free networks} (Barabási, Albert, 1999) have a \emph{power-law} degree
distribution (the WWW, citation and biological networks): a few highly connected \emph{hubs}
and a large number of low-degree nodes. They arise from two mechanisms: \textbf{growth} $+$
\textbf{preferential attachment} --- new nodes tend to connect to already well-connected ones
(\emph{rich-get-richer}, the \emph{Matthew effect}, Price's \emph{cumulative advantage}). The
reason may be \emph{popularity} (the rich get richer), \emph{quality} (the good get better),
or a mixed model --- among similar nodes, those that reach critical mass first become stars
(the \emph{halo effect}). The resulting network has a long-tailed degree distribution and a
small-world structure.
\end{defbox}

\begin{defbox}[Core-periphery structures and centralisation]
\textbf{Centralisation} --- from the differences in node degrees (normalised to $[0,1]$) it
measures how much the network depends on a few hubs: centralised networks are better at
coordination but \emph{more prone to failure} if the key players disconnect. Many large
communities have a densely interconnected \emph{core} critical for connecting a much larger
\emph{periphery}; the cores are found visually, by the location of high-degree nodes, or by
\emph{bow-tie analysis} (the structure of the Web).
\end{defbox}

\subsection{Influence modelling}

\begin{defbox}[Common properties of influence modelling methods]
\begin{itemize}
\item the social network is represented as a \emph{directed} graph, with each actor being one node;
\item each node starts as \emph{active} or \emph{inactive};
\item a node, once activated, will activate its neighbouring nodes;
\item \textbf{once a node is activated, it cannot be deactivated.}
\end{itemize}
\end{defbox}

\begin{defbox}[The Linear Threshold Model (LTM)]
An actor (node) would take an action if the number of their friends who have taken the action exceeds
a certain \emph{threshold}.
\begin{itemize}
\item Each node $v$ chooses a threshold $\theta_v$ \emph{randomly} from a uniform distribution on
      $[0,1]$.
\item In each discrete step, all nodes that were active in the previous step remain active.
\item The nodes satisfying
      \[ \sum_{u\in N_v^{\text{active}}} w_{u,v}\ \ge\ \theta_v
         \sum_{u\in N_v} w_{u,v} \]
      will also be activated.
\end{itemize}
LTM focuses on the \emph{receivers} of the influence.
\end{defbox}

\begin{defbox}[The Independent Cascade Model (ICM)]
It focuses on the \emph{senders} rather than the receivers.
\begin{itemize}
\item A node $v$, once activated at step $t$, has \emph{one} chance to activate each of its
      neighbours, randomly.
\item For a neighbouring node $u$ the activation succeeds with probability $p_{v,u}$ (e.g.\
      $p=0.5$).
\item If the activation succeeds, $u$ becomes active at step $t+1$.
\item In the subsequent rounds, $v$ will \emph{not} attempt to activate $u$ any more.
\item The diffusion process starts with an initial activated set of nodes and continues until no
      further activation is possible.
\end{itemize}
\end{defbox}

\begin{defbox}[Influence maximisation]
\textbf{The task:} given a network and a parameter $k$, which $k$ nodes should be selected to be in
the activation set $B$ in order to maximise the influence in terms of active nodes at the end? Let
$\sigma(B)$ denote the expected number of nodes influenced by $B$:
\[ \max_{|B|\le k}\ \sigma(B). \]
The problem is \textbf{NP-hard}, but it can be proved that the \emph{greedy} approach yields
solutions that are \textbf{63~\%} of the optimal.

\textbf{The greedy approach:} start with $B=\emptyset$; evaluate $\sigma(\{v\})$ for each node and
pick the node with maximum $\sigma$ as the first element of $B$; then repeatedly select the node
that will increase $\sigma(B)$ the most, i.e.\ $\arg\max_{v\in V\setminus B}\sigma(B\cup\{v\})$.
\end{defbox}

\subsection{Influence versus correlation}

User attributes and behaviours quite often tend to correlate with their social networks. Suppose a
binary attribute is associated with each node (e.g.\ whether or not the person is a smoker): if the
attribute is correlated with the network, the actors are expected to share the same attribute values,
positively correlated with the social connections (smokers are more likely to interact with other
smokers).

\begin{defbox}[Test for correlation]
If the fraction of edges linking nodes with \emph{different} attribute values occurs significantly
\emph{less} often than the expected probability, then there is evidence of correlation.

\emph{Example:} if connections are independent of the smoking behaviour, the fraction of smokers is
$p=4/9$ and of non-smokers $1-p=5/9$. One edge connects two smokers with probability $p\cdot p$, two
non-smokers with $(1-p)(1-p)$, and a smoker and a non-smoker with
$2p(1-p)=2\cdot\frac49\cdot\frac59=49\,\%$. In the actual graph this fraction is
$2/14=14\,\%<49\,\%$, so the network \textbf{demonstrates some degree of correlation} with respect to
the smoking behaviour.
\end{defbox}

\begin{defbox}[Three social processes explaining correlation]
\begin{itemize}
\item \textbf{Homophily} --- the tendency to link to others who share a certain similarity with us
      (\uv{birds of a feather flock together}).
\item \textbf{Confounding} --- the correlation can also be forged by \emph{external} influences from
      the environment (\uv{individuals living in the same city are more likely to become friends than
      random individuals}).
\item \textbf{Influence} --- a process causing behavioural correlations between adjacent actors
      (\uv{if most of one's friends switch to another mobile company, one might be influenced and
      switch as well}).
\end{itemize}
\textbf{Why distinguish influence from homophily and confounding?} \emph{Analysis:} to predict the
dynamics of the system --- whether a new norm of behaviour, technology or idea can diffuse like an
epidemic. \emph{Design:} to design a system for inducing a particular behaviour, e.g.\ vaccination
strategies (random, targeting a demographic group, random acquaintances) or viral marketing campaigns.
\end{defbox}

\begin{defbox}[The logistic model of influence and the shuffle test]
In studies of influence modelling, influence is determined by \emph{timestamps}: for an edge $(u,v)$,
$u$ can influence $v$ at discrete time steps $t=0,1,2,\dots$; for each $t$, every inactive node
becomes active with probability $p(a)$, where $a$ is the number of active contacts. The probability of
activation is modelled as a \emph{logistic} function of the number of active friends:
\[ p(a)=\frac{e^{\alpha\ln(a+1)+\beta}}{1+e^{\alpha\ln(a+1)+\beta}},\qquad\text{equivalently}\quad
   \ln\frac{p(a)}{1-p(a)}=\alpha\ln(a+1)+\beta, \]
where $a$ is the number of active friends, $\alpha$ the \emph{social correlation coefficient},
$\beta$ a constant explaining the innate bias for activation, and $\ln(a+1)$ the explanatory
variable.

\textbf{Activation likelihood.} Suppose at a time point $t$, $y_{a,t}$ users with $a$ active friends
become active and $n_{a,t}$ users who also have $a$ active friends stay inactive. The likelihood of
activation at time $t$ is $\prod_a p(a)^{y_{a,t}}\big(1-p(a)\big)^{n_{a,t}}$ and for the whole
sequence $\prod_t\prod_a p(a)^{y_{a,t}}\big(1-p(a)\big)^{n_{a,t}}$. Given the user activity log we
can compute the $\alpha$ and $\beta$ that maximise this likelihood; \textbf{$\alpha$ measures the
correlation between nodes}.

\textbf{The shuffle test (a test for influence).} The key idea: \emph{if influence does not play a
role, the timing of activation should be independent of the timing of other actors} --- even if we
randomly shuffle the timestamps of the user activities, we should obtain a similar social
correlation. \emph{The test:} shuffle the timestamps of the user activities; if the new estimate of
the social correlation is \emph{significantly different} from the estimate based on the activity log,
then there is evidence of \textbf{influence}.

The procedure: \emph{(1)}~for all $a$ compute $y_{a,t}$ and $n_{a,t}$; \emph{(2)}~compute the
$\alpha,\beta$ that maximise the likelihood; \emph{(3)}~pick a random permutation $\pi$ of the
activation times and set the activation time of node $v_i$ to $t'_i=t_{\pi(i)}$; from the new times
compute $y'_{a,t}$, $n'_{a,t}$; \emph{(4)}~compute the $\alpha',\beta'$ that maximise the new
likelihood; \emph{(5)}~the model exhibits \emph{no} social influence if $\alpha$ and $\alpha'$ are
close to each other.
\end{defbox}

\subsection{Homophily, selection and social influence}

\begin{defbox}[The homophily test]
The motivation: \uv{birds of a feather flock together} --- your friends are more similar to you in
age, race, interests and opinions than a random collection of individuals. \emph{Homophily} is the
principle that we tend to be similar to our friends.

\textbf{The test:} if each node were \emph{randomly} assigned an attribute according to the balance
in the real network, the number of \uv{heterogeneous} edges should not change significantly relative
to the real network. If a fraction $p$ are males and $q$ are females, then the probability of a
male--male edge is $p^2$, of a female--female edge $q^2$, and of a \emph{cross-gender} edge $2pq$.

\textbf{If the fraction of heterogeneous (cross-gender) edges is significantly less than $2pq$, then
there is evidence for homophily.}

\emph{Example:} there are 5 cross-gender edges out of 18; $p=6/9=2/3$, $q=3/9=1/3$; without
homophily there should be $2pq=4/9$, i.e.\ 8 out of 18 $\Rightarrow$ \emph{evidence of homophily}.
\end{defbox}

\textbf{Notes on the test.} Significantly \emph{more} than $2pq$ heterogeneous edges $=$
\emph{inverse homophily} (male--female dating relationships). The test works for any
characteristic and also for more than two values --- the number of heterogeneous edges is
compared to a randomly generated graph using the real proportions as probabilities.

\begin{defbox}[Selection vs.\ social influence]
\textbf{Triadic closure}: when two individuals share a common friend, a friendship between
them is more likely --- the link may form even if neither is aware of the mutual friend.
The formation of a link can rarely be attributed to a single factor.

\textbf{Selection} --- people choose friends who are like themselves; it operates actively
(choosing similar friends in a small group) as well as through the environment (a homogeneous
school population).

\textbf{Mutable vs.\ immutable characteristics:} \emph{immutable} ones do not change (gender,
race) or change consistently with the population (age); \emph{mutable} ones change over time
(behaviours, interests, opinions).

\textbf{Social influence} is the \emph{opposite} of selection: selection --- characteristics
drive the formation of links; influence --- existing links shape mutable characteristics.
\end{defbox}

\textbf{Longitudinal studies:} from a single snapshot of a network, selection and influence
cannot be distinguished --- links and behaviour must be tracked over time (\emph{does
behaviour change after changes to the network, or the network after changes in behaviour?}).
The answer matters for designing interventions: an anti-drug programme \uv{friends persuade
friends} only works if the homophily arose by \emph{influence}, not by selection.
\emph{Christakis and Fowler} (obesity, 12\,000 people, 32 years) tested \emph{selection},
\emph{confounding} and \emph{social influence} --- and \textbf{found significant evidence for
influence}: obesity is a \uv{contagion} spreading through social influence.

\subsection{Affiliation networks and closure processes}

\begin{defbox}[Affiliation and social-affiliation network]
An \textbf{affiliation network} is a \emph{bipartite graph} of persons and \emph{affiliations}
(\emph{foci}: activities, organisations, \dots). A \textbf{social-affiliation network} has
\emph{two types of edges}: \emph{person--person} (a social relationship) and
\emph{person--focus} (participation).
\end{defbox}

\begin{defbox}[Three closure processes]
\begin{itemize}
\item \textbf{Triadic closure} --- two people connect because they have a common friend.
\item \textbf{Focal closure} --- a shared focus connects them; closure due to
      \emph{selection}.
\item \textbf{Membership closure} --- a person joins a focus where their friend already is;
      closure due to \emph{social influence}.
\end{itemize}
\end{defbox}

\textbf{Measurement:} from two snapshots of the network ($t_1$, $t_2$), for each $k$ define
$T(k)$ $=$ the fraction of pairs that had exactly $k$ common friends (resp.\ foci) and no edge
at $t_1$ and formed an edge by $t_2$. The baseline model of independent opportunities:
$T_{\text{base}}(k)=1-(1-p)^{k}$.

\textbf{Empirically} (the emails of 22\,000 students --- Kossinets and Watts 2006;
LiveJournal --- Backstrom et al.; Wikipedia --- Crandall et al.) all three closures show the
same pattern: a \emph{significant increase from 1 to 2} common friends/foci and
\emph{diminishing returns} afterwards; sharing a class has nearly the same effect as sharing
a friend. On Wikipedia, moreover, the similarity of editors (the overlap of edited articles)
was measured around the first interaction: the similarity grows \emph{already before} the
link is formed (\emph{selection}) and continues to grow \emph{after} it (\emph{influence}).

\subsection{The Schelling model of segregation}

\begin{defbox}[The Schelling model (T.~Schelling, 1970s)]
A simple model that shows how \emph{global} patterns of self-chosen segregation might occur because of
\emph{local} homophily.
\begin{itemize}
\item The population is formed by two types of individuals --- agents $X$ and $O$; the two types
      represent an \emph{immutable} characteristic (race, ethnicity, country of origin, native
      language).
\item Agents are placed randomly in a grid.
\item An agent is \emph{satisfied} with their location if at least $t$ of their neighbours are agents
      of the same type ($t$ being a pre-set threshold).
\item In each round, first identify all \emph{unsatisfied} agents; then all unsatisfied agents move to
      a location where they are satisfied.
\end{itemize}
Sometimes an agent cannot find a new location that satisfies --- then they are left alone or move to a
random location. The moves may cause a \emph{previously satisfied agent to no longer be satisfied} ---
which is why the process continues.

\textbf{The result: the grid is more segregated than before.}

\textbf{Model variations:} schedule agents to move in random order or in a sweeping motion; agents
move to the nearest satisfying location or to a random one; the threshold could be a percentage or
vary among or within population groups; the world could wrap (the left edge meets up with the right
edge). There are many other variations, but \emph{the results usually end up the same:
self-imposed segregation.}

\textbf{Conclusion:} the Schelling model is an example of how \emph{immutable} characteristics can
become highly correlated with \emph{mutable} ones: the choice of where to live (mutable) over time
conforms to the agents' type (immutable), producing segregation (homophily).
\end{defbox}

\subsection{Positive and negative relationships: structural balance}

Until now a relation in a network has had a positive connotation; there exist, however, social
networks where negative (antagonistic, enemy) relationships can also be expressed. \emph{Positive}
relations $=$ friendly, \emph{negative} $=$ enemy. The question: \textbf{when is the network in a
balanced state?} We shall see how local properties enforce global ones (cf.\ the Schelling model). A
simplification: we consider networks where \emph{all} pairs of nodes are connected by a positive or
negative relation (i.e.\ a complete graph, a clique) --- applicable to small groups of people or to
countries; we relax it later.

\begin{defbox}[The structural balance property]
For three nodes there are (up to symmetry) four sign configurations:
\begin{center}
\small
\begin{tabular}{@{}p{3.0cm}p{9.5cm}p{2.4cm}@{}}
\toprule
\textbf{Signs} & \textbf{Interpretation} & \textbf{State} \\
\midrule
$+,+,+$ & $A$, $B$, $C$ are mutual friends & \emph{balanced} \\
$+,-,-$ & $A$ and $B$ are friends with a mutual enemy $C$ & \emph{balanced} \\
$+,+,-$ & $B$ and $C$ are friends of $A$, but they don't get along with each other & not balanced \\
$-,-,-$ & $A$, $B$, $C$ are mutual enemies & not balanced \\
\bottomrule
\end{tabular}
\end{center}
\textbf{The structural balance property:} for every triangle, all three edges are positive, or
\emph{exactly one} is positive.
\end{defbox}

\begin{defbox}[The Balance Theorem]
\textbf{Theorem.} If a labelled \emph{complete} graph is balanced, then either
\begin{itemize}
\item all pairs of nodes are friends, \emph{or}
\item the nodes can be divided into two groups $X$ and $Y$ such that every pair of nodes in $X$ like
      each other, every pair in $Y$ like each other, and everyone in $X$ is the enemy of everyone in
      $Y$.
\end{itemize}

\textbf{Proof.} Let $G$ be a balanced network. If all edges are positive, we are done. If there is at
least one negative edge, pick a node $A$ and put $X=$ all friends of $A$ (including $A$), $Y=$ all
enemies of $A$. Then: \emph{(1)}~every two nodes in $X$ are friends --- otherwise we would have a
triangle consisting of $A$ and two of its friends with a single negative edge, i.e.\ two positive and
one negative, which is not balanced; \emph{(2)}~every two nodes in $Y$ are friends --- otherwise we
would have a triangle of $A$ and two of its enemies with three negative edges, which is not balanced;
\emph{(3)}~every node in $X$ is an enemy of every node in $Y$ --- otherwise the triangle of $A$, a
friend from $X$ and an enemy from $Y$ would have two positive and one negative edge. $\square$
\end{defbox}

\textbf{Applications of structural balance.} \emph{The evolution of alliances in Europe, 1872--1907}
(solid dark edges indicate friendship, dotted red edges enmity): the Three Emperors' League
1872--81, the Triple Alliance 1882, the lapse of the German--Russian treaty 1890, the
French--Russian alliance 1891--94, the Entente Cordiale 1904, the British--Russian alliance 1907 ---
the network gradually \emph{slides into a balanced labelling}, and into the First World War.
\emph{The conflict over Bangladesh's separation from Pakistan (1972):} the United States's somewhat
surprising support of Pakistan becomes less surprising when one considers that the USSR was China's
enemy, China was India's foe, and India had traditionally bad relations with Pakistan --- since the
U.S.\ was at that time improving its relations with China, it supported the enemies of China's
enemies.

\begin{defbox}[Weak structural balance and generalisations]
\textbf{The weak structural balance property:} there is no set of three nodes whose edges consist of
\emph{exactly two positive edges and one negative edge}. (The \uv{three negative} configuration is
therefore allowed --- it corresponds to a pair making an alliance against the third one; \uv{two
positive and one negative} is a stronger imbalance, because $B$ and $C$ try to resolve their
animosity.)

\textbf{Theorem.} If a labelled complete graph is weakly balanced, then its nodes can be divided into
\emph{groups of mutual friends} such that every two nodes belonging to different groups are enemies.

\textbf{Proof.} \emph{(1)}~Let $A$ be a node and $X$ the set of friends of $A$, including $A$.
\emph{(2)}~All $A$'s friends are friends with each other. \emph{(3)}~$A$ and all its friends are
enemies with everyone else. \emph{(4)}~Remove $X$ from the network and continue with step 1.
$\square$

\textbf{Generalisation to non-complete graphs.} \emph{(a)}~Not all pairs in a network are connected
--- a (non-complete) graph is \emph{balanced} if it can be \uv{completed} by adding edges to form a
signed complete graph that is balanced. \emph{(b)}~Not all triangles must be balanced --- if
\emph{most} triangles are balanced, then the network can be \emph{approximately} divided into two
fractions.
\end{defbox}

\subsection{Link analysis: hubs, authorities and the HITS algorithm}

The year 1998 was an eventful one for web link analysis models --- \emph{both} the PageRank and the
HITS algorithms were reported in it. The connections between them are quite striking; since then
PageRank has emerged as the dominant model, due to its \emph{query-independence}, its ability to
combat spamming, and Google's business success.

\begin{defbox}[Authorities and hubs]
An \textbf{authority} --- roughly, a page with many in-links. The idea: the page may have good or
authoritative content on some topic and thus many people trust it and link to it.

A \textbf{hub} --- a page with many \emph{out-links}. It serves as an organiser of the information on
a particular topic and points to many good authority pages.

\textbf{The key idea of HITS:} \emph{a good hub points to many good authorities} and \emph{a good
authority is pointed to by many good hubs.} Authorities and hubs thus have a \emph{mutual
reinforcement} relationship (a densely linked bipartite sub-graph of hubs and authorities).

\textbf{Why separate hubs from authorities?} Competing firms will not link to each other and cannot be
viewed as directly endorsing each other. In PageRank, by contrast, endorsement passes directly from
one prominent page to another (a page is important if it is cited by other important pages) --- this
is the dominant mode of endorsement in academic and governmental pages, among bloggers, in personal
pages, and in the scientific literature.
\end{defbox}

\begin{defbox}[The HITS algorithm (Hypertext Induced Topic Search)]
Unlike PageRank, which is a \emph{static} ranking algorithm, HITS is search-query \emph{dependent}.
When the user issues a query, HITS first expands the list of relevant pages returned by a search
engine and then produces \emph{two} rankings of the expanded set: an authority ranking and a hub
ranking.

\textbf{Grabbing pages.} Given a broad query $q$: send $q$ to a search engine; collect the $t$
highest-ranked pages ($t=200$ in the HITS paper) --- this set is called the \emph{root set} $S$; grow
$S$ by including any page pointed to by a page in $S$ and any page that points to a page in $S$ ---
this gives the larger \emph{base set} $B$.

\textbf{Scores.} Let $G=(V,E)$ be the hyperlink graph of $B$ with $n$ pages and adjacency matrix $L$
($L_{ij}=1$ if there is an edge $(i,j)$, 0 otherwise). Let $a(i)$ be the authority score and $h(i)$
the hub score of page $i$. The mutual reinforcement:
\[ a(i)=\sum_{(j,i)\in E}h(j),\qquad h(i)=\sum_{(i,j)\in E}a(j), \]
in matrix form, with $a=(a(1),\dots,a(n))^{\!\top}$ and $h=(h(1),\dots,h(n))^{\!\top}$:
\[ a=L^{\!\top}h,\qquad h=La. \]

\textbf{Computation by power iteration:}
\begin{algorithmic}[1]
\State $a_0\gets h_0\gets (1,1,\dots,1)^{\!\top}$; $k\gets 1$
\Repeat
  \State $a_k\gets L^{\!\top}h_{k-1}$; \quad $h_k\gets La_k$
  \State $a_k\gets a_k/\|a_k\|_1$; \quad $h_k\gets h_k/\|h_k\|_1$ \Comment{normalisation}
  \State $k\gets k+1$
\Until{$\|a_k-a_{k-1}\|_1<\varepsilon$ \textbf{and} $\|h_k-h_{k-1}\|_1<\varepsilon$}
\State \Return $a$ and $h$
\end{algorithmic}
\textbf{Normalisation is indispensable}, because the magnitude of the hub and authority values tends
to grow with each update; convergence occurs only when normalisation is taken into account (the
scores are divided by the sum of all scores).
\end{defbox}

\begin{defbox}[Spectral analysis of HITS]
After $k$ steps (the \emph{spectrum} of a matrix is the set of its eigenvalues),
\[ a_k=\big(L^{\!\top}L\big)^{k-1}a_0,\qquad h_k=\big(LL^{\!\top}\big)^{k}h_0. \]
The vector $a_k$ therefore has to converge to an \emph{eigenvector} of the matrix $L^{\!\top}L$:
there are normalising constants $c_k$ such that $a_k/c_k$ converges to a limit, and the direction
does not change when multiplied by the matrix exactly when the vector is an eigenvector.

The matrix $L^{\!\top}L$ is \emph{symmetric}, so it has $n$ eigenvectors $z_1,\dots,z_n$ with
eigenvalues $c_1,\dots,c_n$, $|c_1|\ge|c_2|\ge\cdots\ge|c_n|$. Expanding
$a_0=x_1z_1+\cdots+x_nz_n$, we get
\[ a_k=c_1^{k}x_1z_1+\cdots+c_n^{k}x_nz_n,\qquad
   \frac{a_k}{c_1^{k}}=x_1z_1+x_2\Big(\frac{c_2}{c_1}\Big)^{\!k}z_2+\cdots \]
Assuming $c_1>c_2\ge\cdots$, every term except the first goes to 0 as $k\to\infty$, hence
$a_k/c_1^k$ converges to $x_1z_1$ --- to the \emph{principal} eigenvector. (Convergence regardless of
the initial vector and the relaxation of the assumption $c_1>c_2$ are not shown here.)
\end{defbox}

\textbf{Strengths and weaknesses of HITS.} \emph{Strength:} its ability to rank pages according to the
\emph{query topic}, which may provide more relevant authority and hub pages. \emph{Weaknesses:}
\emph{(1)}~it can be easily spammed --- adding out-links in one's own page is very easy;
\emph{(2)}~\emph{topic drift} --- many pages in the expanded set may not be on topic;
\emph{(3)}~inefficiency at query time --- collecting the root set, expanding it and performing the
eigenvector computation are all expensive operations.

\subsection{PageRank}
\label{ssec:pagerank}

\begin{defbox}[Rank prestige]
In prestige measures an important factor is considered: \emph{the prominence of the individual actors
who do the \uv{voting}}. In the real world, a person chosen by an important person is more prestigious
than one chosen by a less important person (a CEO's vote counts for much more than a worker's). If
one's circle of influence is full of prestigious actors, then one's own prestige is also high.

The rank prestige $P(i)$ is defined as a linear combination of the ranks of those who point to $i$:
\[ P(i)=A_{1i}P(1)+A_{2i}P(2)+\cdots+A_{ni}P(n),\qquad\text{in matrix form}\quad P=A^{\!\top}P, \]
where $A_{ji}=1$ if $j$ points to $i$ and 0 otherwise. This is the \emph{characteristic equation} of
the eigensystem of $A^{\!\top}$; $P$ is an eigenvector of $A^{\!\top}$.
\end{defbox}

\begin{defbox}[Basic PageRank]
PageRank interprets a hyperlink from page $u$ to page $v$ as a \emph{vote}, by page $u$, for page $v$.
It looks at more than the sheer number of votes, though --- it also analyses the page that casts the
vote: \emph{votes cast by important pages weigh more heavily and help make other pages more
important}. Since a page may point to many others, its prestige score has to be \emph{shared}.

The Web as a directed graph $G=(V,E)$, $n=|V|$. The PageRank of page $i$:
\[ P(i)=\sum_{(j,i)\in E}\frac{P(j)}{O_j}, \]
where $O_j$ is the number of out-links of $j$. In matrix form: let $A_{ij}=1/O_i$ for $(i,j)\in E$ and
0 otherwise; then $P=A^{\!\top}P$.

\textbf{The intuition:} PageRank is a kind of \uv{fluid} that circulates through the network, passing
from node to node across edges and pooling at the most important nodes. \emph{The total PageRank in
the network remains constant} --- each page takes its PageRank, divides it up and passes it along its
links; PageRank is never created nor destroyed, just moved around. Therefore \emph{no normalisation is
needed} --- unlike HITS.
\end{defbox}

\textbf{The problem with the basic definition.} In many networks the \uv{wrong} nodes end up with all
the PageRank. If $F$ and $G$ point to each other rather than to $A$, the PageRank that flows to them
from $C$ can never circulate back --- for large $k$ we get $\frac12$ for each of $F$ and $G$ and 0 for
all the others. This is a \uv{slow leak}: small sets of nodes that can be reached from the rest of the
graph but have no paths back. \textbf{The solution:} send a small fraction of PageRank to \emph{all}
nodes.

\begin{defbox}[Scaled PageRank and the damping factor]
Pick a \emph{scaling factor} $d$ between 0 and 1 and replace the basic update rule:
\begin{enumerate}
\item first apply the basic PageRank update rule;
\item scale down all PageRank values by a factor of $d$ (the total PageRank in the network has shrunk
      from 1 to $d$);
\item divide the residual $1-d$ units of PageRank \emph{equally} over all nodes, giving $(1-d)/n$ to
      each.
\end{enumerate}
\textbf{Why it works:} it is a \uv{water cycle} that evaporates $1-d$ units of PageRank in each step
and rains it down uniformly across all nodes.

Repeated application of the scaled rule \emph{converges} to a set of limiting PageRank values and the
equilibrium is \emph{unique} --- but the values depend on the choice of $d$. In practice $d$ is
between 0.8 and 0.9 ($d=0.85$ in the PageRank paper); $d$ is called the \emph{damping factor}.
\end{defbox}

\begin{defbox}[PageRank as a Markov chain]
The Web is modelled as a \emph{stochastic process}: each page is a \emph{state}, a hyperlink is a
\emph{transition} with a certain probability. \emph{Random surfing}: each transition probability is
$1/O_i$ if we assume the surfer clicks the hyperlinks uniformly at random (the \uv{back} button is not
used and no URL is typed in).

Let $A$ be the state transition probability matrix and $p_0$ the initial probability distribution.
Then $p_1=A^{\!\top}p_0$ and in general $p_k=A^{\!\top}p_{k-1}$. \emph{A theorem of Markov chains:} a
finite Markov chain defined by a stochastic matrix $A$ has a \textbf{unique stationary probability
distribution} if $A$ is \emph{irreducible} and \emph{aperiodic}. A stationary distribution means that
$p_k$ converges to a steady-state vector $\pi$ regardless of the choice of $p_0$. At the steady state
$\pi=A^{\!\top}\pi$, so $\pi$ is the principal eigenvector of $A^{\!\top}$ with eigenvalue 1 --- and
it is $\pi$ that is used as the PageRank vector. This is reasonable and intuitive: it reflects the
\emph{long-run probabilities} that a random surfer will visit the pages, and a page has high prestige
if the probability of visiting it is high.

\textbf{Neither condition, however, holds for the web graph:}
\begin{itemize}
\item \emph{$A$ is not stochastic} --- many pages have no out-links, which is reflected in the matrix
      by rows of complete zeros (\emph{dangling pages}). Solutions: \emph{(1)}~remove such pages
      during the computation (they do not directly affect the ranking of any other page);
      \emph{(2)}~add a complete set of outgoing links from each such page to all pages, i.e.\
      $A_{ij}=1/n$.
\item \emph{$A$ is not irreducible} --- irreducibility means that the web graph is \emph{strongly
      connected} (for each pair $i,j$ there is a path from $i$ to $j$); a general web graph is not.
\item \emph{$A$ is not aperiodic} --- a state $i$ is \emph{periodic} with period $k>1$ if $k$ is the
      smallest number such that all paths leading from $i$ back to $i$ have a length that is a
      multiple of $k$; if a state is not periodic ($k=1$) it is aperiodic, and the chain is aperiodic
      if all states are. An example of a periodic chain with $k=3$: starting from state 1, the only
      path back to 1 is $1\to2\to3\to1$, so any return takes 3 transitions.
\end{itemize}
\textbf{Both of the latter problems are dealt with by a single strategy:} add a link from each page to
every page and give each such link a small transition probability controlled by a parameter $d$ ---
the augmented transition matrix is obviously irreducible and aperiodic.
\end{defbox}

\begin{defbox}[The final PageRank]
After this augmentation the random surfer has two options at a page: with probability $d$ they
randomly choose an out-link to follow, and with probability $1-d$ they jump to a random page:
\[ P=\Big(\frac{1-d}{n}\Big)\,\mathbf{E}+d\,A^{\!\top}P, \]
where $\mathbf E=\mathbf{ee}^{\!\top}$ ($\mathbf e$ being a column vector of all 1's), i.e.\
$\mathbf E$ is an $n\times n$ matrix of all 1's. The matrix
$\big(\frac{1-d}{n}\big)\mathbf E+dA^{\!\top}$ is stochastic (transposed), irreducible and
aperiodic. Scaling so that $\mathbf e^{\!\top}P=n$:
\[ P=(1-d)\,\mathbf e+d\,A^{\!\top}P,\qquad\text{i.e.\ for each page } i: \quad
   P(i)=(1-d)+d\sum_{j=1}^{n}A_{ji}P(j), \]
which is equivalent to the formula given in the PageRank paper
\[ P(i)=(1-d)+d\sum_{(j,i)\in E}\frac{P(j)}{O_j}. \]

\textbf{Computation by power iteration:}
\begin{algorithmic}[1]
\State $P_0\gets \mathbf e/n$; $k\gets 1$
\Repeat
  \State $P_k\gets (1-d)\mathbf e+dA^{\!\top}P_{k-1}$
  \State $k\gets k+1$
\Until{$\|P_k-P_{k-1}\|_1<\varepsilon$}
\State \Return $P_k$
\end{algorithmic}
\end{defbox}

\textbf{Advantages of PageRank.} \emph{Fighting spam:} a page is important if the pages pointing to it
are important --- and since it is not easy for a page owner to add \emph{in-links} from other
important pages, it is not easy to influence PageRank. \emph{A global, query-independent measure:} the
PageRank values of all pages are computed and saved \emph{off-line} rather than at query time.
\textbf{Criticism:} precisely the query-independence --- PageRank could not distinguish between pages
that are authoritative in general and pages that are authoritative on the query topic.

\textbf{Summary.} Within the framework of link analysis we have seen the principles of
\emph{centrality} and \emph{prestige} and the
algorithms \emph{HITS} and \emph{PageRank} (which powers Google). Yahoo! and MSN have their own
unpublished link-based algorithms. \textbf{Important:} hyperlink-based ranking is not the only
algorithm used in search engines --- it is combined with many content-based factors to produce the
final ranking presented to the user.

\subsection{Themed communities and their discovery}

Links can also be used to find \emph{themed communities} --- groups of content creators or people
sharing some common interests (web communities, email communities, named entity communities). A
related notion is \emph{focused crawling}: combining contents and links to crawl web pages of a
specific topic --- follow links and use learning/classification to determine whether a page is on
topic.

\begin{defbox}[Themed community]
Given a finite set of entities $E=\{e_1,\dots,e_n\}$ \emph{of the same type}, a \emph{themed
community} is a pair $K=(T,M)$, where $T$ is the community theme and $M\subseteq E$ is the set of all
entities in $E$ that share the theme $T$. If $e\in M$, then $e$ is said to be a \emph{member} of the
community $K$.
\begin{itemize}
\item \emph{A theme defines a community:} given a theme $T$, the set of members is uniquely
      determined, so two communities are equal if they have the same theme.
\item \emph{A theme can be defined arbitrarily:} it can be an event (a sport event, a scandal) or a
      concept (web mining).
\item \emph{An entity can be in any number of communities:} communities can therefore \emph{overlap}
      and multiple communities may share members.
\item \emph{The entities in $E$ are of the same type:} people and organisations thus cannot be in the
      same community.
\end{itemize}
Communities can form a \emph{hierarchical structure}: a themed community $(T,M)$ may include a set of
themed \emph{sub-communities} $\{(T_1,M_1),\dots,(T_k,M_k)\}$, where $T_i$ is a sub-theme of $T$ and
$M_i\subseteq M$; $(T,M)$ is then a \emph{super-community}. Each sub-community can be further
decomposed.

\textbf{The objective of discovery:} given a data set containing entities, discover the hidden
communities --- for each community find the \emph{theme} (usually represented by a set of keywords) and
its \emph{members}.
\end{defbox}

\textbf{Community manifestation in the data:} community members are linked in some way and the
associated texts contain words indicative of the theme. \emph{Web pages:} the members are more likely
to be mutually interconnected through \emph{hyperlinks}; the pages contain \emph{words} related to the
theme. \emph{Emails:} the members are more likely to communicate with one another; their emails contain
words related to the theme. \emph{Text documents:} the members are more likely to appear
\emph{together} in the same sentence and/or document; the words in the sentences indicate the theme.

\begin{defbox}[Bipartite core communities]
An $(i,j)$-\emph{bipartite core} is a complete bipartite sub-graph with at least $i$ \emph{fan} nodes
and at least $j$ \emph{center} nodes, i.e.\ it contains all possible edges between the $i$ fan
vertices and the $j$ center vertices. The fans represent \emph{specialised hubs}, the centers
correspond to \emph{authorities}.

\textbf{The procedure for finding bipartite cores:}
\begin{itemize}
\item \emph{Step 1 --- pruning.} Prune by in-degree: delete all pages that are highly referenced,
      i.e.\ whose number of in-links is greater than some $k$ (e.g.\ $k=50$). Then prune fans and
      centers iteratively (to find the $(i,j)$-cores): prune all fans with an out-degree smaller than
      $j$; prune all centers with an in-degree smaller than $i$; delete the edges associated with the
      pruned nodes.
\item \emph{Step 2 --- generating all $(i,j)$-cores.} Find all $(1,j)$-cores, i.e.\ the set of all
      vertices with out-degree at least $j$; construct all $(2,j)$-cores by checking every fan that
      also points to any center in a $(1,j)$-core; continue in the same fashion up to $i$.
\end{itemize}
\textbf{Limitation:} the algorithm finds only the \emph{core} pages of the communities, not all the
member pages, and it finds neither the \emph{themes} of the communities nor their hierarchical
structure.
\end{defbox}

\begin{defbox}[Overlapping communities of named entities]
\textbf{The objective:} find entity communities from a text corpus. \textbf{The main idea:} two
entities are linked if they \emph{co-occur in the same sentence} (if two people are mentioned in the
same sentence, there is usually a relationship between them); a single entity can appear in multiple
communities; the top keywords are assumed to represent the theme of the community. There are promising
results when applied to finding communities of political figures and celebrities from web documents.

\textbf{The algorithm:}
\begin{enumerate}
\item \emph{Build a link graph} --- parse each document and, for each sentence, identify the named
      entities contained in it; if a sentence has more than one named entity, these entities are
      \emph{pair-wise} linked; the keywords in the sentence are attached to the linked pairs as
      textual content; discard all the other sentences.
\item \emph{Find all triangles} --- a triangle consists of three entities bound together and
      represents the basic building block of communities; the observation is that a community tends to
      expand by triangles sharing a common edge.
\item \emph{Find community cores} --- a community core is a group of tightly bound triangles, i.e.\ a
      relaxed complete sub-graph (clique); intuitively, a core consists of a set of tightly connected
      members of a community.
\item \emph{Cluster around community cores} --- the triangles and entity pairs that are in no core are
      assigned to cores according to their \emph{textual content similarities} with the discovered
      cores.
\end{enumerate}
\end{defbox}

\subsection{Community detection}

\begin{defbox}[Community and community detection]
\emph{Communities} are formed by individuals such that those \emph{within} a group interact with each
other \emph{more frequently} than with those outside the group. In different contexts they are also
called a \emph{group}, \emph{cluster}, \emph{cohesive subgroup} or \emph{module}.

\emph{Community detection} is the discovery of groups in a network where the individuals' group
memberships are \emph{not explicitly given}.

\textbf{The definition of a community can be subjective:} a densely knit sub-graph may be a community,
but so may every connected component.

\textbf{Two types of groups in social media:} \emph{explicit} groups formed by user subscriptions, and
\emph{implicit} groups formed by social interactions. Why extract groups from the topology when some
sites allow people to join groups? Not all sites provide a community platform; not all people want to
make an effort to join groups; groups can change \emph{dynamically}. Network interaction provides
information about the relationships between users, complements other kinds of information, helps
network visualisation and navigation, and provides basic information for other tasks.
\end{defbox}

\begin{defbox}[Minimum cut methods]
\textbf{Minimum cut} (Ahuja et al., 1993). A \emph{cut} is a partition of the vertices of a graph into
two disjoint sets; the \emph{minimum cut problem} is to find the partition that minimises the number of
edges running between the two sets (or the sum of the weights in weighted networks) --- the so-called
\emph{cut size}. The motivation: most interactions occur \emph{within} a group and interactions
between groups are rare $\Rightarrow$ community detection $\approx$ the minimum cut problem.

For $k>2$ the strategy of \emph{iterative bisecting} is implemented (in the second and further
iterations, groups are sequentially divided into two sub-groups).

\textbf{Drawback: it produces groups of unbalanced sizes} (one set is often a singleton).

\textbf{Complexity:} a simple solution using the max-flow based $s$-$t$ cut algorithm has complexity
$O(n^5)$; an implementation using Karger's algorithm has $O(n^4)$ but is not guaranteed always to find
the minimum cut. For unit weights $w_{ij}=1$ and no balancing constraints, the 2-way cut problem is
polynomially solvable in $O(n^2\log n)$ (Karger's randomised minimum cut). Arbitrary edge weights
\emph{and} balancing constraints make the problem \textbf{NP-hard}.
\end{defbox}

\begin{defbox}[Ratio cut and normalised cut]
The minimum cut often returns an imbalanced partition (one set being a singleton) $\Rightarrow$ we
change the objective function to consider the \emph{community size}. For a partition
$\pi=\{C_1,\dots,C_k\}$:
\[ \mathrm{RatioCut}(\pi)=\frac{1}{k}\sum_{i=1}^{k}\frac{\mathrm{cut}(C_i,\bar C_i)}{|C_i|},\qquad
   \mathrm{NCut}(\pi)=\frac{1}{k}\sum_{i=1}^{k}\frac{\mathrm{cut}(C_i,\bar C_i)}{\mathrm{vol}(C_i)}, \]
where $|C_i|$ is the number of nodes in $C_i$, $\mathrm{vol}(C_i)$ the sum of degrees in $C_i$ and
$\mathrm{cut}(A,B)=\sum_{i\in A,\,j\in B}a_{ij}$. \textbf{Both measures prefer a balanced partition.}
\end{defbox}

\begin{defbox}[The Girvan--Newman algorithm (2002)]
A \emph{divisive hierarchical} algorithm (a top-down approach): it decomposes the initial full graph
into progressively smaller connected pieces, until there are no edges to remove and each node
represents a community itself.

\textbf{The idea:} identify the edges that connect vertices belonging to \emph{different} communities
(the so-called bridges) and iteratively remove them. \textbf{The criterion:} \emph{edge betweenness}
--- the fraction of shortest paths that pass along the edge; edges with high betweenness lie on a
large number of shortest paths between vertices and are therefore bridges.

\begin{algorithmic}[1]
\State compute the betweenness of all edges in the network
\State remove the edge with the \emph{highest} betweenness
\State repeat the previous steps until there are no edges left to remove
\end{algorithmic}
\textbf{Input:} the entire graph. \textbf{Output:} a hierarchical structure, represented e.g.\ as a
dendrogram.

\textbf{Selecting the number of communities:} the algorithm returns a \emph{set} of possible
solutions but does not indicate which one is the best. To select the best partition, compute the
\emph{modularity} of each network division and choose the partition with the highest modularity.

\textbf{Drawback:} suitable only for networks of \emph{moderate size} due to the high computational
costs --- $O(m^2n)$ for a graph with $m$ edges and $n$ vertices, resp.\ $O(n^3)$ for sparse networks.
\end{defbox}

\emph{An example of computing edge betweenness:} the edge betweenness of $e(1,2)$ is 4, as all the
shortest paths from 2 to $\{4,5,6,7,8,9\}$ have to pass either $e(1,2)$ or $e(2,3)$, and $e(1,2)$ is
the shortest path between 1 and 2. After removing $e(4,5)$, the betweenness of $e(4,6)$ becomes 20
(the highest); after removing $e(4,6)$, the edge $e(7,9)$ has the highest betweenness, 4.

\begin{defbox}[Modularity $Q$]
It quantifies the quality of a given graph division into communities. \textbf{The basic idea:} a
network has a meaningful community structure if the number of edges \emph{between} communities is
smaller than would be expected based on a random choice. The expected number of edges between nodes
$i$ and $j$ with degrees $k_i$, $k_j$ is $k_ik_j/2m$.
\[ Q=\frac{1}{2m}\sum_{i,j}\Big(A_{ij}-\frac{k_ik_j}{2m}\Big)\delta(c_i,c_j), \]
where $A_{ij}$ is the entry of the adjacency matrix (the number of edges between $i$ and $j$), $m$ the
number of edges, $k_i$ the degree of vertex $i$, $c_i$ the group to which $i$ belongs, and $\delta$
the Kronecker delta.

\textbf{Values:} $Q\in[-1,1]$ and can be either positive or negative. If positive, there is a
possibility of finding a community structure in the network; if the values are large and positive, the
corresponding partition may reflect the true community structure. For meaningful communities the
modularity should be $Q\ge 0.3$ (Clauset et al., 2004).
\end{defbox}

\begin{defbox}[The Louvain method (Blondel et al., 2008)]
A \emph{greedy} optimisation method for \emph{agglomerative hierarchical} modularity maximisation. It
comprises two phases:

\textbf{Phase 1} --- consider each node as a single community:
\begin{enumerate}
\item compute the modularity $Q$;
\item move the isolated node from its community to a neighbouring community;
\item compute the gain/loss in modularity yielded by the changed assignment;
\item if the modularity increases (there is a gain), keep the node in the \uv{new} community;
\item repeat until no further improvements are possible.
\end{enumerate}
\emph{Output:} the $k$-th level partition ($k$ being the number of the iteration).

\textbf{Phase 2} --- create a new network (a \emph{supergraph}) derived from the original one, where
each new node (a \emph{supervertex}) is the aggregation of the nodes assigned to a given community
found in Phase 1; two supervertices are connected by an edge if there is at least one edge linking two
vertices inside the corresponding communities.

\textbf{Iteration:} repeat Phases 1 and 2 until the maximum modularity is attained.

\textbf{Advantages:} good performance also on large network graphs at low computational costs ---
$O(n\log n)$, where $n$ is the number of vertices. \textbf{Drawback:} the results are
\emph{order-sensitive}.

\textbf{Further drawbacks:} \emph{disconnected communities} --- when a node is moved to a different
community, the original community may become \emph{internally disconnected}, yet its nodes are still
locally optimally assigned and will therefore remain in it; \emph{getting stuck in local optima};
\emph{failing to refine communities} after the initial merging.
\end{defbox}

\begin{defbox}[The Leiden algorithm (Traag, Waltman, van Eck, 2019)]
A refinement of the Louvain method for agglomerative hierarchical modularity maximisation. It has
three phases: \emph{(1)}~local moving of nodes (similar to Louvain); \emph{(2)}~\textbf{refinement of
the partition}; \emph{(3)}~network aggregation (similar to Louvain).
Refinement means that every community from phase~1 is split from the inside into connected,
well-linked subcommunities (nodes merge only within their own community) and only these are
aggregated, so the disconnected communities known from Louvain cannot arise.

\textbf{Advantages:} more efficient than the Louvain method, but it may lead to extended processing
times for massive graphs; parallel multicore implementations may boost the speed. \textbf{Drawback:} a
\emph{hard partition} --- the nodes can belong to only one community.
\end{defbox}

\textbf{Challenges of community detection.} Despite the many algorithms, community detection is still
considered a challenging problem, for two main reasons: \emph{(1)}~typically the number of communities
is unknown and has to be determined; \emph{(2)}~communities are usually \emph{heterogeneous} in terms
of size and density.

\subsubsection{A taxonomy of community detection methods}

Four categories according to the level at which the criterion is placed: \emph{node-centric}
(each node of the group satisfies a property), \emph{group-centric} (the group as a whole
does), \emph{network-centric} (partition the whole network into disjoint sets) and
\emph{hierarchy-centric} (a hierarchy of communities).

\begin{defbox}[Node-centric community detection]
\begin{itemize}
\item \textbf{Complete mutuality --- cliques.} A \emph{clique} $=$ a maximal complete
      subgraph; finding the maximum clique is \textbf{NP-hard}. \emph{Recursive pruning:} in
      a clique of size $k$ each node has degree $\ge k-1$, so lower-degree nodes can be
      removed repeatedly (in scale-free networks most get pruned). Cliques serve rather as
      \emph{cores} (seeds) of larger communities.
\item \textbf{Clique percolation} (CPM) --- finds \emph{overlapping} communities: find all
      cliques of size $k$, build the \emph{clique graph} (two cliques are adjacent iff they
      share $k-1$ nodes); the communities $=$ the connected components of the clique graph.
\item \textbf{Reachability.} A \emph{$k$-clique} --- a maximal subgraph with pairwise
      geodesic distance $\le k$ (beware: the paths may run \emph{outside} the subgraph, the
      inner diameter can exceed $k$); a \emph{$k$-club} --- diameter $\le k$; a
      \emph{$k$-clan} --- a $k$-clique that is also a $k$-club.
\item \textbf{Nodal degrees.} A \emph{$k$-core} --- each node connected to at least $k$
      members; a \emph{$k$-plex} --- in a group of $n_s$ nodes each node connected to at
      least $n_s-k$ members. The definitions are complementary, but the sets of all
      $k$-cores and $(n_s-k)$-plexes differ ($n_s$ is not fixed).
\item \textbf{Within vs.\ outside ties.} An \emph{LS set} --- every proper subset has more
      ties inside the group than outside; too strict for practice. A relaxation is the
      \emph{$k$-component} --- connected even after removing fewer than $k$ vertices or
      edges (computed via minimum cut / maximum flow).
\end{itemize}
\textbf{Limitations:} the definitions are strict, do not scale, and may not match the
properties of scale-free networks --- they are mainly useful as community cores.
\end{defbox}

\begin{defbox}[Group-centric detection: quasi-cliques]
The condition is placed on the group as a whole, typically its density: a subgraph
$G_S=(V_S,E_S)$ is a \emph{$\gamma$-dense quasi-clique} if
$|E_S|\ge\gamma\cdot\frac{|V_S|(|V_S|-1)}{2}$. It is found by local search over a sample with
heuristic pruning (remove nodes with degree $<\gamma(k-1)$, i.e.\ below the average degree).
\end{defbox}

\begin{defbox}[Network-centric community detection]
The connections are considered \emph{globally}; the goal is to partition the whole network
into \emph{disjoint} sets:
\begin{itemize}
\item \textbf{vertex similarity} --- $k$-means over the similarity of \emph{neighbourhoods}
      ($\mathrm{Jaccard}=\frac{|N_i\cap N_j|}{|N_i\cup N_j|}$,
      $\cos=\frac{|N_i\cap N_j|}{\sqrt{|N_i||N_j|}}$); structural equivalence is too
      restrictive;
\item \textbf{latent space models} --- MDS: map the nodes into a low-dimensional space
      ($S=V\Lambda^{1/2}$ from the double-centred proximity matrix), then $k$-means;
\item \textbf{block models} --- $A\approx S\,\Sigma\,S^{\!\top}$ ($S$ the community
      indicator matrix, $\Sigma$ the mixing matrix); after relaxing $S$, the optimum is the
      top eigenvectors of $A$;
\item \textbf{spectral clustering} --- the ratio and normalised cut $=$ minimising
      $\mathrm{tr}(S^{\!\top}LS)$ with the \emph{graph Laplacian} $L=D-A$, resp.\
      $D^{-1/2}(D-A)D^{-1/2}$; after relaxation the optimum is the eigenvectors with the
      \emph{smallest} eigenvalues (the first one is useless);
\item \textbf{modularity maximisation} --- $Q=\frac{1}{2m}\sum_C\sum_{i,j\in C}
      \big(A_{ij}-\frac{d_id_j}{2m}\big)$; the optimum is the \emph{top} eigenvectors of the
      \emph{modularity matrix} $B=A-\frac{dd^{\!\top}}{2m}$, followed by $k$-means.
\end{itemize}
\textbf{A unified view:} all four matrix approaches are factorisations
$X\approx S\,\Sigma\,S^{\!\top}$, with $X$ being, in turn, the proximity matrix, the
sociomatrix $A$, the graph Laplacian and the modularity matrix. \textbf{Limitation:} the
number of communities must be specified beforehand.
\end{defbox}

\begin{defbox}[Hierarchy-centric community detection]
Builds a hierarchy of communities (analysis at different \emph{resolutions}).
\textbf{Divisive} --- e.g.\ Girvan--Newman ($O(m^2n)$, for sparse networks $O(n^3)$):
repeatedly remove the \uv{weakest} edge (the highest \emph{edge betweenness} $=$ a bridge)
until the network decomposes into the desired number of components.
\textbf{Agglomerative} --- e.g.\ Louvain ($O(n\log n)$): start with the nodes as communities
and merge them by the \emph{modularity increase}; \emph{Walktrap} merges according to the
results of short random walks.
\end{defbox}

\subsection{Link prediction}

\begin{defbox}[The link prediction task]
\textbf{The objective:} predict \emph{future} links between pairs of nodes (friend
recommendation on Facebook or LinkedIn). \textbf{Content-based measures} build on
\emph{homophily} (similar content $\to$ a link); they are, however, sensitive to the network
at hand (noisy tweets), and content-based homophily does \emph{not} guarantee structural
connectivity. \textbf{Structural measures} build on \emph{triadic closure} (a shared
neighbourhood $\to$ a link); they work across domains and are \emph{far more powerful} ---
structural connectivity usually implies content-based homophily. They include the
\emph{neighbourhood-based} measures, the \emph{Katz} measure and \emph{random walk}-based
measures.
\end{defbox}

\begin{defbox}[Neighbourhood-based link prediction measures]
$S_i$ $=$ the neighbour set of node $i$; each measure fixes a weakness of the previous one:
\begin{itemize}
\item \textbf{common neighbours} $\mathrm{CN}(i,j)=|S_i\cap S_j|$ --- ignores the degrees:
      two popular figures have many common neighbours \emph{just by chance};
\item \textbf{Jaccard} $\frac{|S_i\cap S_j|}{|S_i\cup S_j|}$ --- normalises for the degrees
      of the end nodes, but not of the \emph{intermediate} neighbours: a very popular common
      neighbour is common to many pairs and hence less important for the prediction;
\item \textbf{Adamic--Adar} $\mathrm{AA}(i,j)=\sum_{k\in S_i\cap S_j}\frac{1}{\log|S_k|}$
      --- the weight of a common neighbour decreases with its degree.
\end{itemize}
\end{defbox}

\begin{defbox}[The Katz measure]
With few common neighbours the neighbourhood-based measures cannot distinguish; the Katz
measure captures \emph{indirect} connectivity through walks of all lengths ($W_{ij}^{(t)}$
$=$ the number of walks of length $t$, $\beta<1$ a \emph{discount factor} de-emphasising
long walks):
\[ \mathrm{Katz}(i,j)=\sum_{t=1}^{\infty}\beta^{t}\,W_{ij}^{(t)},\qquad
   K=\sum_{t=1}^{\infty}\beta^{t}A^{t}=(I-\beta A)^{-1}-I. \]
The sum of the Katz coefficients of a node with respect to the other nodes is its
\textbf{Katz centrality}.
\end{defbox}

\begin{defbox}[Link prediction as a classification problem]
The presence of an edge $=$ a binary class; the features of a pair $=$ the similarity
measures (neighbourhood-based, Katz, walk-based) $+$ preferential-attachment features (node
degrees). It is a \emph{positive-unlabelled} problem: pairs with an edge are positive, and
approximate \uv{negative} examples are sampled from the unlabelled ones (a sparse network has
far too many). The base classifier is typically \emph{logistic regression}, in a
cost-sensitive version due to the imbalance. \textbf{Advantages:} content features are easy
to add; the classifier learns the relevance of the features \emph{by itself}; for directed
networks the features can be extracted \emph{asymmetrically} (in- and out-degrees,
personalised PageRank).
\end{defbox}

\textbf{Summary:} \emph{neighbourhood-based} measures are cheap even for extensive data and
perform almost as well as the other unsupervised ones; the \emph{Katz} and \emph{walk}-based
measures help in \emph{very sparse} networks with few common neighbours; \emph{supervision}
is more accurate and adaptable but expensive; content enhances the prediction, yet structural
measures dominate.

\subsection{Summary for the examination}

\begin{itemize}
\item \textbf{SNA} $=$ social science $+$ network analysis $+$ graph theory; it focuses on
      \emph{relationships}, not on entities and their attributes. Representations: a drawing, an edge
      list, an adjacency matrix (asymmetric for directed graphs). \emph{Ego network} vs.\ \uv{whole}
      network --- the \emph{boundary specification problem}: no studied network is whole.
\item \textbf{Tie strength}: strong vs.\ weak; \emph{the strength of weak ties} (Granovetter 1973).
      Homophily, transitivity and bridging; homophily $+$ transitivity $\to$ \emph{cliques}. A bridge
      $=$ an edge whose removal disconnects its endpoints; a \emph{shortcut bridge} $=$ the distance
      after removal grows beyond 2. \textbf{Neighbourhood overlap}
      $O(i,j)=\frac{|N_i\cap N_j|}{|N_i\cup N_j|-2}$ --- the larger, the stronger the tie.
\item \textbf{Four centrality measures}: \emph{degree} (how many people it reaches directly),
      \emph{betweenness} (on how many shortest paths it lies; where the network breaks apart),
      \emph{closeness} (the reciprocal mean shortest path length; the speed of dissemination),
      \emph{eigenvector} ($Ax=\lambda x$, only the \emph{greatest} eigenvalue gives a non-negative
      vector; how well it is connected to well-connected nodes). Know the interpretation table and be
      able to explain why one measure is not enough (nodes 3 and 5 \emph{together} $>$ node 10).
      \textbf{Structural equivalence}: the positions can be swapped without changing the structure.
\item \textbf{Cohesion}: reciprocity (directed graphs only), density
      ($|E|/\binom{|V|}{2}$; a clique $=1$; directed graphs have half), the \textbf{clustering
      coefficient} $C_i=\frac{2|\{e_{jk}\}|}{k_i(k_i-1)}$ and the network average, the
      \emph{diameter} $=$ the longest shortest path, the average distance, eccentricity. A
      \textbf{small world} $=$ high $C$ $+$ a short average path.
\item \textbf{Six properties of real-world complex networks}: the small-world effect (Milgram, median
      6), transitivity/clustering, power-law degree distributions, resilience (robust to random,
      vulnerable to targeted removal of hubs), mixing patterns (assortative mixing $=$ homophily),
      community structure. \textbf{Scale-free networks}: growth $+$ \emph{preferential attachment}
      (rich-get-richer, the Matthew effect, cumulative advantage); the reasons --- popularity, quality,
      a mixed model (the halo effect). \textbf{Centralisation} and the \emph{core-periphery} structure
      (bow-tie analysis).
\item \textbf{Influence modelling}: the common features (a directed graph, active/inactive, \emph{an
      activated node cannot be deactivated}). \textbf{LTM} --- a threshold $\theta_v$ uniform on
      $[0,1]$, activation when $\sum_{\text{active}}w_{u,v}\ge\theta_v\sum w_{u,v}$; focuses on the
      \emph{receivers}. \textbf{ICM} --- a \emph{single} chance to activate each neighbour with
      probability $p_{v,u}$; focuses on the \emph{senders}. \textbf{Influence maximisation} is NP-hard
      and the greedy approach gives \textbf{63~\%} of the optimum.
\item \textbf{Influence vs.\ correlation}: the \emph{test for correlation} (the fraction of edges
      between different attribute values vs.\ $2p(1-p)$; the smoker example $14\,\%<49\,\%$); the three
      processes --- homophily, \emph{confounding}, influence; the logistic model
      $\ln\frac{p(a)}{1-p(a)}=\alpha\ln(a+1)+\beta$ with $\alpha$ $=$ the social correlation
      coefficient; the \textbf{shuffle test} --- shuffle the timestamps, and if $\alpha'$ differs
      significantly from $\alpha$ this is evidence of \emph{influence}.
\item \textbf{The homophily test}: significantly fewer heterogeneous edges than $2pq$ $\Rightarrow$
      homophily; significantly \emph{more} $\Rightarrow$ \emph{inverse} homophily (dating
      relationships). Be able to give the example: 5 out of 18 vs.\ $2pq=8$ out of 18.
\item \textbf{Selection vs.\ social influence}: selection $=$ characteristics drive link formation;
      influence $=$ links shape \emph{mutable} characteristics. \emph{Immutable} (gender, race, age)
      vs.\ \emph{mutable} (behaviour, opinions). They can be told apart only by \emph{longitudinal}
      studies (obesity: Christakis \& Fowler, 12\,000 people over 32 years $\to$ evidence for
      \emph{social influence}).
\item \textbf{Affiliation networks}: bipartite (persons $\times$ foci); a social-affiliation network has
      two types of edges. \textbf{Three closures}: triadic (a common friend), \emph{focal} (a common
      focus; \emph{selection}), \emph{membership} (a friend already in the focus; \emph{social
      influence}). The methodology for measuring $T(k)$ from two snapshots; the baseline model
      $T_{\text{base}}(k)=1-(1-p)^k$; empirically a large increase from 1 to 2 and then
      \emph{diminishing returns}.
\item \textbf{The Schelling model}: agents of two types in a grid, satisfied with $\ge t$ neighbours of
      the same type, the unsatisfied move $\Rightarrow$ \emph{self-imposed segregation}. It shows how an
      \emph{immutable} characteristic becomes highly correlated with a \emph{mutable} one (the choice of
      where to live) and how local homophily creates a global pattern.
\item \textbf{Structural balance}: the balanced triangles are $+{+}{+}$ and $+{-}{-}$ (all three
      positive, or exactly one positive). \textbf{The Balance Theorem}: a complete balanced graph
      $\Rightarrow$ all friends, or a division into two groups of internal friends who are mutual
      enemies --- be able to give the proof ($X$ $=$ the friends of $A$, $Y$ $=$ the enemies of $A$,
      three steps). \textbf{Weak balance}: only the \emph{two positive and one negative} configuration
      is forbidden $\Rightarrow$ a division into \emph{any number} of groups. Applications: the European
      alliances 1872--1907, Bangladesh 1972.
\item \textbf{HITS}: query-\emph{dependent}; the root set $S$ ($t=200$) $\to$ the base set $B$;
      $a=L^{\!\top}h$, $h=La$, power iteration with \textbf{indispensable normalisation};
      $a_k=(L^{\!\top}L)^{k-1}a_0$ converges to the principal eigenvector of $L^{\!\top}L$.
      Weaknesses: easily spammed, topic drift, inefficiency at query time.
\item \textbf{PageRank}: query-\emph{independent}, $P(i)=\sum_{(j,i)\in E}P(j)/O_j$, \emph{needs no
      normalisation} (the total PageRank is constant --- the \uv{fluid}). The \uv{slow leak} problem
      $\Rightarrow$ the \emph{scaled} rule with the damping factor $d$ (0.8--0.9; $0.85$ in the paper),
      the \uv{water cycle}. A Markov chain needs a \emph{stochastic, irreducible and aperiodic} matrix
      --- the web graph satisfies none of them (dangling pages, not strongly connected, periodicity);
      all of it is dealt with by a \emph{single} strategy (a link from each page to every page with a
      small probability). The final form $P(i)=(1-d)+d\sum_{(j,i)\in E}P(j)/O_j$, computed by power
      iteration. Advantages: resistance to spam (in-links cannot easily be obtained), a global offline
      measure; the criticism: query-independence.
\item \textbf{Themed communities}: $K=(T,M)$, the theme defines the community, communities can
      \emph{overlap}, a hierarchy of sub-communities. \emph{Bipartite cores} --- an $(i,j)$-core, fans
      $=$ hubs, centers $=$ authorities; pruning $+$ core generation; it finds only the cores, neither
      the themes nor the hierarchy. \emph{Named entity communities}: a graph from co-occurrences in a
      sentence $\to$ triangles $\to$ cores $\to$ clustering by textual similarity.
\item \textbf{Community detection --- the definition is subjective}; explicit vs.\ implicit groups.
      \emph{Minimum cut}: unbalanced groups, $O(n^5)$ / Karger $O(n^4)$, \textbf{NP-hard} with general
      weights and balancing constraints $\Rightarrow$ the \emph{ratio} and \emph{normalised cut}
      (dividing by $|C_i|$, resp.\ $\mathrm{vol}(C_i)$).
\item \textbf{Girvan--Newman}: divisive, removes the edges with the highest \emph{betweenness},
      $O(m^2n)$ / $O(n^3)$; the number of communities is chosen by the highest \textbf{modularity}
      $Q=\frac{1}{2m}\sum_{i,j}(A_{ij}-\frac{k_ik_j}{2m})\delta(c_i,c_j)$, $Q\in[-1,1]$, meaningful
      communities at $Q\ge 0.3$.
\item \textbf{Louvain}: greedy agglomerative modularity maximisation in two phases (local moves $\to$
      supergraph), $O(n\log n)$; drawbacks: order sensitivity, \emph{disconnected communities}, local
      optima. \textbf{Leiden} adds a \emph{partition refinement} phase; hard partition only. The
      challenges: the unknown number of communities, heterogeneous sizes and densities.
\item \textbf{Four categories of methods}: \emph{node-centric} (cliques --- NP-hard, recursive pruning
      of nodes with degree $<k-1$, \textbf{CPM} for \emph{overlapping} communities via a clique graph
      of cliques sharing $k-1$ nodes; $k$-clique / $k$-club / $k$-clan; $k$-core and $k$-plex are
      complementary; LS sets and $k$-components), \emph{group-centric} ($\gamma$-dense quasi-cliques),
      \emph{network-centric} (vertex similarity --- Jaccard, cosine $+$ $k$-means; MDS; block models;
      spectral clustering via the graph Laplacian with the \emph{smallest} eigenvalues; modularity
      maximisation via the modularity matrix with the \emph{largest} --- all unified as the matrix
      factorisation $X\approx S\Sigma S^{\!\top}$, which requires the number of communities to be
      given), \emph{hierarchy-centric} (divisive $=$ Girvan--Newman, agglomerative $=$ Louvain,
      Walktrap).
\item \textbf{Link prediction}: content-based measures (homophily) vs.\ structural ones (triadic
      closure; stronger). Neighbourhood-based measures: \emph{common neighbours} $|S_i\cap S_j|$ (does
      not account for degrees), \emph{Jaccard} $\frac{|S_i\cap S_j|}{|S_i\cup S_j|}$ (normalises the
      degrees of the two nodes but not of the intermediate neighbours), \emph{Adamic--Adar}
      $\sum_{k}\frac{1}{\log|S_k|}$ (the weight decreases with the degree of the common neighbour).
      \textbf{The Katz measure} $\sum_t\beta^tW_{ij}^{(t)}=(I-\beta A)^{-1}-I$ for \emph{sparse}
      networks with few common neighbours; the sum over the other nodes $=$ \emph{Katz centrality}.
      \textbf{The classification view}: a positive-unlabelled problem, a sample of negatives, logistic
      regression as the base classifier, cost-sensitive variants; the advantage --- the classifier
      learns the relevance of the features itself and \emph{asymmetric} features can be used for
      directed networks.
\end{itemize}

%=====================================================================
\newpage
\section{Practical use of the techniques}
\label{sec:praxe}
%=====================================================================

The last sentence of the topic asks for both halves to be brought together: what the techniques
covered are \emph{actually} used for. In the examination it helps to be able to say, for each
application, \emph{which task} it solves (Chapter~\ref{sec:paradigmata}), \emph{by which method}
(Chapter~\ref{sec:metody}) and \emph{how it is evaluated} (Chapter~\ref{sec:reprez}).

\subsection{Data mining}

\begin{center}
\small
\begin{tabular}{@{}p{3.5cm}p{4.4cm}p{6.9cm}@{}}
\toprule
\textbf{Application} & \textbf{Task type / method} & \textbf{Note} \\
\midrule
Market basket analysis & association rules, Apriori / FP-Growth &
  planning and layout of shops, coupons and time-limited discounts, \uv{packaging} of products;
  comparison of new and existing shops using \emph{virtual items} \\
Fraud detection & classification, outlier detection &
  20 million procedures billed to a health insurance company $\to$ 214 thousand odd ones $\to$ 14
  thousand passed on for manual review $\to$ 8 thousand confirmed frauds; the machine
  \emph{focuses} the expert's attention \\
Credit scoring & decision trees, naive Bayes, logistic regression &
  the running example throughout the text; \emph{comprehensibility} is required (regulation), as is
  the handling of missing values \\
Customer segmentation & cluster analysis ($k$-means, hierarchical) &
  representing a segment by its \emph{centroid} (the \uv{typical customer}) is very intuitive for
  business \\
Churn prediction & classification on imbalanced data &
  accuracy is not enough $\Rightarrow$ precision/recall, $F_1$, ROC/AUC; \emph{lift analysis} for
  deciding how many customers to contact \\
Targeted marketing campaigns & scoring and ranking, lift &
  the economic reasoning over the lift curve (the watch example: bin 1 $+\$550$, bin 10
  $-\$475$) \\
Text classification & SVM with kernels &
  among the best classifiers for text (high-dimensional data); comprehensibility is not expected \\
Time series prediction & regression, neural networks &
  the series is transformed by a sliding window into a table of inputs and outputs \\
Classification of chemical compounds & classification of structural data &
  representation by SMILES strings or facts in Prolog \\
Web usage mining & sequential patterns (GSP) &
  navigational patterns of users based on the sequences of page visits \\
Recommender systems \mimo{} & cluster analysis, classification, link prediction &
  collaborative filtering (similarity of users or items) vs.\ content-based filtering; the cold-start
  problem is addressed by content \\
\bottomrule
\end{tabular}
\end{center}

\textbf{A note on the motivation for decision trees.} The lecture cites the case of Barclays bank,
which introduced a telephone system routing callers by their credit rating: clients with considerable
savings or a sufficient credit limit were connected to a British operator, the others to an Indian
call centre. The aim was to \uv{filter out} the customers who do not have the money to buy the bank's
products. The criticism: the system is preferential towards wealthier customers and is not intended
to give the best advice. It is a good illustration that a \emph{technically correct} classification
model can have problematic consequences --- and that the \emph{Evaluation} phase of CRISP-DM must
include a view beyond the accuracy of the model.

\subsection{Social network analysis}

\begin{center}
\small
\begin{tabular}{@{}p{3.7cm}p{11.2cm}@{}}
\toprule
\textbf{Area} & \textbf{Use} \\
\midrule
Businesses & analysis and improvement of the communication flow in an organisation, with partners or
  customers; marketing campaigns based on social networking (visualisation of the communication flows
  before and after the introduction of a content management system) \\
Law enforcement agencies and the army & identification of criminal and terrorist networks from traces
  of collected communication; identification of the \emph{key players} in them (centrality measures,
  Chapter~\ref{sec:sna}) \\
Social network sites (Facebook) & identification and recommendation of potential friends based on
  \uv{friends-of-friends} --- precisely the \emph{link prediction} task \\
Civil society organisations & uncovering conflicts of interest in hidden connections between
  government bodies, lobbies and businesses \\
Network operators & optimisation of the structure and capacity of communication networks \\
Medicine & contagious diseases and their prevention; sexually transmitted diseases --- the syphilis
  outbreak in Rockdale County (Atlanta, 1996): a group of young white girls about 16 years old, a
  group of white boys aged 17 to 21, and a group of African-American boys aged 17 to 21; syphilis was
  diagnosed in 6 white females, 2 white males and 2 African-American males \\
Scale-free networks & \emph{computing:} networks with an SF architecture (the WWW) are extremely
  resistant to random failures but very vulnerable to coordinated attacks --- the eradication of
  internet viruses is essentially impossible; \emph{medicine:} vaccination campaigns focused on hubs
  (people with many contacts --- though identifying them is difficult), new drugs targeting the key
  molecules (hubs) involved in a disease, minimising the side effects by using the link networks found
  inside the cells; \emph{business:} understanding the mutual links between companies, industry and
  the economy in order to monitor and eliminate cascaded financial crashes; marketing --- exploring
  the spread of contagion in SF networks and more efficient advertising of new products \\
Community detection & \emph{A Belgian mobile phone operator} (the Louvain method): about 2 million
  customers; the size of a node is proportional to the number of individuals in the community and the
  colour represents the main language spoken in the community (red for French, green for Dutch); only
  communities of more than 100 customers are plotted. A zoom at higher resolution reveals an
  \emph{intermediate} community made of sub-communities with less apparent language separation \\
\bottomrule
\end{tabular}
\end{center}

\textbf{Why and when to use SNA.} To study a social network (offline or online) or improve its
effectiveness; to visualise data so as to uncover patterns in relationships or interactions; to follow
the paths along which information flows in social networks; in quantitative research (though for
qualitative research a network perspective is also valuable) --- the range of actions and
opportunities afforded to individuals is often a function of their \emph{positions} in social
networks, and uncovering these positions (say as fathers, mothers, teachers, workers) can yield
interesting and sometimes surprising results; a quantitative analysis of a network helps identify
different types of actors (key players); to analyse social media in general, identify the causes of
dysfunctional communities or networks, and promote social cohesion and growth in an online community.

\begin{poznbox}[Thoughts on design]
SNA theory has inspired some of the first social network sites (e.g.\ SixDegrees), but it is still not
used so often in conjunction with design decisions. Open questions:
\begin{itemize}
\item How can an online social media platform (or its administrators) leverage the methods and
      insights of SNA?
\item How can a platform encourage its users to act (\emph{locally}) in such a way that would expand
      the community?
\item What measures can an online community take to optimise its structure? Example: \emph{cliques}
      can be undesirable because they shun newcomers.
\item What are the desirable structures for different types of online platforms?
\item How can online communities identify and utilise key players for their benefit?
\end{itemize}
\end{poznbox}

\subsection{Operational and ethical aspects of deployment}

\mimo{} Deploying a model (the \emph{Deployment} phase of CRISP-DM,
Section~\ref{ssec:metodologie}) is not the end of the project:
\begin{itemize}
\item \textbf{Monitoring and drift.} The distribution of the data changes over time (\emph{data
      drift}), as does the relation between the inputs and the target (\emph{concept drift}) --- the
      model has to be retrained regularly and the drop in quality on production data monitored.
\item \textbf{Feedback of the model into the data.} A deployed model influences the very reality it
      measures: a scoring model for approving loans only ever sees the applicants to whom it granted a
      loan, so the further training data is biased by the model itself.
\item \textbf{Privacy.} Social network analysis works with data \emph{about relationships} that the
      user often could not influence; even an anonymised graph can be de-anonymised from its
      structure. Link prediction and influence modelling are essentially tools for predicting and
      influencing the behaviour of concrete people.
\item \textbf{Fairness and discrimination.} A model learned on historical data reproduces historical
      inequality (see the Barclays case above); a protected attribute can often be
      \emph{reconstructed} from the other attributes, so simply omitting it is not enough.
\item \textbf{The legal framework.} GDPR Art.~22 (the right to an explanation of an automated
      decision) and the AI Act for high-risk systems --- in regulated domains interpretability is a
      legal requirement, not a convenience.
\end{itemize}

\subsection{Summary for the examination}

\begin{itemize}
\item For each application, be able to give the \textbf{triple}: the type of task (descriptive /
      predictive, supervised / unsupervised), the method used, and the way it is evaluated.
\item \textbf{Flagship data mining applications}: market basket analysis (association rules),
      fraud detection (the health insurance example --- the machine \emph{focuses} the expert's
      attention rather than replacing them), credit scoring (trees and Bayes; the comprehensibility
      requirement), customer segmentation (clustering, the centroid $=$ the typical customer), churn
      and targeted marketing (imbalanced data $\to$ ROC/AUC and \emph{lift}), text classification
      (SVM), time series prediction (a sliding window), web usage mining (sequential patterns).
\item \textbf{Flagship SNA applications}: identification of key players in criminal and terrorist
      networks (centrality), friend recommendation (link prediction), optimisation of communication
      networks, epidemiology (Rockdale County 1996), vaccination targeted at \emph{hubs}, marketing
      and cascaded financial crashes in scale-free networks, segmentation of an operator's customers by
      language communities (Louvain, 2 million customers).
\item \textbf{The key property of SF networks for practice}: resilience to random failures (up to
      80~\% of nodes) vs.\ vulnerability to a targeted attack (5--15~\% of hubs) --- from which both
      the vaccination strategy and the impossibility of eradicating internet viruses follow.
\item \textbf{Deployment is not the end}: monitoring of drift, feedback of the model into the data,
      privacy (even an anonymised graph can be de-anonymised), fairness (a protected attribute can be
      reconstructed from the others), GDPR Art.~22 and the AI Act.
\end{itemize}

\end{document}
